\documentclass[14pt,letterpaper]{article}

% ======================s======================================
% PACKAGES
% ============================================================

% Page layout
\usepackage[margin=1in]{geometry}
\usepackage{setspace}
\setlength{\parskip}{1em}
\setlength{\parindent}{0pt}
\onehalfspacing

% Mathematical packages
\usepackage{amsmath}
\usepackage{amssymb}
\usepackage{amsthm}
\usepackage{amsfonts}
\usepackage{mathtools}
\usepackage{bm}
\usepackage{comment}

% Graphics and numerical illustrations
\usepackage{graphicx}
\usepackage{float}
\usepackage{subcaption}
\usepackage{booktabs}
\usepackage{array}

% Mathematical packages
\usepackage{amsmath}
\usepackage{amssymb}
\usepackage{amsthm}
\usepackage{amsfonts}
\usepackage{mathtools}
\usepackage{bm}
\usepackage{mathrsfs}  % Required for \mathscr

% Colors
\usepackage[dvipsnames]{xcolor}

% MATLAB code formatting
\usepackage{listings}

% References and hyperlinks
\usepackage[
    colorlinks=true,
    linkcolor=MidnightBlue,
    citecolor=ForestGreen,
    urlcolor=RoyalBlue
]{hyperref}

% Better cross-referencing
\usepackage[nameinlink,noabbrev]{cleveref}

% Typography
\usepackage{microtype}
\usepackage{enumitem}

% ============================================================
% THEOREM ENVIRONMENTS
% ============================================================

\theoremstyle{plain}
\newtheorem{theorem}{Theorem}[section]
\newtheorem{lemma}[theorem]{Lemma}
\newtheorem{proposition}[theorem]{Proposition}
\newtheorem{corollary}[theorem]{Corollary}

\theoremstyle{definition}
\newtheorem{definition}[theorem]{Definition}
\newtheorem{example}[theorem]{Example}

\theoremstyle{remark}
\newtheorem{remark}[theorem]{Remark}

% ============================================================
% MATHEMATICAL NOTATION
% ============================================================

\newcommand{\R}{\mathbb{R}}
\newcommand{\C}{\mathbb{C}}

\newcommand{\norm}[1]{\left\lVert #1 \right\rVert}
\newcommand{\abs}[1]{\left| #1 \right|}
\newcommand{\inner}[2]{\left\langle #1,#2 \right\rangle}

\DeclareMathOperator{\rank}{rank}
\DeclareMathOperator{\tr}{tr}
\DeclareMathOperator{\diag}{diag}
\DeclareMathOperator{\range}{range}
\DeclareMathOperator{\Span}{span}
\DeclareMathOperator{\dist}{dist}
\DeclareMathOperator{\spec}{spec}
\DeclareMathOperator{\argmax}{arg\,max}
\DeclareMathOperator{\argmin}{arg\,min}

% ============================================================
% MATLAB LISTINGS STYLE
% ============================================================

\lstdefinestyle{matlabstyle}{
    language=Matlab,
    basicstyle=\ttfamily\small,
    keywordstyle=\color{MidnightBlue}\bfseries,
    commentstyle=\color{ForestGreen},
    stringstyle=\color{BrickRed},
    numbers=left,
    numberstyle=\tiny\color{gray},
    numbersep=8pt,
    stepnumber=1,
    breaklines=true,
    breakatwhitespace=false,
    frame=single,
    framerule=0.4pt,
    rulecolor=\color{gray!50},
    backgroundcolor=\color{gray!4},
    showstringspaces=false,
    tabsize=4,
    keepspaces=true,
    columns=fullflexible,
    captionpos=b
}

\lstset{style=matlabstyle}

% ============================================================
% DOCUMENT INFORMATION
% ============================================================

\title{
    \textbf{Topics in Scientific Computing}\\[0.5em]
    \large
    Subspace Geometry, the Courant--Fischer Theorem,\\
    Weyl's Inequality, and the Davis--Kahan Theorem
}

\author{Biajid Chowdhury}

\date{\today}

% ============================================================
% BEGIN DOCUMENT
% ============================================================

\begin{document}

\maketitle

% ============================================================
% ABSTRACT
% ============================================================

\begin{abstract}

Eigenvalues, eigenvectors, and invariant subspaces play a
fundamental role in scientific computing, numerical linear
algebra, and data analysis. An important question in these
fields is how reliably such spectral information can be
computed and how it changes when the underlying matrix
is subjected to small perturbations.

This note develops several fundamental ideas connecting
the geometry of subspaces, variational characterizations
of eigenvalues, and perturbation theory for real symmetric
matrices. We begin with a brief review of essential concepts
from linear algebra, including the spectral theorem,
the Rayleigh quotient, the power method, singular value
decomposition, orthogonal projections, and metric spaces.

We then introduce a natural notion of distance between
subspaces through their orthogonal projection matrices.
We establish that this distance satisfies the axioms of
a metric and investigate its geometric interpretation
through principal angles. In particular, we explain how
the singular values of a matrix formed from two orthonormal
bases characterize the angular separation between their
corresponding subspaces.

Building on these geometric foundations, we develop the
Courant--Fischer min--max and max--min characterizations
of eigenvalues and examine their connection to numerical
computation. We subsequently establish Weyl's eigenvalue
perturbation inequality and the Davis--Kahan
$\sin\Theta$ theorem, demonstrating how the stability of
eigenvalues and invariant subspaces depends on the
magnitude of perturbations and the separation of
eigenvalues.

Throughout the discussion, emphasis is placed on
mathematical intuition, detailed derivations, and
step-by-step proofs. Illustrative examples and MATLAB
implementations are included to connect the theoretical
results with numerical computation.

The primary objective is to provide a self-contained
and accessible treatment that enables students with
a basic background in linear algebra to develop a
clear understanding of these fundamental results and
their significance in scientific computing.

\end{abstract}

% ============================================================
% TABLE OF CONTENTS
% ============================================================

\newpage
\tableofcontents

\newpage

% ============================================================
% SECTION 1: INTRODUCTION
% ============================================================


% ============================================================
% SECTION 1: INTRODUCTION
% ============================================================
%```latex
% ============================================================
% SECTION 1: INTRODUCTION
% ============================================================
    

\begin{comment}
    

\section{Introduction}
\label{sec:introduction}

Nature loves symmetry!
From the arrangement of atoms in a crystal to the laws
governing physical motion, symmetry appears repeatedly
as a fundamental organizing principle. In mathematics,
symmetry is more than an appealing geometric property:
it often reveals hidden structure, simplifies complicated
problems, and makes otherwise difficult computations
considerably more manageable.

Many mathematical models arising in natural science,
engineering, and scientific computing possess some form
of symmetry. For example, symmetric matrices appear
naturally in covariance analysis, quadratic optimization,
mechanical systems, and the discretization of many
self-adjoint differential operators. Even when a problem
is not symmetric in its original formulation, it may
sometimes be reformulated to expose a related symmetric
structure.

An important example arises in the study of a general
real matrix $B$. Although $B$ need not be symmetric,
the matrices
\[
B^TB
\qquad\text{and}\qquad
BB^T
\]
are always symmetric and positive semidefinite.
Their eigenvalues and eigenvectors are closely related
to the singular values and singular vectors of $B$.
This observation provides one of the fundamental
connections between eigenvalue problems and the
singular value decomposition.

Of course, transforming a problem into a symmetric
one does not always make it computationally preferable.
For example, forming $B^TB$ may worsen numerical
conditioning. Nevertheless, such transformations
illustrate the importance of symmetric structure
in understanding more general matrix problems.

\subsection{Why Are Symmetric Matrices Special?}

Consider a real square matrix
\[
A\in\R^{n\times n}.
\]

We call $A$ symmetric if
\[
\boxed{A=A^T.}
\]

At first glance, this condition appears to impose
only a simple relationship between matrix entries:
\[
a_{ij}=a_{ji}.
\]

However, this seemingly elementary requirement has
remarkable consequences.

The spectral theorem guarantees that every real
symmetric matrix possesses real eigenvalues and
an orthonormal basis of eigenvectors. Consequently,
such a matrix can always be diagonalized by an
orthogonal transformation:
\[
\boxed{
A=V\Lambda V^T,
}
\]
where
\[
V^TV=VV^T=I
\]
and $\Lambda$ is a real diagonal matrix containing
the eigenvalues of $A$.

This decomposition is especially valuable because
orthogonal transformations preserve Euclidean lengths
and angles. In particular,
\[
\norm{Vx}_2=\norm{x}_2
\qquad
\text{for every }x\in\R^n.
\]

Thus, the spectral decomposition does not merely
express a matrix in a simpler algebraic form.
It identifies a coordinate system in which the
action of the matrix becomes particularly transparent.

Indeed, if
\[
A=V\Lambda V^T,
\]
then applying $A$ to a vector can be interpreted
through three operations: expressing the vector
in the eigenvector coordinates, scaling its
components by the corresponding eigenvalues,
and transforming the result back into the
original coordinate system.

This geometric interpretation is one reason
symmetric matrices play such a central role
in numerical linear algebra.

Moreover, symmetry allows us to develop powerful
mathematical tools for estimating eigenvalues,
studying invariant subspaces, and understanding
the effects of numerical errors.

The Rayleigh quotient, the Courant--Fischer theorem,
Weyl's inequality, and the Davis--Kahan theorem
are among the most important results arising
from this structure.

\subsection{From Eigenvalue Computation to Stability}

Suppose we wish to solve the eigenvalue problem
\[
\boxed{
Av=\lambda v,
}
\]
where $v\neq0$.

At a purely algebraic level, the problem asks us
to determine the scalars $\lambda$ and nonzero
vectors $v$ satisfying this equation.

However, scientific computation introduces
additional questions.

In practical applications, the entries of a matrix
may come from experimental measurements, numerical
approximations, statistical estimates, or the
discretization of continuous mathematical models.

Consequently, the matrix available to us may
not be the exact matrix $A$, but rather a
perturbed version
\[
\boxed{
\widetilde A=A+E,
}
\]
where $E$ represents measurement errors,
rounding errors, approximation errors, or
other perturbations.

Even if $\norm{E}_2$ is very small, we must ask:

\begin{quote}
\emph{
If two symmetric matrices are nearly identical,
must their eigenvalues and eigenvectors also
be nearly identical?
}
\end{quote}

This question is fundamental to the reliability
of numerical computation.

Interestingly, the answer is different for
eigenvalues and eigenvectors.

For real symmetric matrices, eigenvalues exhibit
a remarkable stability property. As we shall
establish using Weyl's inequality,
\[
\boxed{
\abs{
\lambda_i(A+E)-\lambda_i(A)
}
\leq\norm{E}_2,
}
\]
where the eigenvalues of both matrices are
arranged in the same order.

Thus, a small symmetric perturbation cannot
cause an arbitrarily large change in an
individual ordered eigenvalue.

The corresponding eigenvectors, however,
require a more delicate analysis.

When two eigenvalues are very close,
even a small perturbation may produce a
substantial change in their associated
eigenvectors. Conversely, when an eigenvalue
is well separated from the rest of the
spectrum, its corresponding eigenspace
tends to be considerably more stable.

This distinction leads naturally to the
study of spectral gaps and invariant
subspace perturbation theory.

\subsection{Why Do We Study Subspaces Rather Than
Individual Eigenvectors?}

An eigenvector is not uniquely determined,
even when its corresponding eigenvalue is simple.

Indeed, if
\[
Av=\lambda v,
\]
then
\[
A(cv)=\lambda(cv)
\]
for every nonzero scalar $c$.

In particular, both $v$ and $-v$ are
eigenvectors associated with the same eigenvalue.

Although these two vectors point in opposite
directions, they span exactly the same
one-dimensional subspace:
\[
\Span\{v\}=\Span\{-v\}.
\]

Consequently, directly comparing eigenvectors
using ordinary Euclidean distance may produce
a misleading description of their geometric
relationship.

For example, if $\norm{v}_2=1$, then
\[
\norm{v-(-v)}_2=2,
\]
even though the two vectors determine
the same eigenspace.

This observation suggests that, when studying
eigenvector stability, the more natural
mathematical objects are often the invariant
subspaces generated by eigenvectors rather
than the individual vectors themselves.

But this introduces another question:

\begin{quote}
\emph{
How should we define the distance between
two subspaces?
}
\end{quote}

Unlike ordinary vectors, subspaces cannot
simply be subtracted from one another.

We therefore need a suitable mathematical
representation that allows us to compare
them in a meaningful and basis-independent way.

As we shall see, orthogonal projection matrices
provide precisely such a representation.

If $\mathcal U$ and $\mathcal W$ are subspaces
with corresponding orthogonal projectors
$P_{\mathcal U}$ and $P_{\mathcal W}$,
we may define
\[
\boxed{
\dist(\mathcal U,\mathcal W)
=
\norm{
P_{\mathcal U}-P_{\mathcal W}
}_2.
}
\]

We shall prove that this definition satisfies
the axioms of a metric and investigate its
relationship with the principal angles
between subspaces.

In particular, for two subspaces of the
same dimension, their distance can be
expressed as
\[
\boxed{
\dist(\mathcal U,\mathcal W)
=
\sin\theta_{\max},
}
\]
where $\theta_{\max}$ denotes their largest
principal angle.

This geometric framework will become
essential when we study the Davis--Kahan
theorem.

\subsection{The Variational Point of View}

There is another important perspective
that connects eigenvalues with optimization.

For a real symmetric matrix $A$, consider
the Rayleigh quotient
\[
\boxed{
R_A(x)=\frac{x^TAx}{x^Tx},
\qquad x\neq0.
}
\]

The Rayleigh quotient associates a real
number with every nonzero vector.

Its extreme values are precisely the
smallest and largest eigenvalues of $A$.

More remarkably, intermediate eigenvalues
can also be characterized through
optimization over appropriately chosen
subspaces.

This is the central idea of the
Courant--Fischer theorem.

The theorem transforms an algebraic
eigenvalue problem into a geometric
optimization problem involving the
Rayleigh quotient and the dimensions
of subspaces.

Besides providing a deeper understanding
of eigenvalues, this variational
characterization becomes a powerful tool
for establishing perturbation bounds.

In particular, we shall use the
Courant--Fischer theorem to derive
Weyl's inequality, thereby connecting
the geometry of the Rayleigh quotient
with the stability of eigenvalues.

\subsection{Objectives and Organization of the Note}

The purpose of this note is to develop
a clear and connected understanding
of several fundamental results in
symmetric matrix theory and scientific
computation.

Rather than presenting these results
as isolated theorems, we shall build
them progressively, emphasizing the
mathematical ideas that connect them.

We begin with a review of the necessary
background, including eigenvalues and
eigenvectors, the spectral theorem,
the Rayleigh quotient, the power method,
singular value decomposition, matrix
norms, metric spaces, and orthogonal
projections.

We then develop the geometry of
subspaces, introduce a rigorous
definition of subspace distance,
establish its metric properties,
and investigate its relationship
with principal angles and singular
values.

Next, we study the Courant--Fischer
min--max and max--min principles,
providing detailed proofs and
explaining their variational meaning.

With these foundations established,
we turn to eigenvalue perturbation
theory and prove Weyl's inequality.

Finally, we investigate the
Davis--Kahan $\sin\Theta$ theorem,
which describes how invariant
subspaces change under symmetric
perturbations and explains the
fundamental role of spectral gaps.

Throughout the note, the theoretical
discussion will be complemented
by illustrative examples and MATLAB
implementations whenever they help
clarify the underlying mathematics.

Our guiding principle is simple:
a mathematical result becomes much
more useful when we understand not
only what it states, but also why
it is true, how it is derived, and
when it can be applied.

\end{comment}

%```latex
% ============================================================
% SECTION 1: INTRODUCTION
% ============================================================

\section{Introduction}
\label{sec:introduction}

Nature loves symmetry!
From the arrangement of atoms in crystalline structures
to the laws governing physical phenomena, symmetry appears
repeatedly as an organizing principle of the natural world.
What makes symmetry particularly interesting is that its
importance extends far beyond geometric beauty. In
mathematics, the presence of symmetry often reveals
structure that would otherwise remain hidden, allowing
complicated problems to be understood through simpler
and more elegant mathematical descriptions.

The importance of symmetry becomes especially apparent
in scientific computation. Many mathematical models
arising in physics, engineering, statistics, and applied
mathematics possess underlying symmetric structures.
Others, although not symmetric in their original form,
can sometimes be reformulated in ways that allow
such structures to be exploited. In either case,
recognizing symmetry can substantially simplify the
analysis of a problem and, when used appropriately,
improve the efficiency and reliability of its
numerical solution.

Among the mathematical objects that exhibit these properties, real symmetric matrices occupy a
particularly important position. They arise naturally
in numerous scientific applications, including
mechanical and structural analysis, optimization,
statistical modeling, and the numerical approximation
of differential equations. Their special algebraic
and geometric properties provide a powerful foundation
for understanding a wide range of computational
problems.

One of the most significant benefits of symmetry
is that it makes the spectral structure of a matrix
especially transparent. Questions concerning
eigenvalues, eigenvectors, and invariant subspaces
become more accessible, both theoretically and
computationally. This allows us to develop powerful
connections between linear algebra, geometry,
optimization, and numerical analysis.

However, obtaining a numerical solution is only
one part of scientific computation. Equally important
is understanding how reliable that solution is.

In practical applications, mathematical models
are rarely represented with perfect accuracy.
Experimental measurements contain uncertainty,
statistical estimates depend on finite samples,
continuous models must often be approximated by
discrete systems, and numerical computations are
performed using finite-precision arithmetic.

Consequently, the matrix that we analyze may differ
slightly from the matrix that describes the ideal
mathematical problem. Such differences raise an
important question:

\begin{quote}
\emph{
If a matrix is changed by a small amount,
how much can its spectral properties change?
}
\end{quote}

At first, it may seem reasonable to expect that
a small change in a matrix should produce only
small changes in all its eigenvalues and
eigenvectors. However, the situation is more
subtle than this naive intuition suggests.

For symmetric matrices, eigenvalues possess
remarkable stability properties. Eigenvectors
and invariant subspaces, on the other hand,
may respond quite differently to perturbations.
Their behavior depends not only on the size
of a perturbation but also on the separation
between the relevant eigenvalues.

This distinction is fundamental in scientific
computing. A numerical algorithm may produce
highly accurate eigenvalues while the associated
eigenvectors are considerably more sensitive
to small changes in the underlying data.
Understanding this phenomenon requires us
to examine not only the algebraic properties
of matrices but also the geometry of their
invariant subspaces.

These considerations lead to several closely
related mathematical questions. How should
we measure the distance between two subspaces?
How can eigenvalues be characterized through
optimization principles? How much can the
eigenvalues of a symmetric matrix change under
a perturbation? And under what conditions
will its invariant subspaces remain stable?

The purpose of this note is to develop a
connected understanding of these questions
through four fundamental mathematical topics:
subspace geometry and distance, the
Courant--Fischer theorem, Weyl's eigenvalue
perturbation inequality, and the
Davis--Kahan $\sin\Theta$ theorem.

Although these topics are sometimes studied
separately, they are closely related.
The geometry of subspaces provides a meaningful
way to compare eigenspaces, while variational
principles connect eigenvalues with optimization.
These ideas, in turn, establish the foundations
for analyzing the stability of spectral
information under perturbations.

Our approach throughout this note will be
deliberately elementary and constructive.
We shall begin with a review of the necessary
mathematical background and then develop
the principal results through detailed
derivations, geometric intuition, and
step-by-step proofs.

Whenever appropriate, MATLAB implementations
and numerical examples will be used to
illustrate how the theoretical conclusions
appear in actual computations.

The intention is not merely to collect
several classical theorems in one place,
but to explain the mathematical reasoning
that connects them. We hope that readers
with a basic knowledge of linear algebra and numerical computation
will be able to follow the development
without repeatedly consulting external
references.

\begin{comment}
Ultimately, our objective is to demonstrate
how symmetry, geometry, and variational
principles work together to provide a
deeper understanding of numerical computation
and the stability of its results.
\end{comment}

%```latex
% ============================================================
% SECTION 2: MATHEMATICAL BACKGROUND
% ============================================================

\section{Mathematical Background}
\label{sec:background}

Before developing the principal results of this note,
we briefly review several fundamental concepts from
linear algebra, geometry, and numerical computation.
Readers who are already familiar with these concepts
may skip this section.

Throughout the note, unless otherwise stated, all
matrices and vector spaces are assumed to be real.

% ============================================================
% 2.1 EIGENVALUES, EIGENVECTORS, AND INVARIANT SUBSPACES
% ============================================================

\subsection{Eigenvalues, Eigenvectors, and Invariant Subspaces}
\label{subsec:eigen_background}

Let
\[
A\in\R^{n\times n}
\]
be a square matrix.

A nonzero vector $v\in\R^n$ is called an
\emph{eigenvector} of $A$ if there exists a scalar
$\lambda\in\R$ such that
\[
\boxed{Av=\lambda v.}
\]

The scalar $\lambda$ is called the corresponding
\emph{eigenvalue}.

Geometrically, an eigenvector identifies a direction
that is preserved by the linear transformation
represented by matrix $A$. The matrix may stretch, compress,
or reverse this direction, but it does not move the
vector outside the line that it spans.

The scalar $\lambda$ determines the corresponding
scaling factor.

For example, if
\[
Av=3v,
\]
then the action of $A$ stretches $v$ by a factor
of three.

On the other hand, if
\[
Av=-2v,
\]
then the length is multiplied by two while
the direction is reversed.

\begin{definition}[Eigenspace]
For an eigenvalue $\lambda$ of $A$, the associated
eigenspace is
\[
\boxed{
\mathcal E_{\lambda}
=
\ker(A-\lambda I).
}
\]

Equivalently,
\[
\mathcal E_{\lambda}
=
\{x\in\R^n:Ax=\lambda x\}.
\]
\end{definition}

Notice that an eigenspace contains the zero vector,
whereas the definition of an eigenvector excludes it.

More generally, we may consider eigenspaces as the subspaces that are
preserved by the action of a matrix.

\begin{definition}[Invariant Subspace]
A linear subspace $\mathcal S\subseteq\R^n$ is called
an invariant subspace of $A$ if
\[
\boxed{
A\mathcal S\subseteq\mathcal S.
}
\]

In other words, whenever $x\in\mathcal S$, we also have
\[
Ax\in\mathcal S.
\]
\end{definition}

Every eigenspace is an invariant subspace.
However, invariant subspaces need not consist
entirely of eigenvectors associated with a single
eigenvalue.

For instance, if $v_1$ and $v_2$ are eigenvectors
of $A$, possibly associated with different eigenvalues,
then
\[
\mathcal S=\Span\{v_1,v_2\}
\]
is an invariant subspace.

Indeed, for any
\[
x=c_1v_1+c_2v_2,
\]
we have
\[
Ax=c_1\lambda_1v_1+c_2\lambda_2v_2,
\]
which remains in $\mathcal S$.

This distinction between eigenvectors and invariant
subspaces will become important when studying
the Davis--Kahan theorem.

% ============================================================
% 2.2 SYMMETRIC MATRICES AND SPECTRAL THEOREM
% ============================================================

\subsection{Symmetric Matrices and the Spectral Theorem}
\label{subsec:spectral_background}

A real square matrix $A$ is symmetric if
\[
\boxed{A=A^T.}
\]

Symmetric matrices possess several remarkable
properties that distinguish them from general
square matrices.

In particular, the spectral theorem provides
a complete description of their eigen-structure.

\begin{theorem}[Spectral Theorem]
Let $A\in\R^{n\times n}$ be symmetric.

Then all eigenvalues of $A$ are real, and
$\R^n$ possesses an orthonormal basis consisting
of eigenvectors of $A$.

Consequently, there exists an orthogonal matrix
$V$ and a real diagonal matrix $\Lambda$ such that
\[
\boxed{
A=V\Lambda V^T.
}
\]
\end{theorem}

Here,
\[
V=
\begin{bmatrix}
v_1 & v_2 & \cdots & v_n
\end{bmatrix}
\]
contains the orthonormal eigenvectors of $A$, and
\[
\Lambda=
\diag(\lambda_1,\lambda_2,\ldots,\lambda_n)
\]
contains the corresponding eigenvalues.

Throughout this note, we adopt the convention
\[
\boxed{
\lambda_1\leq\lambda_2\leq\cdots\leq\lambda_n.
}
\]

Thus, $\lambda_1$ denotes the smallest eigenvalue
and $\lambda_n$ denotes the largest.

The orthogonality of $V$ means that
\[
V^TV=VV^T=I.
\]

In particular,
\[
v_i^Tv_j=
\begin{cases}
1, & i=j,\\
0, & i\neq j.
\end{cases}
\]

\paragraph{Orthogonality of eigenvectors.}

One useful property is that eigenvectors associated
with distinct eigenvalues of a symmetric matrix
must be orthogonal.

Indeed, suppose
\[
Av_i=\lambda_i v_i,
\qquad
Av_j=\lambda_j v_j.
\]

Since $A=A^T$,
\[
v_i^TAv_j=(Av_i)^Tv_j.
\]

Substituting the eigenvalue equations gives
\[
\lambda_j v_i^Tv_j
=
\lambda_i v_i^Tv_j.
\]

Therefore,
\[
(\lambda_i-\lambda_j)v_i^Tv_j=0.
\]

If $\lambda_i\neq\lambda_j$, we conclude that
\[
\boxed{v_i^Tv_j=0.}
\]

This simple argument illustrates how strongly
symmetry constrains the geometry of eigenvectors.

\paragraph{Positive semidefinite matrices.}

A symmetric matrix $A$ is called positive
semidefinite if
\[
x^TAx\geq0
\qquad
\text{for every }x\in\R^n.
\]

For a real symmetric matrix, this condition is
equivalent to requiring that all its eigenvalues
be nonnegative.

Similarly, $A$ is positive definite if
\[
x^TAx>0
\qquad
\text{for every }x\neq0,
\]
which is equivalent to all its eigenvalues
being strictly positive.

Positive definite and positive semidefinite
matrices occur frequently in optimization,
covariance analysis, and numerical methods
for differential equations.

% ============================================================
% 2.3 RAYLEIGH QUOTIENT
% ============================================================
\begin{comment}
\subsection{The Rayleigh Quotient}
\label{subsec:rayleigh_background}

The Rayleigh quotient provides an important
connection between eigenvalues and optimization.

\begin{definition}[Rayleigh Quotient]
For a real symmetric matrix $A$ and a nonzero
vector $x\in\R^n$, the Rayleigh quotient is
defined by
\[
\boxed{
R_A(x)=\frac{x^TAx}{x^Tx}.
}
\]
\end{definition}

Notice that the quotient is unchanged when
$x$ is multiplied by any nonzero scalar $c$:
\[
R_A(cx)
=
\frac{c^2x^TAx}{c^2x^Tx}
=
R_A(x).
\]

Therefore,
\[
\boxed{R_A(cx)=R_A(x),\qquad c\neq0.}
\]

This means that the Rayleigh quotient depends
on the direction of a vector rather than
its length.

If $x$ is an eigenvector associated with
$\lambda$, then
\[
Ax=\lambda x.
\]

Consequently,
\[
R_A(x)
=
\frac{x^T(\lambda x)}{x^Tx}
=
\lambda.
\]

Thus,
\[
\boxed{
R_A(v_i)=\lambda_i.
}
\]

The Rayleigh quotient also provides information
when $x$ is not an eigenvector.

To see this, use the spectral theorem to express
an arbitrary nonzero vector as
\[
x=\sum_{i=1}^{n}c_i v_i.
\]

Since the eigenvectors are orthonormal,
\[
x^Tx=\sum_{i=1}^{n}c_i^2.
\]

Similarly,
\[
x^TAx=\sum_{i=1}^{n}\lambda_i c_i^2.
\]

Therefore,
\[
\boxed{
R_A(x)
=
\frac{
\displaystyle\sum_{i=1}^{n}\lambda_i c_i^2
}{
\displaystyle\sum_{i=1}^{n}c_i^2
}.
}
\]

This expression is a weighted average of the
eigenvalues, where the weights are nonnegative.

Because
\[
\lambda_1\leq\lambda_i\leq\lambda_n,
\]
we obtain
\[
\boxed{
\lambda_1\leq R_A(x)\leq\lambda_n.
}
\]

Moreover, equality is attained by the
corresponding eigenvectors.

Hence,
\[
\boxed{
\lambda_1=\min_{x\neq0}R_A(x),
\qquad
\lambda_n=\max_{x\neq0}R_A(x).
}
\]

These are the simplest variational
characterizations of eigenvalues.

The Courant--Fischer theorem extends this
principle to every intermediate eigenvalue.

We shall study that extension and its proof
in detail in a later section.

\end{comment}


% ============================================================
% 2.3 RAYLEIGH QUOTIENT
% ============================================================

\subsection{The Rayleigh Quotient}
\label{subsec:rayleigh_background}

The Rayleigh quotient is one of the most important
tools in spectral theory. It establishes a connection
between eigenvalues, eigenvectors, and optimization,
allowing us to study eigenvalue problems from a
variational point of view.

\begin{definition}[Rayleigh Quotient]
Let $A\in\R^{n\times n}$ be a real symmetric matrix.
For any nonzero vector $x\in\R^n$, the
\emph{Rayleigh quotient} of $A$ at $x$ is defined by
\[
\boxed{
R_A(x)=\frac{x^TAx}{x^Tx}.
}
\]
\end{definition}

Although the quotient itself can be defined for
any real square matrix, the eigenvalue bounds
developed below depend fundamentally on
the symmetry of $A$.

\paragraph{Homogeneity of the Rayleigh quotient.}

An important property is that the Rayleigh
quotient remains unchanged when its argument
is multiplied by a nonzero scalar.

Indeed, for any $c\neq0$,
\[
\begin{aligned}
R_A(cx)
&=\frac{(cx)^TA(cx)}{(cx)^T(cx)}\\[4pt]
&=\frac{c^2x^TAx}{c^2x^Tx}\\[4pt]
&=R_A(x).
\end{aligned}
\]

Therefore,
\[
\boxed{
R_A(cx)=R_A(x),\qquad c\neq0.
}
\]

This means that the Rayleigh quotient depends
only on the direction of a vector, not its
length. Consequently, when studying its
extreme values, we may restrict our attention
to unit vectors.

\paragraph{Evaluation at an eigenvector.}

Suppose $v_i$ is an eigenvector of $A$
associated with eigenvalue $\lambda_i$.

Since
\[
Av_i=\lambda_i v_i,
\]
we obtain
\[
\begin{aligned}
R_A(v_i)
&=\frac{v_i^TAv_i}{v_i^Tv_i}\\[4pt]
&=\frac{v_i^T(\lambda_i v_i)}{v_i^Tv_i}\\[4pt]
&=\lambda_i.
\end{aligned}
\]

Hence,
\[
\boxed{
R_A(v_i)=\lambda_i.
}
\]

This observation is not restricted to
symmetric matrices: whenever a real
matrix possesses a real eigenvector
with a real eigenvalue, the corresponding
Rayleigh quotient returns that eigenvalue.

However, the more powerful property of
the Rayleigh quotient emerges when
we evaluate it at an arbitrary vector.

% ------------------------------------------------------------
% BOUNDS ON THE RAYLEIGH QUOTIENT
% ------------------------------------------------------------

\paragraph{Bounds on the Rayleigh quotient.}

We now establish one of its fundamental
properties for real symmetric matrices.

\begin{proposition}[Rayleigh Quotient Bounds]
Let $A\in\R^{n\times n}$ be symmetric,
with eigenvalues ordered as
\[
\lambda_1\leq\lambda_2\leq\cdots\leq\lambda_n.
\]

Then, for every nonzero $x\in\R^n$,
\[
\boxed{
\lambda_1
\leq
R_A(x)
\leq
\lambda_n.
}
\]

Furthermore,
\[
\boxed{
\lambda_1=\min_{x\neq0}R_A(x),
\qquad
\lambda_n=\max_{x\neq0}R_A(x).
}
\]
\end{proposition}

\begin{proof}

Since $A$ is real and symmetric, the
spectral theorem guarantees an
orthonormal basis of eigenvectors
\[
v_1,v_2,\ldots,v_n.
\]

Therefore, any nonzero vector
$x\in\R^n$ can be expressed as
\[
\boxed{
x=\sum_{i=1}^{n}c_i v_i,
}
\]
where the coefficients $c_i$ are real
and not all zero.

We first calculate the denominator
of the Rayleigh quotient.

Using the orthonormality of the
eigenvectors,
\[
\begin{aligned}
x^Tx
&=
\left(\sum_{i=1}^{n}c_i v_i\right)^T
\left(\sum_{j=1}^{n}c_j v_j\right)\\[4pt]
&=
\sum_{i=1}^{n}
\sum_{j=1}^{n}
c_i c_j v_i^Tv_j\\[4pt]
&=
\sum_{i=1}^{n}c_i^2.
\end{aligned}
\]

Thus,
\[
\boxed{
x^Tx=\sum_{i=1}^{n}c_i^2>0.
}
\]

Next, consider the numerator.

Since
\[
Av_i=\lambda_i v_i,
\]
we have
\[
\begin{aligned}
x^TAx
&=
\left(\sum_{i=1}^{n}c_i v_i\right)^T
A\left(\sum_{j=1}^{n}c_j v_j\right)\\[4pt]
&=
\left(\sum_{i=1}^{n}c_i v_i\right)^T
\left(\sum_{j=1}^{n}
c_j\lambda_j v_j\right)\\[4pt]
&=
\sum_{i=1}^{n}
\sum_{j=1}^{n}
c_i c_j\lambda_j v_i^Tv_j\\[4pt]
&=
\sum_{i=1}^{n}\lambda_i c_i^2.
\end{aligned}
\]

Consequently,
\[
\boxed{
R_A(x)
=
\frac{
\displaystyle\sum_{i=1}^{n}\lambda_i c_i^2
}{
\displaystyle\sum_{i=1}^{n}c_i^2
}.
}
\]

We now establish the lower and upper
bounds separately.

\medskip

\noindent
\textbf{Lower bound.}

By our ordering convention,
\[
\lambda_1\leq\lambda_i
\qquad
\text{for every }i.
\]

Because $c_i^2\geq0$, multiplying
the inequality by $c_i^2$ preserves
its direction:
\[
\lambda_1c_i^2\leq\lambda_i c_i^2.
\]

Summing over all indices gives
\[
\sum_{i=1}^{n}\lambda_1c_i^2
\leq
\sum_{i=1}^{n}\lambda_i c_i^2.
\]

Since $\lambda_1$ is constant,
\[
\lambda_1\sum_{i=1}^{n}c_i^2
\leq
\sum_{i=1}^{n}\lambda_i c_i^2.
\]

Now divide both sides by
\[
\sum_{i=1}^{n}c_i^2>0.
\]

We obtain
\[
\lambda_1
\leq
\frac{
\displaystyle\sum_{i=1}^{n}\lambda_i c_i^2
}{
\displaystyle\sum_{i=1}^{n}c_i^2
}.
\]

Therefore,
\[
\boxed{
\lambda_1\leq R_A(x).
}
\]

\medskip

\noindent
\textbf{Upper bound.}

Similarly, since
\[
\lambda_i\leq\lambda_n
\qquad
\text{for every }i,
\]
we have
\[
\lambda_i c_i^2
\leq
\lambda_n c_i^2.
\]

Summing over all indices,
\[
\sum_{i=1}^{n}\lambda_i c_i^2
\leq
\lambda_n\sum_{i=1}^{n}c_i^2.
\]

Again, dividing by the strictly
positive denominator gives
\[
\frac{
\displaystyle\sum_{i=1}^{n}\lambda_i c_i^2
}{
\displaystyle\sum_{i=1}^{n}c_i^2
}
\leq
\lambda_n.
\]

Hence,
\[
\boxed{
R_A(x)\leq\lambda_n.
}
\]

Combining the two inequalities,
we conclude that
\[
\boxed{
\lambda_1
\leq
R_A(x)
\leq
\lambda_n
\qquad
\text{for every }x\neq0.
}
\]

Finally, we already know that
\[
R_A(v_1)=\lambda_1,
\qquad
R_A(v_n)=\lambda_n.
\]

Thus, both bounds are attained,
and we obtain
\[
\boxed{
\lambda_1=\min_{x\neq0}R_A(x),
\qquad
\lambda_n=\max_{x\neq0}R_A(x).
}
\]

\end{proof}

% ------------------------------------------------------------
% GEOMETRIC INTERPRETATION
% ------------------------------------------------------------

\paragraph{Interpretation as a weighted average.}

The preceding proof also reveals
an important geometric interpretation.

Define
\[
\omega_i
=
\frac{c_i^2}{
\displaystyle\sum_{j=1}^{n}c_j^2
}.
\]

Notice that
\[
\omega_i\geq0,
\qquad
\sum_{i=1}^{n}\omega_i=1.
\]

Therefore, the Rayleigh quotient
can be written as
\[
\boxed{
R_A(x)=\sum_{i=1}^{n}\omega_i\lambda_i.
}
\]

Thus, for a real symmetric matrix,
the Rayleigh quotient is a
\emph{convex combination} of its
eigenvalues.

This provides another explanation
for the bounds we just established:
a weighted average of real numbers
with nonnegative weights summing
to one cannot exceed their maximum
or fall below their minimum.

The coefficients $c_i$ describe
how strongly $x$ is aligned with
the corresponding eigenvectors.
Their squares determine how much
each eigenvalue contributes to
the Rayleigh quotient.

% ------------------------------------------------------------
% IMPORTANT WARNING
% ------------------------------------------------------------

\begin{remark}[Why Symmetry Is Essential]
\label{rem:rayleigh_symmetry}

A common mistake is to assume that
\[
\lambda_{\min}(A)
\leq R_A(x)\leq
\lambda_{\max}(A)
\]
holds for every real square matrix.

\textbf{This is not true in general.}

Our proof depends on the existence
of an orthonormal eigenbasis, which
is guaranteed for real symmetric
matrices but not for arbitrary
real square matrices.

To see what can go wrong, consider
the nonsymmetric matrix
\[
A=
\begin{bmatrix}
1&4\\
0&2
\end{bmatrix}.
\]

Since $A$ is upper triangular,
its eigenvalues are
\[
\lambda_1=1,
\qquad
\lambda_2=2.
\]

Now choose
\[
x=
\begin{bmatrix}
1\\
1
\end{bmatrix}.
\]

We calculate
\[
Ax=
\begin{bmatrix}
5\\
2
\end{bmatrix}.
\]

Therefore,
\[
x^TAx
=
\begin{bmatrix}1&1\end{bmatrix}
\begin{bmatrix}5\\2\end{bmatrix}
=7,
\]
while
\[
x^Tx=2.
\]

Hence,
\[
\boxed{
R_A(x)=\frac72=3.5.
}
\]

But the largest eigenvalue of $A$ is
\[
\lambda_{\max}(A)=2.
\]

Consequently,
\[
\boxed{
R_A(x)=3.5>2=\lambda_{\max}(A).
}
\]

The familiar Rayleigh quotient
bound has failed, even though
this nonsymmetric matrix has
two real eigenvalues.

For a general real square matrix,
we may write
\[
A=
\underbrace{\frac{A+A^T}{2}}_{S}
+
\underbrace{\frac{A-A^T}{2}}_{K},
\]
where $S$ is symmetric and
$K$ is skew-symmetric.

Since
\[
x^TKx=0
\]
for every real vector $x$, we have
\[
x^TAx=x^TSx.
\]

Therefore,
\[
\boxed{
R_A(x)=R_S(x),
\qquad
S=\frac{A+A^T}{2}.
}
\]

In particular,
\[
\boxed{
\lambda_{\min}(S)
\leq R_A(x)\leq
\lambda_{\max}(S).
}
\]

Thus, for a general real matrix,
the Rayleigh quotient is bounded
by the extreme eigenvalues of
its \emph{symmetric part},
not necessarily by the
eigenvalues of the original matrix.

This distinction is important
and should not be overlooked.

\end{remark}

% ------------------------------------------------------------
% CONNECTION TO COURANT--FISCHER
% ------------------------------------------------------------

\paragraph{Connection to the Courant--Fischer theorem.}

We have established that the smallest
and largest eigenvalues of a real
symmetric matrix can be characterized
through the optimization of its
Rayleigh quotient:
\[
\boxed{
\lambda_1=\min_{x\neq0}R_A(x),
\qquad
\lambda_n=\max_{x\neq0}R_A(x).
}
\]

A natural question now arises:

\begin{quote}
\emph{
Can the intermediate eigenvalues
also be characterized through
optimization?
}
\end{quote}

The answer is yes.

The Courant--Fischer theorem extends
these elementary variational principles
to every eigenvalue by considering
the Rayleigh quotient over subspaces
of prescribed dimensions.

We shall develop both the min--max
and max--min formulations of this
theorem, together with their proofs,
in a later section.
% ------------------------------------------------------------
% ============================================================
% 2.4 POWER METHOD
% ============================================================

\subsection{The Power Method}
\label{subsec:power_background}

The power method is one of the simplest iterative
algorithms for approximating an eigenvalue
and its corresponding eigenvector.

Given an initial nonzero vector $x_0$, the method
repeatedly applies the matrix $A$ and normalizes
the result:
\[
\boxed{
x_{k+1}
=
\frac{Ax_k}{\norm{Ax_k}_2},
}
\]
provided $Ax_k\neq0$.

To understand the underlying principle, suppose
$A$ is symmetric and
\[
x_0=\sum_{i=1}^{n}c_i v_i.
\]

After applying $A$ repeatedly,
\[
A^k x_0
=
\sum_{i=1}^{n}c_i\lambda_i^k v_i.
\]

Suppose one eigenvalue $\lambda_*$ has
strictly larger absolute value than all
the others:
\[
|\lambda_*|>|\lambda_j|
\qquad
\text{for every other eigenvalue }\lambda_j.
\]

If the initial vector has a nonzero component
along the corresponding eigenvector $v_*$,
then repeated multiplication increasingly
emphasizes that eigendirection.

After normalization, the iterates approach
the one-dimensional subspace
\[
\Span\{v_*\}.
\]

The convergence speed is governed by the
ratio between the second-largest and largest
eigenvalue magnitudes.

When this ratio is close to one, convergence
may be slow.

When it is substantially smaller than one,
convergence is generally faster.

\begin{remark}
The ordinary power method targets the eigenvalue
of largest \emph{magnitude}, which is not
necessarily the largest algebraic eigenvalue.

For example, if the eigenvalues are
\[
-10,\quad 2,\quad 5,
\]
the dominant eigenvalue in magnitude is $-10$,
not $5$.

Moreover, if the dominant eigenvalue is negative,
successive normalized vectors may alternate in sign.
The corresponding one-dimensional eigenspace can
still converge.
\end{remark}

Once an approximate eigenvector $x_k$ is obtained,
its eigenvalue may be estimated using the
Rayleigh quotient:
\[
\boxed{
\lambda_*\approx R_A(x_k).
}
\]

We shall later encounter another important
connection between eigenvalue computation and
optimization through the Courant--Fischer theorem.

% ============================================================
% 2.5 SVD AND MATRIX NORMS
% ============================================================

\subsection{Singular Value Decomposition and Matrix Norms}
\label{subsec:svd_background}

Unlike eigendecomposition, singular value
decomposition is available for every real
matrix, including rectangular matrices.

\begin{theorem}[Singular Value Decomposition]
Let
\[
B\in\R^{m\times n}.
\]

Then there exist orthogonal matrices
\[
U\in\R^{m\times m},
\qquad
V\in\R^{n\times n},
\]
and a rectangular diagonal matrix
\[
\Sigma\in\R^{m\times n}
\]
such that
\[
\boxed{
B=U\Sigma V^T.
}
\]
\end{theorem}

The nonnegative diagonal entries of $\Sigma$
are called the singular values of $B$.

We arrange them in decreasing order:
\[
\boxed{
\sigma_1\geq\sigma_2\geq\cdots
\geq\sigma_{\min(m,n)}\geq0.
}
\]

The singular values describe how strongly
the matrix stretches vectors along certain
orthogonal directions.

In particular, the spectral norm of $B$ is
defined by
\[
\boxed{
\norm{B}_2
=
\max_{\norm{x}_2=1}\norm{Bx}_2.
}
\]

The singular value decomposition gives
\[
\boxed{
\norm{B}_2=\sigma_1(B).
}
\]

Thus, the spectral norm measures the largest
possible stretching produced by the matrix.

Another commonly used matrix norm is the
Frobenius norm:
\[
\boxed{
\norm{B}_F
=
\sqrt{
\sum_{i=1}^{m}
\sum_{j=1}^{n}b_{ij}^2
}.
}
\]

Equivalently,
\[
\norm{B}_F^2
=
\sum_i\sigma_i(B)^2.
\]

For a real symmetric matrix $A$, the singular
values are the absolute values of its eigenvalues.

Therefore,
\[
\boxed{
\norm{A}_2
=
\max_i|\lambda_i(A)|.
}
\]

This identity will be especially important
when we study matrix perturbations.

\paragraph{Connection to symmetric matrices.}

For any real matrix $B$, both
\[
B^TB
\qquad\text{and}\qquad
BB^T
\]
are symmetric and positive semidefinite.

Their nonzero eigenvalues are the squares
of the nonzero singular values of $B$.

Thus, singular value decomposition connects
general matrices with the spectral theory
of symmetric matrices.

Later, we shall also see that singular values
provide a natural way to measure the
angular separation between subspaces.

% ============================================================
% 2.6 METRIC SPACES
% ============================================================

\subsection{Metric Spaces}
\label{subsec:metric_background}

In ordinary Euclidean geometry, the distance
between two points is familiar.

However, in mathematics, we may wish to
measure distances between more general objects,
including functions, matrices, and subspaces.

The abstract notion of a metric provides
a rigorous framework for doing so.

\begin{definition}[Metric Space]
Let $X$ be a nonempty set.

A function
\[
d:X\times X\longrightarrow[0,\infty)
\]
is called a metric if, for every
$x,y,z\in X$, it satisfies the following
properties:

\begin{enumerate}[label=(\roman*)]

\item \textbf{Nonnegativity:}
\[
d(x,y)\geq0.
\]

\item \textbf{Identity of indiscernibles:}
\[
d(x,y)=0
\quad\Longleftrightarrow\quad
x=y.
\]

\item \textbf{Symmetry:}
\[
d(x,y)=d(y,x).
\]

\item \textbf{Triangle inequality:}
\[
d(x,z)\leq d(x,y)+d(y,z).
\]

\end{enumerate}

The pair $(X,d)$ is called a metric space.
\end{definition}

The most familiar example is Euclidean space,
where
\[
X=\R^n
\]
and
\[
\boxed{
d(x,y)=\norm{x-y}_2.
}
\]

The significance of the metric-space concept
is that the objects being compared need not
be ordinary vectors.

For instance, we may consider a collection
whose elements are themselves linear subspaces.

If we can construct a distance function
satisfying the metric axioms, then we have
a mathematically consistent way to compare
those subspaces.

This idea will be central to our development
of subspace distance.

% ============================================================
% 2.7 ORTHOGONAL PROJECTIONS
% ============================================================

\subsection{Orthogonal Projections}
\label{subsec:projection_background}

Let $\mathcal S$ be a $k$-dimensional
subspace of $\R^n$, and suppose
\[
Q\in\R^{n\times k}
\]
has orthonormal columns spanning $\mathcal S$.

Thus,
\[
Q^TQ=I_k,
\qquad
\mathcal S=\range(Q).
\]

The orthogonal projection of a vector
$x\in\R^n$ onto $\mathcal S$ is given by
\[
\boxed{
P_{\mathcal S}x=QQ^Tx.
}
\]

The matrix
\[
\boxed{
P_{\mathcal S}=QQ^T
}
\]
is called the orthogonal projection matrix
onto $\mathcal S$.

It satisfies two important properties:
\[
\boxed{
P_{\mathcal S}^T=P_{\mathcal S}
}
\]
and
\[
\boxed{
P_{\mathcal S}^2=P_{\mathcal S}.
}
\]

The first property follows from
\[
(QQ^T)^T=QQ^T.
\]

The second follows from
\[
(QQ^T)(QQ^T)
=
Q(Q^TQ)Q^T
=
QQ^T.
\]

Thus, an orthogonal projector is symmetric
and idempotent.

Every vector $x\in\R^n$ can be decomposed as
\[
\boxed{
x=P_{\mathcal S}x+(I-P_{\mathcal S})x.
}
\]

The first component belongs to $\mathcal S$,
while the second belongs to its orthogonal
complement $\mathcal S^\perp$.

Moreover, these two components are orthogonal.

Consequently,
\[
\boxed{
\norm{x}_2^2
=
\norm{P_{\mathcal S}x}_2^2
+
\norm{(I-P_{\mathcal S})x}_2^2.
}
\]

This is simply the Pythagorean theorem
expressed in terms of orthogonal projections.

% ============================================================
% 2.8 VECTOR DISTANCE AND SUBSPACE DISTANCE
% ============================================================

\subsection{Distances Between Vectors and Subspaces}
\label{subsec:distance_background}

Before proceeding, it is important to distinguish
three related but different geometric questions:

\begin{enumerate}[label=(\roman*)]

\item What is the distance between two vectors?

\item What is the distance from a vector to a subspace?

\item What is the distance between two subspaces?

\end{enumerate}

These questions involve different mathematical
objects and should not be confused.

\paragraph{Distance between two vectors.}

For $x,y\in\R^n$, their Euclidean distance is
\[
\boxed{
d(x,y)=\norm{x-y}_2.
}
\]

This measures the length of the vector
joining $x$ and $y$.

\paragraph{Distance from a vector to a subspace.}

Let $\mathcal S\subseteq\R^n$ be a subspace.

The distance from a vector $x$ to $\mathcal S$
is defined by
\[
\boxed{
d(x,\mathcal S)
=
\min_{s\in\mathcal S}\norm{x-s}_2.
}
\]

Because the orthogonal projection
$P_{\mathcal S}x$ is the closest point
in $\mathcal S$ to $x$, we obtain
\[
\boxed{
d(x,\mathcal S)
=
\norm{x-P_{\mathcal S}x}_2.
}
\]

Equivalently,
\[
\boxed{
d(x,\mathcal S)
=
\norm{(I-P_{\mathcal S})x}_2.
}
\]

Thus, this distance measures the magnitude
of the component of $x$ perpendicular
to the subspace.

\paragraph{Distance between two subspaces.}

Now suppose $\mathcal U$ and $\mathcal W$
are two subspaces of the same ambient
Euclidean space.

Unlike vectors, subspaces are sets
containing infinitely many vectors.

Therefore, the expression
\[
\norm{\mathcal U-\mathcal W}_2
\]
does not define an ordinary Euclidean
distance between them.

Furthermore, the choice of a particular
basis should not affect how far apart
two subspaces are considered to be.

For example, consider
\[
u=
\begin{bmatrix}
1\\0
\end{bmatrix},
\qquad
w=
\begin{bmatrix}
-1\\0
\end{bmatrix}.
\]

Their Euclidean distance is
\[
\norm{u-w}_2=2.
\]

Nevertheless,
\[
\boxed{
\Span\{u\}=\Span\{w\}.
}
\]

Both vectors represent exactly the same
one-dimensional subspace: the horizontal
axis in $\R^2$.

A reasonable notion of subspace distance
must therefore assign distance zero
to these two identical subspaces.

This illustrates why ordinary vector
distance is not appropriate for
comparing eigenspaces.

In the next section, we shall construct
a natural distance between subspaces
using orthogonal projection matrices.

We shall verify that this distance
satisfies the axioms of a metric and
develop its geometric interpretation
through principal angles and the
singular value decomposition.

% ============================================================
% END OF SECTION 2
% ============================================================



% ============================================================
% SECTION 2: MATHEMATICAL BACKGROUND
% ============================================================

%% ============================================================
% SECTION 3: SUBSPACE GEOMETRY AND DISTANCE
% ============================================================
\begin{comment}
\section{Subspace Geometry and Distance}
\label{sec:subspace_geometry}

In the preceding section, we distinguished the
distance between two vectors from the distance
between two subspaces. We now develop a rigorous
mathematical framework for comparing subspaces
of a Euclidean space.

Our principal objective is to construct a notion
of distance that depends only on the subspaces
themselves, rather than on the particular vectors
chosen to represent them.

Orthogonal projection matrices provide a natural
solution to this problem. As we shall establish,
they allow us to define a distance that satisfies
the axioms of a metric and possesses a useful
geometric interpretation in terms of angles.

% ============================================================
% 3.1 SUBSPACES AND THEIR BASES
% ============================================================

\subsection{Subspaces and Their Bases}
\label{subsec:subspace_bases}

A fundamental distinction in linear algebra is
the distinction between a vector space and a
particular basis used to describe it.

Consider a $k$-dimensional subspace
\[
\mathcal W\subseteq\R^n.
\]

Suppose
\[
W=
\begin{bmatrix}
w_1 & w_2 & \cdots & w_k
\end{bmatrix}
\in\R^{n\times k}
\]
has orthonormal columns spanning $\mathcal W$.

Thus,
\[
W^TW=I_k,
\qquad
\mathcal W=\range(W).
\]

The columns of $W$ provide a coordinate
description of $\mathcal W$, but this
description is not unique.

Indeed, another matrix
\[
U=
\begin{bmatrix}
u_1 & u_2 & \cdots & u_k
\end{bmatrix}
\in\R^{n\times k}
\]
may also have orthonormal columns and span
exactly the same subspace.

In that case,
\[
\range(W)=\range(U)=\mathcal W.
\]

Although the matrices $W$ and $U$ may
be different, they represent the same
geometric object.

This observation is particularly important
when comparing invariant subspaces.

A numerical eigensolver may return a
different orthonormal basis for the same
eigenspace without changing the eigenspace
itself.

Therefore, measuring the difference
between two basis matrices does not
necessarily measure the difference
between their corresponding subspaces.

\begin{example}[Different Bases for the Same Subspace]

Consider the vectors
\[
w=
\begin{bmatrix}
1\\0
\end{bmatrix},
\qquad
u=
\begin{bmatrix}
-1\\0
\end{bmatrix}.
\]

Clearly,
\[
w\neq u,
\]
and
\[
\norm{w-u}_2=2.
\]

Nevertheless,
\[
\Span\{w\}
=
\Span\{u\}.
\]

Thus, the two vectors represent the
same one-dimensional subspace.

Any meaningful distance between subspaces
must assign distance zero to these
two representations.

\end{example}

The preceding example suggests that
we should search for an object that
uniquely represents a subspace,
independently of the basis used.

Before constructing such an object,
we establish an important relationship
between two orthonormal bases of the
same subspace.

\begin{proposition}[Change of Orthonormal Basis]
\label{prop:orthonormal_change_basis}

Let $W,U\in\R^{n\times k}$ have
orthonormal columns.

Then
\[
\range(W)=\range(U)
\]
if and only if there exists an
orthogonal matrix
\[
Q\in\R^{k\times k}
\]
such that
\[
\boxed{
U=WQ.
}
\]

\end{proposition}

\begin{proof}

First suppose
\[
\range(W)=\range(U).
\]

Since every column of $U$ belongs
to the column space of $W$,
there exists a matrix
$Q\in\R^{k\times k}$ such that
\[
U=WQ.
\]

Multiplying by $W^T$ gives
\[
W^TU=W^TWQ.
\]

Since $W^TW=I_k$,
\[
\boxed{
Q=W^TU.
}
\]

We now show that $Q$ is orthogonal.

Using $U=WQ$ and the orthonormality
of the columns of $U$, we obtain
\[
\begin{aligned}
I_k
&=U^TU\\
&=(WQ)^T(WQ)\\
&=Q^TW^TWQ\\
&=Q^TQ.
\end{aligned}
\]

Therefore,
\[
Q^TQ=I_k.
\]

Since $Q$ is square, it follows
that $Q$ is orthogonal.

Conversely, suppose
\[
U=WQ
\]
for some orthogonal matrix $Q$.

Because $Q$ is invertible,
\[
W=UQ^T.
\]

Consequently, every column of $U$
is a linear combination of the
columns of $W$, and every column
of $W$ is a linear combination
of the columns of $U$.

Thus,
\[
\range(U)\subseteq\range(W)
\]
and
\[
\range(W)\subseteq\range(U).
\]

We conclude that
\[
\boxed{
\range(W)=\range(U).
}
\]

\end{proof}

This result has an important
geometric interpretation.

Two orthonormal bases of the same
subspace differ only by an
orthogonal change of coordinates.

The basis vectors may change,
but the subspace remains unchanged.

Our next objective is to construct
a matrix that remains unchanged
under precisely such transformations.

% ============================================================
% 3.2 ORTHOGONAL PROJECTORS AND BASIS INDEPENDENCE
% ============================================================

\subsection{Orthogonal Projectors as Unique
Representations of Subspaces}
\label{subsec:projector_representation}

Let $\mathcal W$ be a $k$-dimensional
subspace of $\R^n$, and let
\[
W\in\R^{n\times k}
\]
have orthonormal columns spanning
$\mathcal W$.

Recall that the orthogonal projection
matrix onto $\mathcal W$ is
\[
\boxed{
P_{\mathcal W}=WW^T.
}
\]

At first, this expression might
appear to depend on the particular
matrix $W$.

After all, a subspace admits
infinitely many possible
orthonormal bases.

However, the following proposition
shows that this apparent dependence
does not exist.

\begin{proposition}[Basis Independence of
Orthogonal Projection]
\label{prop:projection_basis_independence}

Let $W,U\in\R^{n\times k}$
have orthonormal columns spanning
the same subspace.

Then
\[
\boxed{
WW^T=UU^T.
}
\]

\end{proposition}

\begin{proof}

Since $W$ and $U$ span the
same subspace, the preceding
proposition guarantees the
existence of an orthogonal
matrix $Q$ such that
\[
U=WQ.
\]

Therefore,
\[
\begin{aligned}
UU^T
&=(WQ)(WQ)^T\\
&=WQQ^TW^T.
\end{aligned}
\]

Because $Q$ is orthogonal,
\[
QQ^T=I_k.
\]

Hence,
\[
\begin{aligned}
UU^T
&=WI_kW^T\\
&=WW^T.
\end{aligned}
\]

Thus,
\[
\boxed{
UU^T=WW^T.
}
\]

\end{proof}

This proposition establishes that
the orthogonal projection matrix
depends only on the subspace,
not on the basis used to describe it.

For example, the vectors
\[
w=
\begin{bmatrix}
1\\0
\end{bmatrix},
\qquad
u=
\begin{bmatrix}
-1\\0
\end{bmatrix}
\]
produce the same projector:
\[
ww^T
=
uu^T
=
\begin{bmatrix}
1&0\\
0&0
\end{bmatrix}.
\]

The projection matrix therefore
recognizes that these two vectors
span the same line.

But we can establish something
stronger: two subspaces are equal
if and only if their orthogonal
projection matrices are equal.

\begin{proposition}[Uniqueness of the
Orthogonal Projector]
\label{prop:unique_projector}

Let $\mathcal U$ and $\mathcal W$
be linear subspaces of $\R^n$.

Then
\[
\boxed{
\mathcal U=\mathcal W
\quad\Longleftrightarrow\quad
P_{\mathcal U}=P_{\mathcal W}.
}
\]

\end{proposition}

\begin{proof}

Suppose first that
\[
\mathcal U=\mathcal W.
\]

Since the orthogonal projection
matrix is independent of the
chosen orthonormal basis,
the projectors onto these
subspaces must be equal:
\[
P_{\mathcal U}=P_{\mathcal W}.
\]

Conversely, suppose
\[
P_{\mathcal U}=P_{\mathcal W}.
\]

Recall that the range of an
orthogonal projection matrix
is precisely the subspace
onto which it projects.

To see this, let $P_{\mathcal U}$
denote the projector onto
$\mathcal U$.

By definition, for every
$x\in\R^n$,
\[
P_{\mathcal U}x\in\mathcal U.
\]

Hence,
\[
\range(P_{\mathcal U})
\subseteq\mathcal U.
\]

On the other hand, if
$u\in\mathcal U$, then
orthogonal projection leaves
$u$ unchanged:
\[
P_{\mathcal U}u=u.
\]

Therefore,
\[
\mathcal U
\subseteq\range(P_{\mathcal U}).
\]

Combining these inclusions gives
\[
\boxed{
\range(P_{\mathcal U})=\mathcal U.
}
\]

Similarly,
\[
\range(P_{\mathcal W})=\mathcal W.
\]

Since the projectors are equal,
their ranges must be equal.

Consequently,
\[
\boxed{
\mathcal U=\mathcal W.
}
\]

\end{proof}

We have now established the
essential property required
for comparing subspaces:

\begin{quote}
\emph{
Every linear subspace of a
Euclidean space is uniquely
determined by its orthogonal
projection matrix.
}
\end{quote}

This allows us to replace
the problem of comparing
subspaces with the more
familiar problem of
comparing matrices.

% ============================================================
% 3.3 A METRIC ON THE SET OF SUBSPACES
% ============================================================

\subsection{A Metric Defined by Orthogonal Projectors}
\label{subsec:projector_metric}

We now introduce the central
definition of subspace distance
used throughout this note.

\begin{definition}[Projection Distance]
\label{def:projection_distance}

Let $\mathcal U$ and $\mathcal W$
be linear subspaces of $\R^n$,
and let $P_{\mathcal U}$ and
$P_{\mathcal W}$ denote their
orthogonal projection matrices.

We define the projection distance
between $\mathcal U$ and
$\mathcal W$ by
\[
\boxed{
d(\mathcal U,\mathcal W)
=
\norm{
P_{\mathcal U}-P_{\mathcal W}
}_2.
}
\]

Here, $\norm{\cdot}_2$ denotes
the spectral matrix norm.

\end{definition}

To understand this definition,
recall that
\[
\norm{B}_2
=
\max_{\norm{x}_2=1}\norm{Bx}_2.
\]

Therefore,
\[
\boxed{
d(\mathcal U,\mathcal W)
=
\max_{\norm{x}_2=1}
\norm{
P_{\mathcal U}x-P_{\mathcal W}x
}_2.
}
\]

This expression provides an
intuitive interpretation.

For any unit vector $x$,
we project $x$ onto each of
the two subspaces and measure
the difference between the
resulting projected vectors.

The projection distance is
the largest such difference
over all unit vectors.

Thus, the distance measures
how differently the two
subspaces act on vectors
through their orthogonal
projections.

If the subspaces are identical,
their projections agree
for every vector, and the
distance is zero.

If the subspaces are different,
there must be some vector
whose projections are different,
and consequently the distance
is positive.

We now establish that this
construction is a genuine metric.

\begin{theorem}[Projection Distance Is a Metric]
\label{thm:projection_metric}

Let $\mathscr S_n$ denote
the collection of all linear
subspaces of $\R^n$.

Then
\[
d(\mathcal U,\mathcal W)
=
\norm{
P_{\mathcal U}-P_{\mathcal W}
}_2
\]
defines a metric on $\mathscr S_n$.

\end{theorem}

\begin{proof}

We verify the four metric axioms.

\medskip

\noindent
\textbf{1. Nonnegativity.}

Since the spectral norm is
nonnegative,
\[
\norm{
P_{\mathcal U}-P_{\mathcal W}
}_2
\geq0.
\]

Therefore,
\[
\boxed{
d(\mathcal U,\mathcal W)\geq0.
}
\]

\medskip

\noindent
\textbf{2. Identity of indiscernibles.}

Suppose
\[
d(\mathcal U,\mathcal W)=0.
\]

By definition,
\[
\norm{
P_{\mathcal U}-P_{\mathcal W}
}_2=0.
\]

A matrix has zero spectral
norm if and only if it is
the zero matrix.

Consequently,
\[
P_{\mathcal U}-P_{\mathcal W}=0,
\]
or equivalently,
\[
P_{\mathcal U}=P_{\mathcal W}.
\]

By Proposition
\ref{prop:unique_projector},
this implies
\[
\mathcal U=\mathcal W.
\]

Conversely, if
\[
\mathcal U=\mathcal W,
\]
then their orthogonal
projectors are equal.

Therefore,
\[
d(\mathcal U,\mathcal W)
=
\norm{0}_2=0.
\]

Hence,
\[
\boxed{
d(\mathcal U,\mathcal W)=0
\quad\Longleftrightarrow\quad
\mathcal U=\mathcal W.
}
\]

\medskip

\noindent
\textbf{3. Symmetry.}

Observe that
\[
P_{\mathcal U}-P_{\mathcal W}
=
-\left(
P_{\mathcal W}-P_{\mathcal U}
\right).
\]

Using the absolute homogeneity
of the spectral norm,
\[
\begin{aligned}
d(\mathcal U,\mathcal W)
&=
\norm{
P_{\mathcal U}-P_{\mathcal W}
}_2\\
&=
\norm{
-\left(
P_{\mathcal W}-P_{\mathcal U}
\right)
}_2\\
&=
\norm{
P_{\mathcal W}-P_{\mathcal U}
}_2\\
&=
d(\mathcal W,\mathcal U).
\end{aligned}
\]

Thus,
\[
\boxed{
d(\mathcal U,\mathcal W)
=
d(\mathcal W,\mathcal U).
}
\]

\medskip

\noindent
\textbf{4. Triangle inequality.}

Let $\mathcal U,\mathcal V$,
and $\mathcal W$ be arbitrary
linear subspaces of $\R^n$.

We may write
\[
P_{\mathcal U}-P_{\mathcal W}
=
\left(
P_{\mathcal U}-P_{\mathcal V}
\right)
+
\left(
P_{\mathcal V}-P_{\mathcal W}
\right).
\]

Therefore,
\[
\begin{aligned}
d(\mathcal U,\mathcal W)
&=
\norm{
P_{\mathcal U}-P_{\mathcal W}
}_2\\[4pt]
&=
\norm{
(P_{\mathcal U}-P_{\mathcal V})
+
(P_{\mathcal V}-P_{\mathcal W})
}_2.
\end{aligned}
\]

Applying the triangle inequality
for the spectral matrix norm,
\[
\begin{aligned}
d(\mathcal U,\mathcal W)
&\leq
\norm{
P_{\mathcal U}-P_{\mathcal V}
}_2\\
&\quad+
\norm{
P_{\mathcal V}-P_{\mathcal W}
}_2.
\end{aligned}
\]

Hence,
\[
\boxed{
d(\mathcal U,\mathcal W)
\leq
d(\mathcal U,\mathcal V)
+
d(\mathcal V,\mathcal W).
}
\]

We have verified all four
metric axioms.

Therefore,
\[
\boxed{
(\mathscr S_n,d)
\text{ is a metric space.}
}
\]

\end{proof}

An important feature of this
construction is that the
subspaces need not have
the same dimension.

The distance is well-defined
for any two linear subspaces
of the same ambient Euclidean
space.

However, our primary interest
will be in comparing subspaces
of equal dimension, particularly
those arising from eigenvectors
of symmetric matrices.

The collection of all
$k$-dimensional subspaces
of $\R^n$ is called the
Grassmannian and is denoted by
\[
\boxed{
\operatorname{Gr}(k,n).
}
\]

The projection distance
also defines a metric on
this collection.

In subsequent subsections,
we shall discover that this
metric has a particularly
elegant geometric interpretation
on $\operatorname{Gr}(k,n)$.

Specifically, its value is
determined by the largest
principal angle between
the two subspaces.

\end{comment}
% ============================================================
% SECTION 3: SUBSPACE GEOMETRY AND DISTANCE
% ============================================================
%% ============================================================
% SECTION 3: SUBSPACE GEOMETRY AND DISTANCE
% ============================================================
% ============================================================
% SECTION 3: SUBSPACE GEOMETRY AND DISTANCE
% ============================================================

\section{Subspace Geometry and Distance}
\label{sec:subspace_geometry}

In the preceding section, we distinguished the
distance between two vectors from the distance
between two subspaces. We now develop a rigorous
mathematical framework for comparing subspaces
of a Euclidean space.

Our principal objective is to construct a notion
of distance that depends only on the subspaces
themselves, rather than on the particular vectors 
chosen to represent them.

Orthogonal projection matrices provide a natural
solution to this problem. As we shall establish,
they allow us to define a distance that satisfies
the axioms of a metric and possesses a useful
geometric interpretation in terms of angles.

% ============================================================
% 3.1 SUBSPACES AND THEIR BASES
% ============================================================

\subsection{Subspaces and Their Bases}
\label{subsec:subspace_bases}

A fundamental distinction in linear algebra is
the distinction between a vector space and a
particular basis used to describe it.

Consider a $k$-dimensional subspace
\[
\mathcal W\subseteq\R^n.
\]

Suppose
\[
W=
\begin{bmatrix}
w_1 & w_2 & \cdots & w_k
\end{bmatrix}
\in\R^{n\times k}
\]
has orthonormal columns spanning $\mathcal W$.

Thus,
\[
W^TW=I_k,
\qquad
\mathcal W=\range(W).
\]

The columns of $W$ provide a coordinate
description of $\mathcal W$, but this
description is not unique.

Indeed, another matrix
\[
U=
\begin{bmatrix}
u_1 & u_2 & \cdots & u_k
\end{bmatrix}
\in\R^{n\times k}
\]
may also have orthonormal columns and span
exactly the same subspace.

In that case,
\[
\range(W)=\range(U)=\mathcal W.
\]

Although the matrices $W$ and $U$ may
be different, they represent the same
geometric object.

This observation is particularly important
when comparing invariant subspaces.

A numerical eigensolver may return a
different orthonormal basis for the same
eigenspace without changing the eigenspace
itself.

Therefore, measuring the difference
between two basis matrices does not
necessarily measure the difference
between their corresponding subspaces.

\begin{example}[Different Bases for the Same Subspace]

Consider the vectors
\[
w=
\begin{bmatrix}
1\\0
\end{bmatrix},
\qquad
u=
\begin{bmatrix}
-1\\0
\end{bmatrix}.
\]

Clearly,
\[
w\neq u,
\]
and
\[
\norm{w-u}_2=2.
\]

Nevertheless,
\[
\Span\{w\}
=
\Span\{u\}.
\]

Thus, the two vectors represent the
same one-dimensional subspace.

Any meaningful distance between subspaces
must assign distance zero to these
two representations.

\end{example}

The preceding example suggests that
we should search for an object that
uniquely represents a subspace,
independently of the basis used.

Before constructing such an object,
we establish an important relationship
between two orthonormal bases of the
same subspace.

\begin{proposition}[Change of Orthonormal Basis]
\label{prop:orthonormal_change_basis}

Let $W,U\in\R^{n\times k}$ have
orthonormal columns.

Then
\[
\range(W)=\range(U)
\]
if and only if there exists an
orthogonal matrix
\[
Q\in\R^{k\times k}
\]
such that
\[
\boxed{
U=WQ.
}
\]

\end{proposition}

\begin{proof}

First suppose
\[
\range(W)=\range(U).
\]

Since every column of $U$ belongs
to the column space of $W$,
there exists a matrix
$Q\in\R^{k\times k}$ such that
\[
U=WQ.
\]

Multiplying by $W^T$ gives
\[
W^TU=W^TWQ.
\]

Since $W^TW=I_k$,
\[
\boxed{
Q=W^TU.
}
\]

We now show that $Q$ is orthogonal.

Using $U=WQ$ and the orthonormality
of the columns of $U$, we obtain
\[
\begin{aligned}
I_k
&=U^TU\\
&=(WQ)^T(WQ)\\
&=Q^TW^TWQ\\
&=Q^TQ.
\end{aligned}
\]

Therefore,
\[
Q^TQ=I_k.
\]

Since $Q$ is square, it follows
that $Q$ is orthogonal.

Conversely, suppose
\[
U=WQ
\]
for some orthogonal matrix $Q$.

Because $Q$ is invertible,
\[
W=UQ^T.
\]

Consequently, every column of $U$
is a linear combination of the
columns of $W$, and every column
of $W$ is a linear combination
of the columns of $U$.

Thus,
\[
\range(U)\subseteq\range(W)
\]
and
\[
\range(W)\subseteq\range(U).
\]

We conclude that
\[
\boxed{
\range(W)=\range(U).
}
\]

\end{proof}

This result has an important
geometric interpretation.

Two orthonormal bases of the same
subspace differ only by an
orthogonal change of coordinates.

The basis vectors may change,
but the subspace remains unchanged.

Our next objective is to construct
a matrix that remains unchanged
under precisely such transformations.

% ============================================================
% 3.2 ORTHOGONAL PROJECTORS AND BASIS INDEPENDENCE
% ============================================================

\subsection{Orthogonal Projectors as Unique
Representations of Subspaces}
\label{subsec:projector_representation}

Let $\mathcal W$ be a $k$-dimensional
subspace of $\R^n$, and let
\[
W\in\R^{n\times k}
\]
have orthonormal columns spanning
$\mathcal W$.

Recall that the orthogonal projection
matrix onto $\mathcal W$ is
\[
\boxed{
P_{\mathcal W}=WW^T.
}
\]

At first, this expression might
appear to depend on the particular
matrix $W$.

After all, a subspace admits
infinitely many possible
orthonormal bases.

However, the following proposition
shows that this apparent dependence
does not exist.

\begin{proposition}[Basis Independence of
Orthogonal Projection]
\label{prop:projection_basis_independence}

Let $W,U\in\R^{n\times k}$
have orthonormal columns spanning
the same subspace.

Then
\[
\boxed{
WW^T=UU^T.
}
\]

\end{proposition}

\begin{proof}

Since $W$ and $U$ span the
same subspace, the preceding
proposition guarantees the
existence of an orthogonal
matrix $Q$ such that
\[
U=WQ.
\]

Therefore,
\[
\begin{aligned}
UU^T
&=(WQ)(WQ)^T\\
&=WQQ^TW^T.
\end{aligned}
\]

Because $Q$ is orthogonal,
\[
QQ^T=I_k.
\]

Hence,
\[
\begin{aligned}
UU^T
&=WI_kW^T\\
&=WW^T.
\end{aligned}
\]

Thus,
\[
\boxed{
UU^T=WW^T.
}
\]

\end{proof}

This proposition establishes that
the orthogonal projection matrix
depends only on the subspace,
not on the basis used to describe it.

For example, the vectors
\[
w=
\begin{bmatrix}
1\\0
\end{bmatrix},
\qquad
u=
\begin{bmatrix}
-1\\0
\end{bmatrix}
\]
produce the same projector:
\[
ww^T
=
uu^T
=
\begin{bmatrix}
1&0\\
0&0
\end{bmatrix}.
\]

The projection matrix therefore
recognizes that these two vectors
span the same line.

But we can establish something
stronger: two subspaces are equal
if and only if their orthogonal
projection matrices are equal.

\begin{proposition}[Uniqueness of the
Orthogonal Projector]
\label{prop:unique_projector}

Let $\mathcal U$ and $\mathcal W$
be linear subspaces of $\R^n$.

Then
\[
\boxed{
\mathcal U=\mathcal W
\quad\Longleftrightarrow\quad
P_{\mathcal U}=P_{\mathcal W}.
}
\]

\end{proposition}

\begin{proof}

Suppose first that
\[
\mathcal U=\mathcal W.
\]

Since the orthogonal projection
matrix is independent of the
chosen orthonormal basis,
the projectors onto these
subspaces must be equal:
\[
P_{\mathcal U}=P_{\mathcal W}.
\]

Conversely, suppose
\[
P_{\mathcal U}=P_{\mathcal W}.
\]

Recall that the range of an
orthogonal projection matrix
is precisely the subspace
onto which it projects.

To see this, let $P_{\mathcal U}$
denote the projector onto
$\mathcal U$.

By definition, for every
$x\in\R^n$,
\[
P_{\mathcal U}x\in\mathcal U.
\]

Hence,
\[
\range(P_{\mathcal U})
\subseteq\mathcal U.
\]

On the other hand, if
$u\in\mathcal U$, then
orthogonal projection leaves
$u$ unchanged:
\[
P_{\mathcal U}u=u.
\]

Therefore,
\[
\mathcal U
\subseteq\range(P_{\mathcal U}).
\]

Combining these inclusions gives
\[
\boxed{
\range(P_{\mathcal U})=\mathcal U.
}
\]

Similarly,
\[
\range(P_{\mathcal W})=\mathcal W.
\]

Since the projectors are equal,
their ranges must be equal.

Consequently,
\[
\boxed{
\mathcal U=\mathcal W.
}
\]

\end{proof}

We have now established the
essential property required
for comparing subspaces:

\begin{quote}
\emph{
Every linear subspace of a
Euclidean space is uniquely
determined by its orthogonal
projection matrix.
}
\end{quote}

This allows us to replace
the problem of comparing
subspaces with the more
familiar problem of
comparing matrices.

% ============================================================
% 3.3 A METRIC ON THE SET OF SUBSPACES
% ============================================================

\subsection{A Metric Defined by Orthogonal Projectors}
\label{subsec:projector_metric}

We now introduce the central
definition of subspace distance
used throughout this note.

\begin{definition}[Projection Distance]
\label{def:projection_distance}

Let $\mathcal U$ and $\mathcal W$
be linear subspaces of $\R^n$,
and let $P_{\mathcal U}$ and
$P_{\mathcal W}$ denote their
orthogonal projection matrices.

We define the projection distance
between $\mathcal U$ and
$\mathcal W$ by
\[
\boxed{
d(\mathcal U,\mathcal W)
=
\norm{
P_{\mathcal U}-P_{\mathcal W}
}_2.
}
\]

Here, $\norm{\cdot}_2$ denotes
the spectral matrix norm.

\end{definition}

To understand this definition,
recall that
\[
\norm{B}_2
=
\max_{\norm{x}_2=1}\norm{Bx}_2.
\]

Therefore,
\[
\boxed{
d(\mathcal U,\mathcal W)
=
\max_{\norm{x}_2=1}
\norm{
P_{\mathcal U}x-P_{\mathcal W}x
}_2.
}
\]

This expression provides an
intuitive interpretation.

For any unit vector $x$,
we project $x$ onto each of
the two subspaces and measure
the difference between the
resulting projected vectors.

The projection distance is
the largest such difference
over all unit vectors.

Thus, the distance measures
how differently the two
subspaces act on vectors
through their orthogonal
projections.

If the subspaces are identical,
their projections agree
for every vector, and the
distance is zero.

If the subspaces are different,
there must be some vector
whose projections are different,
and consequently the distance
is positive.

We now establish that this
construction is a genuine metric.

\begin{theorem}[Projection Distance Is a Metric]
\label{thm:projection_metric}

Let $\mathscr S_n$ denote
the collection of all linear
subspaces of $\R^n$.

Then
\[
d(\mathcal U,\mathcal W)
=
\norm{
P_{\mathcal U}-P_{\mathcal W}
}_2
\]
defines a metric on $\mathscr S_n$.

\end{theorem}

\begin{proof}

We verify the four metric axioms.

\medskip

\noindent
\textbf{1. Nonnegativity.}

Since the spectral norm is
nonnegative,
\[
\norm{
P_{\mathcal U}-P_{\mathcal W}
}_2
\geq0.
\]

Therefore,
\[
\boxed{
d(\mathcal U,\mathcal W)\geq0.
}
\]

\medskip

\noindent
\textbf{2. Identity of indiscernibles.}

Suppose
\[
d(\mathcal U,\mathcal W)=0.
\]

By definition,
\[
\norm{
P_{\mathcal U}-P_{\mathcal W}
}_2=0.
\]

A matrix has zero spectral
norm if and only if it is
the zero matrix.

Consequently,
\[
P_{\mathcal U}-P_{\mathcal W}=0,
\]
or equivalently,
\[
P_{\mathcal U}=P_{\mathcal W}.
\]

By Proposition
\ref{prop:unique_projector},
this implies
\[
\mathcal U=\mathcal W.
\]

Conversely, if
\[
\mathcal U=\mathcal W,
\]
then their orthogonal
projectors are equal.

Therefore,
\[
d(\mathcal U,\mathcal W)
=
\norm{0}_2=0.
\]

Hence,
\[
\boxed{
d(\mathcal U,\mathcal W)=0
\quad\Longleftrightarrow\quad
\mathcal U=\mathcal W.
}
\]

\medskip

\noindent
\textbf{3. Symmetry.}

Observe that
\[
P_{\mathcal U}-P_{\mathcal W}
=
-\left(
P_{\mathcal W}-P_{\mathcal U}
\right).
\]

Using the absolute homogeneity
of the spectral norm,
\[
\begin{aligned}
d(\mathcal U,\mathcal W)
&=
\norm{
P_{\mathcal U}-P_{\mathcal W}
}_2\\
&=
\norm{
-\left(
P_{\mathcal W}-P_{\mathcal U}
\right)
}_2\\
&=
\norm{
P_{\mathcal W}-P_{\mathcal U}
}_2\\
&=
d(\mathcal W,\mathcal U).
\end{aligned}
\]

Thus,
\[
\boxed{
d(\mathcal U,\mathcal W)
=
d(\mathcal W,\mathcal U).
}
\]

\medskip

\noindent
\textbf{4. Triangle inequality.}

Let $\mathcal U,\mathcal V$,
and $\mathcal W$ be arbitrary
linear subspaces of $\R^n$.

We may write
\[
P_{\mathcal U}-P_{\mathcal W}
=
\left(
P_{\mathcal U}-P_{\mathcal V}
\right)
+
\left(
P_{\mathcal V}-P_{\mathcal W}
\right).
\]

Therefore,
\[
\begin{aligned}
d(\mathcal U,\mathcal W)
&=
\norm{
P_{\mathcal U}-P_{\mathcal W}
}_2\\[4pt]
&=
\norm{
(P_{\mathcal U}-P_{\mathcal V})
+
(P_{\mathcal V}-P_{\mathcal W})
}_2.
\end{aligned}
\]

Applying the triangle inequality
for the spectral matrix norm,
\[
\begin{aligned}
d(\mathcal U,\mathcal W)
&\leq
\norm{
P_{\mathcal U}-P_{\mathcal V}
}_2\\
&\quad+
\norm{
P_{\mathcal V}-P_{\mathcal W}
}_2.
\end{aligned}
\]

Hence,
\[
\boxed{
d(\mathcal U,\mathcal W)
\leq
d(\mathcal U,\mathcal V)
+
d(\mathcal V,\mathcal W).
}
\]

We have verified all four
metric axioms.

Therefore,
\[
\boxed{
(\mathscr S_n,d)
\text{ is a metric space.}
}
\]

\end{proof}

An important feature of this
construction is that the
subspaces need not have
the same dimension.

The distance is well-defined
for any two linear subspaces
of the same ambient Euclidean
space.

However, our primary interest
will be in comparing subspaces
of equal dimension, particularly
those arising from eigenvectors
of symmetric matrices.

It is worth to mention that the collection of all
$k$-dimensional subspaces
of $\R^n$ is called the
\textbf{Grassmannian} and is denoted by
\[
\boxed{
\operatorname{Gr}(k,n).
}
\]

The projection distance
also defines a metric on
this collection.

In subsequent subsections,
we shall discover that this
metric has a particularly
elegant geometric interpretation
on $\operatorname{Gr}(k,n)$.

Specifically, its value is
determined by the largest
principal angle between
the two subspaces.


% ============================================================
% 3.4 THE ONE-DIMENSIONAL CASE
% ============================================================

\subsection{Distance Between One-Dimensional Subspaces}
\label{subsec:one_dimensional_distance}

We have constructed a metric using the difference of
orthogonal projection matrices. But what does this
metric measure geometrically? The simplest case already
reveals the answer: for two lines through the origin,
the projection distance is the sine of the acute angle
between them.

Let $u,w\in\R^n$ be unit vectors, and consider
\[
\mathcal U=\Span\{u\},
\qquad
\mathcal W=\Span\{w\}.
\]

The corresponding orthogonal projectors are rank 1 matrices given by
\[
P_{\mathcal U}=uu^T,
\qquad
P_{\mathcal W}=ww^T.
\]

Notice that each subspace is a \emph{line}, not a
particular choice of a vector along that line. In
particular, replacing $w$ by $-w$ does not change
$\mathcal W$ or its projector, since
\[
(-w)(-w)^T=ww^T.
\]

We therefore measure the angle between the two
\emph{lines} by the acute angle
\[
\boxed{
\theta=\cos^{-1}\abs{u^Tw},
\qquad 0\leq\theta\leq\frac{\pi}{2}.
}
\]

The absolute value is essential: opposite unit vectors
span the same line and should give angle zero.

\begin{theorem}[Projection Distance Between Two Lines]
\label{thm:one_dimensional_distance}

Let $u,w\in\R^n$ be unit vectors, and let
$\mathcal U=\Span\{u\}$ and
$\mathcal W=\Span\{w\}$.

If $\theta$ is their acute angle, then
\[
\boxed{
\begin{aligned}
d(\mathcal U,\mathcal W)
&=\norm{uu^T-ww^T}_2\\
&=\sqrt{1-(u^Tw)^2}\\
&=\sin\theta.
\end{aligned}
}
\]

\end{theorem}

\begin{proof}

Since replacing $w$ with $-w$ does not alter its
projector, we may choose its sign so that
\[
c=u^Tw=\cos\theta\geq0.
\]

Set
\[
s=\sin\theta=\sqrt{1-c^2}.
\]

\noindent\textbf{Case 1: $s=0$.}

Then $c=1$. The Cauchy--Schwarz equality condition
for unit vectors implies $w=u$ (after our sign choice).

Thus $uu^T=ww^T$, and
\[
\norm{uu^T-ww^T}_2=0=\sin\theta.
\]

The desired identity follows.

\medskip

\noindent\textbf{Case 2: $s>0$.}

Define
\[
\boxed{
z=\frac{w-cu}{s}.
}
\]

This vector is orthogonal to $u$, because
\[
u^Tz
=\frac{u^Tw-c\,u^Tu}{s}
=\frac{c-c}{s}=0.
\]

It also has unit length. Indeed,
\[
\begin{aligned}
\norm{z}_2^2
&=\frac{(w-cu)^T(w-cu)}{s^2}\\
&=\frac{1-2c^2+c^2}{s^2}\\
&=\frac{1-c^2}{s^2}=1.
\end{aligned}
\]

Consequently, $u$ and $z$ are orthonormal, and
our definition of $z$ gives
\[
\boxed{w=cu+sz.}
\]

Thus the two lines lie in the two-dimensional plane
\[
\mathcal P=\Span\{u,z\}=\Span\{u,w\}.
\]

%\paragraph{Why the calculation reduces to this plane.}

Now consider the difference of projectors
\[
D=uu^T-ww^T.
\]

For any $x\in\R^n$,
\[
\boxed{
Dx=u(u^Tx)-w(w^Tx).
}
\]

Both terms on the right are scalar multiples of
$u$ or $w$. Therefore,
\[
Dx\in\mathcal P
\qquad\text{for every }x\in\R^n.
\]

Furthermore, if $x\in\mathcal P^\perp$, then
$u^Tx=w^Tx=0$, so
\[
Dx=0.
\]

Hence $D$ acts nontrivially only within $\mathcal P$
and vanishes on its orthogonal complement. We may
therefore compute its nonzero eigenvalues by studying
its action on the orthonormal basis $\{u,z\}$.

\paragraph{Matrix representation in the plane.}

Using $w^Tu=c$ and $w^Tz=s$, we obtain
\[
\begin{aligned}
Du
&=u(u^Tu)-w(w^Tu)\\
&=u-c(cu+sz)\\
&=(1-c^2)u-cs\,z\\
&=s^2u-cs\,z,
\end{aligned}
\]
and similarly,
\[
\begin{aligned}
Dz
&=u(u^Tz)-w(w^Tz)\\
&=-s(cu+sz)\\
&=-cs\,u-s^2z.
\end{aligned}
\]

Therefore, relative to the orthonormal basis
$\{u,z\}$, the restriction of $D$ to $\mathcal P$
has the symmetric matrix
\[
\boxed{
M=\begin{bmatrix}
s^2 & -cs\\
-cs & -s^2
\end{bmatrix}.
}
\]

\paragraph{Eigenvalues and spectral norm.}

The characteristic polynomial of $M$ is
\[
\begin{aligned}
\det(M-\lambda I)
&=(s^2-\lambda)(-s^2-\lambda)-c^2s^2\\
&=\lambda^2-s^4-c^2s^2\\
&=\lambda^2-s^2(s^2+c^2)\\
&=\lambda^2-s^2.
\end{aligned}
\]

Since $c^2+s^2=1$, the eigenvalues of $M$ are
\[
\lambda_+=s,
\qquad
\lambda_-=-s.
\]

The operator $D$ is symmetric and vanishes on
$\mathcal P^\perp$. Thus its eigenvalues are
$\pm s$, together with zeros if the ambient
dimension is greater than two.

Its spectral norm is consequently the largest
absolute eigenvalue:
\[
\boxed{
\norm{D}_2=\max\{\abs{s},\abs{-s}\}=s.
}
\]

Finally, $s=\sqrt{1-c^2}$ and
$c^2=(u^Tw)^2$, so
\[
\boxed{
\norm{uu^T-ww^T}_2
=\sqrt{1-(u^Tw)^2}
=\sin\theta.
}
\]

This proves the theorem.

\end{proof}

\paragraph{Geometric meaning.}

The formula makes the behavior of our metric
transparent.

When the lines coincide, $\theta=0$ and $d=0$.
When they are perpendicular, $\theta=\pi/2$
and $d=1$.

Between these extremes, the metric increases
continuously with their angular separation.

\begin{example}[Two Lines in $\R^2$]

Let
\[
u=\begin{bmatrix}1\\0\end{bmatrix},
\qquad
w=\begin{bmatrix}\cos\alpha\\\sin\alpha\end{bmatrix},
\qquad 0\leq\alpha\leq\frac{\pi}{2}.
\]

The two lines meet at angle $\alpha$.
Because $u^Tw=\cos\alpha$, the theorem gives
\[
\boxed{
d(\Span\{u\},\Span\{w\})
=\sqrt{1-\cos^2\alpha}=\sin\alpha.
}
\]

For instance, a $30^\circ$ angle gives distance
$1/2$, while a $90^\circ$ angle gives distance $1$.

\end{example}

This result is the one-dimensional prototype of the
higher-dimensional formula we will develop next.

For two $k$-dimensional subspaces, there is generally
more than one meaningful angle between them. These
are called \emph{principal angles}, and the largest
of them determines the projection distance.

% ============================================================
% 3.5 PRINCIPAL ANGLES BETWEEN SUBSPACES
% ============================================================

\subsection{Principal Angles Between Subspaces}
\label{subsec:principal_angles}

The preceding subsection showed that the distance between
one-dimensional subspaces is determined by the sine of
the angle between them. For higher-dimensional subspaces,
however, a single angle is generally not sufficient to
describe their relative positions.

For example, two planes in $\R^3$ may share an entire line
while tilting away from one another in a second direction.
Along their common line they agree perfectly, but in the
other direction they differ. We therefore need a way to
describe how the subspaces align in several independent
directions.

The appropriate generalization is the notion of
\emph{principal angles}.

\paragraph{The angle between vectors from two subspaces.}

Let $\mathcal W$ and $\mathcal U$ be two $k$-dimensional
subspaces of $\R^n$. Although a vector $w$ may belong to
$\mathcal W$ and a vector $u$ may belong to $\mathcal U$,
both vectors still belong to the same ambient space $\R^n$.
Consequently, their Euclidean inner product and the angle
between them are well-defined.

For nonzero $w$ and $u$, the usual angle $\phi\in[0,\pi]$
satisfies
\[
\cos\phi=\frac{w^Tu}{\norm{w}_2\norm{u}_2}.
\]

For unit vectors, this simplifies to
\[
\cos\phi=w^Tu.
\]

The central idea is to select directions within the two
subspaces that are as closely aligned as possible, and
then repeat the process under suitable orthogonality
constraints.

\begin{definition}[Principal Angles]
\label{def:principal_angles}

Let $\mathcal W,\mathcal U\subseteq\R^n$ be subspaces of
the same dimension $k$. Their \emph{principal angles}
\[
0\leq\theta_1\leq\theta_2\leq\cdots\leq\theta_k
\leq\frac{\pi}{2}
\]
are defined recursively as follows.

For $i=1,\ldots,k$, choose unit vectors
$w_i\in\mathcal W$ and $u_i\in\mathcal U$ that maximize
$w_i^Tu_i$, subject to
\[
\begin{aligned}
&w_i\perp w_j,\qquad j<i,\\
&u_i\perp u_j,\qquad j<i.
\end{aligned}
\]

The $i$th principal angle is defined by
\[
\boxed{\cos\theta_i=w_i^Tu_i.}
\]

The vectors $w_i$ and $u_i$ are called a pair of
\emph{principal vectors}.

\end{definition}

The first pair of principal vectors therefore identifies
the two directions, one in each subspace, that are most
closely aligned. The second pair is chosen after excluding
the first direction from each subspace, and so on.

The first maximum is nonnegative: if $w^Tu<0$, we can
replace $u$ by $-u$. The same observation applies at each
stage. Thus each principal angle lies between $0$ and
$\pi/2$.

Now consider orthonormal bases for the subspaces.
Let
\[
W=\begin{bmatrix}w^{(0)}_1&\cdots&w^{(0)}_k\end{bmatrix},
\qquad
U=\begin{bmatrix}u^{(0)}_1&\cdots&u^{(0)}_k\end{bmatrix}
\]
be matrices whose columns form orthonormal bases of
$\mathcal W$ and $\mathcal U$, respectively. In particular,
\[
W^TW=U^TU=I_k.
\]

Every unit vector in $\mathcal W$ can be written as
$w=Wa$ for some $a\in\R^k$ with $\norm{a}_2=1$.
Similarly, every unit vector in $\mathcal U$ can be
written as $u=Ub$ with $\norm{b}_2=1$.

Their inner product is
\[
\boxed{w^Tu=a^T(W^TU)b.}
\]

Thus, the matrix
\[
\boxed{C=W^TU\in\R^{k\times k}}
\]
contains precisely the information needed to compare
the directions in the two subspaces.

We now establish the connection rigorously.

\begin{theorem}[Principal Angles and Singular Values]
\label{thm:principal_angles_svd}

Let $W,U\in\R^{n\times k}$ have orthonormal columns,
and let $\mathcal W=\range(W)$ and
$\mathcal U=\range(U)$. Write the singular values of
$W^TU$ in nonincreasing order:
\[
\sigma_1\geq\sigma_2\geq\cdots\geq\sigma_k\geq0.
\]

Then the principal angles between $\mathcal W$ and
$\mathcal U$ satisfy
\[
\boxed{\cos\theta_i=\sigma_i(W^TU),
\qquad i=1,\ldots,k.}
\]

In particular,
\[
\boxed{\theta_{\max}=\theta_k
=\cos^{-1}\bigl(\sigma_k(W^TU)\bigr).}
\]

\end{theorem}

\begin{proof}

Set $C=W^TU$. Take a singular value decomposition
\[
\boxed{C=L\Sigma R^T,}
\]
where $L,R\in\R^{k\times k}$ are orthogonal and
\[
\Sigma=\diag(\sigma_1,\ldots,\sigma_k).
\]

Write the columns of $L$ and $R$ as
$\ell_1,\ldots,\ell_k$ and $r_1,\ldots,r_k$.

First notice that all singular values lie in $[0,1]$.
Indeed, $W^T$ cannot increase Euclidean length because
$W$ has orthonormal columns. For every $b\in\R^k$,
\[
\norm{W^TUb}_2\leq\norm{Ub}_2=\norm{b}_2.
\]

Consequently, $\norm{C}_2\leq1$, so
$0\leq\sigma_i\leq1$ for every $i$.

Define vectors in the original ambient space by
\[
\boxed{p_i=W\ell_i,\qquad q_i=Ur_i.}
\]

These vectors belong to $\mathcal W$ and $\mathcal U$,
respectively. Moreover,
\[
\begin{aligned}
p_i^Tp_j&=\ell_i^TW^TW\ell_j
=\ell_i^T\ell_j=\delta_{ij},\\
q_i^Tq_j&=r_i^TU^TUr_j
=r_i^Tr_j=\delta_{ij}.
\end{aligned}
\]

Thus each family is orthonormal within its own subspace.

Their cross-inner products are
\[
\begin{aligned}
p_i^Tq_j
&=\ell_i^TW^TUr_j\\
&=\ell_i^TCr_j\\
&=\ell_i^TL\Sigma R^Tr_j\\
&=\sigma_i\delta_{ij}.
\end{aligned}
\]

In particular,
\[
\boxed{p_i^Tq_i=\sigma_i.}
\]

It remains to show that these pairs attain the maxima
in the recursive definition of principal angles.

For the first pair, let $w=Wa$ and $u=Ub$ be arbitrary
unit vectors. Write $a=L\alpha$ and $b=R\beta$.
Because $L$ and $R$ are orthogonal,
$\norm{\alpha}_2=\norm{\beta}_2=1$. Then
\[
\begin{aligned}
w^Tu
&=a^TCb\\
&=\alpha^T\Sigma\beta\\
&=\sum_{j=1}^k\sigma_j\alpha_j\beta_j\\
&\leq\sigma_1\sum_{j=1}^k|\alpha_j\beta_j|\\
&\leq\sigma_1\norm{\alpha}_2\norm{\beta}_2\\
&=\sigma_1.
\end{aligned}
\]

The second inequality is Cauchy--Schwarz. Equality is
attained by $a=\ell_1$ and $b=r_1$, which give
$w=p_1$ and $u=q_1$. Hence
\[
\cos\theta_1=\sigma_1.
\]

Now suppose the first $i-1$ principal pairs have been
chosen as $p_j,q_j$ for $j<i$. The orthogonality
constraints on the next pair require
\[
w\perp p_j,\qquad u\perp q_j\qquad(j<i).
\]

Since $w=WL\alpha$ and $u=UR\beta$, these conditions
are equivalent to
\[
\alpha_j=\beta_j=0\qquad(j<i).
\]

Therefore the inner product becomes
\[
\begin{aligned}
w^Tu
&=\sum_{j=i}^k\sigma_j\alpha_j\beta_j\\
&\leq\sigma_i\sum_{j=i}^k|\alpha_j\beta_j|\\
&\leq\sigma_i.
\end{aligned}
\]

The value $\sigma_i$ is attained by choosing
$\alpha=\beta=e_i$, giving $w=p_i$ and $u=q_i$.
Thus
\[
\cos\theta_i=\sigma_i
\qquad(i=1,\ldots,k).
\]

Since cosine decreases on $[0,\pi/2]$, decreasing
singular values correspond to increasing principal
angles. This proves the theorem.

\end{proof}

\paragraph{Now we will show that these angles do not depend on a basis.}

Suppose $\widetilde W=WQ_1$ and
$\widetilde U=UQ_2$ are different orthonormal bases
for the same respective subspaces, with $Q_1,Q_2$
orthogonal. Then
\[
\widetilde W^T\widetilde U=Q_1^T(W^TU)Q_2.
\]

Left or right multiplication by an orthogonal matrix
does not change singular values. Hence the angles
computed from $\widetilde W^T\widetilde U$ agree with
those from $W^TU$, just as a geometric notion of
subspace angle should.

\begin{example}[Two Planes with a Common Line]
\label{ex:principal_angles_planes}

In $\R^3$, let
\[
\mathcal W=\Span\{e_1,e_2\},\qquad
\mathcal U=\Span\left\{e_1,
\frac{e_2+e_3}{\sqrt2}\right\},
\]
where $e_1,e_2,e_3$ are the standard coordinate vectors.

Choose orthonormal basis matrices
\[
W=\begin{bmatrix}
1&0\\0&1\\0&0
\end{bmatrix},
\qquad
U=\begin{bmatrix}
1&0\\0&1/\sqrt2\\0&1/\sqrt2
\end{bmatrix}.
\]

Then
\[
\boxed{
W^TU=\begin{bmatrix}1&0\\0&1/\sqrt2\end{bmatrix}.
}
\]

Its singular values are $\sigma_1=1$ and
$\sigma_2=1/\sqrt2$. Hence
\[
\boxed{\theta_1=0,
\qquad\theta_2=\frac{\pi}{4}=45^\circ.}
\]

The first angle is zero because the planes share the
line $\Span\{e_1\}$. The second angle is $45^\circ$
because their remaining orthogonal directions are
$e_2$ and $(e_2+e_3)/\sqrt2$.

Thus, although the two planes intersect, they are not
the same plane. Their two principal angles distinguish
the common direction from the direction in which the
planes separate.

\end{example}

\begin{remark}[Special Cases and Nonuniqueness]

If all principal angles vanish, then the two
$k$-dimensional subspaces coincide. If a principal angle
is $\pi/2$, the corresponding principal directions are
orthogonal. Repeated singular values may make the
individual principal vectors nonunique, but the list
of principal angles remains well-defined.

For $k=1$, the matrix $W^TU$ has one entry, $w^Tu$.
Its only singular value is $|w^Tu|$, so the theorem
recovers the acute-angle definition used in the
preceding subsection.

\end{remark}

We have now identified the complete angular
information encoded by $W^TU$. One question remains:
why does the \emph{largest} principal angle, rather than
the smallest or an average of the angles, determine the
spectral-norm distance between the corresponding
orthogonal projectors? We will answer this next.



% ============================================================
% 3.6 SUBSPACE DISTANCE AND PRINCIPAL ANGLES
% ============================================================

\subsection{Subspace Distance and Principal Angles}
\label{subsec:distance_principal_angles}

We have developed two descriptions of the relative position
of subspaces. The first uses orthogonal projection matrices,
while the second uses principal angles. We now prove that
they describe the same geometric phenomenon.

For two lines, the relationship was
\[
\norm{uu^T-ww^T}_2=\sin\theta.
\]
Our aim is to obtain the corresponding result for two
subspaces of the same dimension, without assuming that
they share any particular coordinate directions.

\begin{theorem}[Projection Distance and the Largest Principal Angle]
\label{thm:projection_distance_principal_angles}

Let $\mathcal W,\mathcal U\subseteq\R^n$ be subspaces of
the same dimension $k\geq1$. Let $W,U\in\R^{n\times k}$
have orthonormal columns spanning $\mathcal W$ and
$\mathcal U$, respectively. Write their principal angles as
\[
0\leq\theta_1\leq\cdots\leq\theta_k\leq\frac{\pi}{2}.
\]

Then
\[
\boxed{
\begin{aligned}
d(\mathcal W,\mathcal U)
&=\norm{WW^T-UU^T}_2\\[3pt]
&=\sqrt{1-\sigma_k^2(W^TU)}\\[3pt]
&=\sin\theta_k.
\end{aligned}
}
\]

In particular, the projection distance is determined by
the \emph{largest} principal angle.

\end{theorem}

\begin{proof}

To keep the notation simple, write
\[
P=WW^T,\qquad Q=UU^T,\qquad D=P-Q.
\]

Because $P$ and $Q$ are orthogonal projectors,
\[
P^T=P,\quad Q^T=Q,
\qquad P^2=P,\quad Q^2=Q.
\]

Consequently, $D$ is symmetric. Its spectral norm is
therefore related to the square of the operator by
\[
\boxed{\norm{D}_2^2=\norm{D^2}_2.}
\]

Indeed, if the eigenvalues of $D$ are $\lambda_i$, those
of $D^2$ are $\lambda_i^2$, and the spectral norm of a
symmetric matrix is its largest absolute eigenvalue.

\paragraph{Step 1: Decomposing the square of the projector difference.}

Expanding and using idempotence gives
\[
\begin{aligned}
D^2
&=(P-Q)(P-Q)\\
&=P^2-PQ-QP+Q^2\\
&=P+Q-PQ-QP.
\end{aligned}
\]

We now rewrite this expression as
\[
\boxed{
D^2=P(I-Q)P+(I-P)Q(I-P).
}
\]

To verify the identity, expand the right-hand side:
\[
\begin{aligned}
&P(I-Q)P+(I-P)Q(I-P)\\
&\quad=P-PQP+Q-PQ-QP+PQP\\
&\quad=P+Q-PQ-QP=D^2.
\end{aligned}
\]

The purpose of this decomposition is geometric: its two
terms act on mutually orthogonal parts of $\R^n$.

\paragraph{Step 2: Understanding the two orthogonal components.}

Set
\[
H=P(I-Q)P,
\qquad
K=(I-P)Q(I-P).
\]

The first operator has a projector $P$ on its left.
Therefore $Hx\in\mathcal W$ for every $x\in\R^n$.
It also has $P$ on its right, so if
$x\in\mathcal W^\perp$, then $Px=0$ and hence $Hx=0$.

Conversely, the second operator has $I-P$ on its left,
so $Kx\in\mathcal W^\perp$ for every $x$. It also
has $I-P$ on its right, so $Kx=0$ whenever
$x\in\mathcal W$.

Thus, relative to the orthogonal decomposition
\[
\R^n=\mathcal W\oplus\mathcal W^\perp,
\]
the matrix $D^2=H+K$ has two diagonal blocks and no
cross terms. Therefore
\[
\boxed{
\norm{D^2}_2=\max\{\norm{H}_2,\norm{K}_2\}.
}
\]

This equality is worth emphasizing: $H$ and $K$ do not
compete in the same direction. Each acts on one
orthogonal component and vanishes on the other.

Next, the identities $(I-Q)^2=I-Q$ and $Q^2=Q$ give
\[
\begin{aligned}
H&=[(I-Q)P]^T[(I-Q)P],\\
K&=[Q(I-P)]^T[Q(I-P)].
\end{aligned}
\]

For any matrix $B$, we have
$\norm{B^TB}_2=\norm{B}_2^2$. Hence
\[
\begin{aligned}
\norm{H}_2&=\norm{(I-Q)P}_2^2,\\
\norm{K}_2&=\norm{Q(I-P)}_2^2
              =\norm{(I-P)Q}_2^2,
\end{aligned}
\]
where the last equality uses
$\norm{B^T}_2=\norm{B}_2$.

Putting these observations together, we obtain
\[
\boxed{
\norm{P-Q}_2^2
=\max\!\left\{
\norm{(I-Q)P}_2^2,
\norm{(I-P)Q}_2^2
\right\}.
}
\]

We have reduced the difference of two projectors to
two quantities measuring how much one subspace
projects outside the other.

\paragraph{Step 3: Expressing the two quantities through singular values.}

Because $P=WW^T$ and the columns of $W$ are
orthonormal,
\[
\boxed{
\norm{(I-Q)P}_2=\norm{(I-Q)W}_2.
}
\]

For completeness, let us justify this equality. For
any unit vector $x\in\R^n$,
\[
\begin{aligned}
\norm{(I-Q)Px}_2
&=\norm{(I-Q)WW^Tx}_2\\
&\leq\norm{(I-Q)W}_2\,\norm{W^Tx}_2\\
&\leq\norm{(I-Q)W}_2.
\end{aligned}
\]

Conversely, choose $x=Wa$ with $\norm{a}_2=1$.
Then $\norm{x}_2=1$, $Px=x$, and
\[
\norm{(I-Q)Px}_2=\norm{(I-Q)Wa}_2.
\]

Maximizing over $a$ proves the reverse inequality.
The same argument, with $U$ in place of $W$, yields
\[
\boxed{
\norm{(I-P)Q}_2=\norm{(I-P)U}_2.
}
\]

Now introduce the $k\times k$ matrix
\[
C=W^TU.
\]

The squared norm of $(I-P)U$ satisfies
\[
\begin{aligned}
\norm{(I-P)U}_2^2
&=\norm{U^T(I-P)U}_2\\
&=\norm{I_k-U^TWW^TU}_2\\
&=\norm{I_k-C^TC}_2.
\end{aligned}
\]

Since $C^TC$ is symmetric with eigenvalues
$\sigma_1^2(C),\ldots,\sigma_k^2(C)$, and
$0\leq\sigma_i(C)\leq1$, the eigenvalues of
$I_k-C^TC$ are
\[
1-\sigma_1^2(C),\ldots,1-\sigma_k^2(C).
\]

These numbers are nonnegative. Because the singular
values are ordered in decreasing order, the largest
of these numbers is $1-\sigma_k^2(C)$. Thus
\[
\boxed{
\norm{(I-P)U}_2^2=1-\sigma_k^2(C).
}
\]

Similarly,
\[
\begin{aligned}
\norm{(I-Q)W}_2^2
&=\norm{W^T(I-Q)W}_2\\
&=\norm{I_k-W^TUU^TW}_2\\
&=\norm{I_k-CC^T}_2\\
&=1-\sigma_k^2(C).
\end{aligned}
\]

The last equality follows because $CC^T$ and $C^TC$
have the same eigenvalues: the squares of the singular
values of the square matrix $C$.

Therefore, for equal-dimensional subspaces,
\[
\boxed{
\norm{(I-Q)P}_2^2
=\norm{(I-P)Q}_2^2
=1-\sigma_k^2(W^TU).
}
\]

Substituting into the identity from Step 2 gives
\[
\norm{P-Q}_2^2=1-\sigma_k^2(W^TU).
\]

Since norms are nonnegative,
\[
\boxed{
\norm{P-Q}_2=\sqrt{1-\sigma_k^2(W^TU)}.
}
\]

Finally, by Theorem~\ref{thm:principal_angles_svd},
\[
\sigma_k(W^TU)=\cos\theta_k.
\]

Consequently,
\[
\begin{aligned}
\norm{P-Q}_2
&=\sqrt{1-\cos^2\theta_k}\\
&=\sin\theta_k,
\end{aligned}
\]
where $\theta_k\in[0,\pi/2]$ ensures that its sine
is nonnegative. Hence
\[
\boxed{
d(\mathcal W,\mathcal U)
=\norm{WW^T-UU^T}_2
=\sin\theta_{\max}.
}
\]

This completes the proof.
\end{proof}

\paragraph{Geometric interpretation.}

We can extract another useful identity from the proof:
\[
\boxed{
\begin{aligned}
d(\mathcal W,\mathcal U)
&=\norm{(I-UU^T)W}_2\\
&=\norm{(I-WW^T)U}_2.
\end{aligned}
}
\]

For example, $(I-UU^T)W$ measures the components of
unit vectors in $\mathcal W$ that lie perpendicular
to $\mathcal U$. Its norm is the largest such component.
Thus, the largest principal angle determines the
\emph{worst-case misalignment} of the two subspaces.

\begin{example}[Revisiting Two Planes in $\R^3$]
\label{ex:projection_distance_planes}

Consider the two planes from
Example~\ref{ex:principal_angles_planes}:
\[
\mathcal W=\Span\{e_1,e_2\},
\qquad
\mathcal U=\Span\left\{e_1,\frac{e_2+e_3}{\sqrt2}\right\}.
\]

Their orthogonal projectors are
\[
P=\begin{bmatrix}
1&0&0\\
0&1&0\\
0&0&0
\end{bmatrix},
\qquad
Q=\begin{bmatrix}
1&0&0\\
0&\tfrac12&\tfrac12\\
0&\tfrac12&\tfrac12
\end{bmatrix}.
\]

Therefore,
\[
P-Q=\begin{bmatrix}
0&0&0\\
0&\tfrac12&-\tfrac12\\
0&-\tfrac12&-\tfrac12
\end{bmatrix}.
\]

The lower-right $2\times2$ block has eigenvalues
$\pm1/\sqrt2$; the remaining eigenvalue is zero.
Consequently,
\[
\boxed{\norm{P-Q}_2=\frac1{\sqrt2}.}
\]

This agrees with the principal-angle calculation,
because $\theta_1=0$ and $\theta_2=\pi/4$:
\[
\boxed{
\sin\theta_{\max}=\sin\frac\pi4=\frac1{\sqrt2}.
}
\]

The shared direction $e_1$ contributes no separation.
The distance is determined entirely by the direction
in which the two planes tilt apart.

\end{example}

\begin{remark}[What If the Dimensions Are Different?]
\label{rem:different_subspace_dimensions}

So far, our relationship between subspace distance
and principal angles has been developed for two
subspaces of the same dimension.

However, the projection-based metric
\[
d(\mathcal U,\mathcal W)
=
\norm{P_{\mathcal U}-P_{\mathcal W}}_2
\]
remains well-defined even when the dimensions of
the two subspaces are different.

Interestingly, in this case, the distance
always attains its maximum possible value:
\[
\boxed{
\dim\mathcal U\neq\dim\mathcal W
\quad\Longrightarrow\quad
d(\mathcal U,\mathcal W)=1.
}
\]

Let us explain why this happens.

\paragraph{Geometric intuition.}

Suppose that
\[
\dim\mathcal U>\dim\mathcal W.
\]

The subspace $\mathcal U$ contains more linearly
independent directions than $\mathcal W$.

Consequently, there must exist at least one
nonzero direction in $\mathcal U$ that is
orthogonal to every vector in $\mathcal W$.

This means that we can find a unit vector
$u\in\R^n$ satisfying
\[
\boxed{
u\in\mathcal U
\qquad\text{and}\qquad
u\in\mathcal W^\perp.
}
\]

In other words,
\[
\boxed{
u\in\mathcal U\cap\mathcal W^\perp.
}
\]

We now justify this observation using a
simple dimension-counting argument.

\paragraph{Why must such a vector exist?}

Recall the dimension formula for two
subspaces $\mathcal A,\mathcal B\subseteq\R^n$:
\[
\dim(\mathcal A+\mathcal B)
=
\dim\mathcal A+\dim\mathcal B
-\dim(\mathcal A\cap\mathcal B).
\]

Since
\[
\dim(\mathcal A+\mathcal B)\leq n,
\]
we obtain
\[
\boxed{
\dim(\mathcal A\cap\mathcal B)
\geq
\dim\mathcal A+\dim\mathcal B-n.
}
\]

Now choose
\[
\mathcal A=\mathcal U,
\qquad
\mathcal B=\mathcal W^\perp.
\]

Because
\[
\dim\mathcal W^\perp
=
n-\dim\mathcal W,
\]
the dimension inequality gives
\[
\begin{aligned}
\dim(\mathcal U\cap\mathcal W^\perp)
&\geq
\dim\mathcal U+\dim\mathcal W^\perp-n\\
&=
\dim\mathcal U+(n-\dim\mathcal W)-n\\
&=
\dim\mathcal U-\dim\mathcal W.
\end{aligned}
\]

But our assumption is that
\[
\dim\mathcal U>\dim\mathcal W.
\]

Therefore,
\[
\boxed{
\dim(\mathcal U\cap\mathcal W^\perp)>0.
}
\]

Hence, the intersection contains a nonzero
vector, which we may normalize to obtain
a unit vector $u$.

This proves the existence of the direction
described above.

\paragraph{Step 1: Establishing the lower bound.}

Since $u\in\mathcal U$, its orthogonal
projection onto $\mathcal U$ is itself:
\[
P_{\mathcal U}u=u.
\]

On the other hand, because
$u\in\mathcal W^\perp$, its orthogonal
projection onto $\mathcal W$ is zero:
\[
P_{\mathcal W}u=0.
\]

Consequently,
\[
\begin{aligned}
(P_{\mathcal U}-P_{\mathcal W})u
&=P_{\mathcal U}u-P_{\mathcal W}u\\
&=u-0\\
&=u.
\end{aligned}
\]

Taking Euclidean norms gives
\[
\norm{(P_{\mathcal U}-P_{\mathcal W})u}_2
=
\norm{u}_2
=
1.
\]

Recall that the spectral norm is defined by
\[
\norm{B}_2
=
\max_{\norm{x}_2=1}\norm{Bx}_2.
\]

Since we have found a unit vector $u$
for which the projector difference has
norm exactly one, its spectral norm
cannot be smaller than one.

Thus,
\[
\boxed{
\norm{P_{\mathcal U}-P_{\mathcal W}}_2
\geq1.
}
\]

\paragraph{Step 2: Establishing the upper bound.}

We now show that the difference of two
orthogonal projectors can never have
spectral norm greater than one.

Let $x\in\R^n$ be any unit vector.

Because orthogonal projectors are symmetric
and idempotent,
\[
\begin{aligned}
x^TP_{\mathcal U}x
&=x^TP_{\mathcal U}^TP_{\mathcal U}x\\
&=\norm{P_{\mathcal U}x}_2^2.
\end{aligned}
\]

Similarly,
\[
x^TP_{\mathcal W}x
=
\norm{P_{\mathcal W}x}_2^2.
\]

Orthogonal projection cannot increase
the Euclidean length of a vector.

Since $\norm{x}_2=1$, we have
\[
0\leq\norm{P_{\mathcal U}x}_2^2\leq1
\]
and
\[
0\leq\norm{P_{\mathcal W}x}_2^2\leq1.
\]

Therefore,
\[
\begin{aligned}
x^T(P_{\mathcal U}-P_{\mathcal W})x
&=
\norm{P_{\mathcal U}x}_2^2\\
&\quad-
\norm{P_{\mathcal W}x}_2^2.
\end{aligned}
\]

Since both terms lie in $[0,1]$,
their difference must lie in $[-1,1]$.

Hence,
\[
\boxed{
\abs{x^T(P_{\mathcal U}-P_{\mathcal W})x}
\leq1.
}
\]

Now observe that
$P_{\mathcal U}-P_{\mathcal W}$
is symmetric.

For a real symmetric matrix $B$,
the spectral theorem and the Rayleigh
quotient characterization imply that
\[
\norm{B}_2
=
\max_{\norm{x}_2=1}\abs{x^TBx}.
\]

Applying this identity gives
\[
\boxed{
\norm{P_{\mathcal U}-P_{\mathcal W}}_2
\leq1.
}
\]

\paragraph{Conclusion.}

Combining the lower and upper bounds,
we obtain
\[
1
\leq
\norm{P_{\mathcal U}-P_{\mathcal W}}_2
\leq1.
\]

Therefore,
\[
\boxed{
d(\mathcal U,\mathcal W)
=
\norm{P_{\mathcal U}-P_{\mathcal W}}_2
=1.
}
\]

The case
$\dim\mathcal W>\dim\mathcal U$
follows by interchanging the two
subspaces.

\paragraph{A simple example.}

Consider the subspaces of $\R^3$
\[
\mathcal W=\Span\{e_1\},
\qquad
\mathcal U=\Span\{e_1,e_2\}.
\]

Here,
\[
\dim\mathcal W=1,
\qquad
\dim\mathcal U=2.
\]

Notice that
\[
e_2\in\mathcal U,
\qquad
e_2\perp\mathcal W.
\]

The corresponding projectors are
\[
P_{\mathcal W}
=
\begin{bmatrix}
1&0&0\\
0&0&0\\
0&0&0
\end{bmatrix},
\qquad
P_{\mathcal U}
=
\begin{bmatrix}
1&0&0\\
0&1&0\\
0&0&0
\end{bmatrix}.
\]

Thus,
\[
P_{\mathcal U}-P_{\mathcal W}
=
\begin{bmatrix}
0&0&0\\
0&1&0\\
0&0&0
\end{bmatrix}.
\]

Its spectral norm is one, confirming that
\[
\boxed{
d(\mathcal U,\mathcal W)=1.
}
\]

Notice that $\mathcal W$ is actually
contained in $\mathcal U$.

Nevertheless, their projection distance
is maximal because the larger subspace
contains an additional direction that
is completely orthogonal to the smaller
subspace.

\paragraph{Connection to the Courant--Fischer theorem.}

The key ingredient in this argument was
the dimension inequality
\[
\boxed{
\dim(\mathcal A\cap\mathcal B)
\geq
\dim\mathcal A+\dim\mathcal B-n.
}
\]

In particular, whenever
\[
\dim\mathcal A+\dim\mathcal B>n,
\]
the two subspaces must have a
nontrivial intersection.

This seemingly elementary observation
will play an important role when we
study the Courant--Fischer min--max
and max--min principles.


\end{remark}
We have completed the connection between the
projection-based metric, the SVD, and the geometry of
principal angles. These ideas provide the tools needed
to study how invariant subspaces move when the matrix
defining them is perturbed, a question we shall revisit
in the Davis--Kahan theorem.
% ============================================================
% SECTION 4: COURANT--FISCHER THEOREM
% ============================================================
% ============================================================
% SECTION 4: COURANT--FISCHER THEOREM
% ============================================================
\begin{comment}
    

\section{The Courant--Fischer Theorem}
\label{sec:courant_fischer}

The Rayleigh quotient provides a variational description
of the smallest and largest eigenvalues of a real symmetric
matrix. We now extend that description to all the
intermediate eigenvalues.

This extension is especially interesting because it
connects three seemingly different ideas: the spectral
theorem, optimization, and the geometry of subspaces.
The essential tool will be a dimension-counting argument
that we have already encountered when studying subspace
distance.

% ============================================================
% 4.1 MOTIVATION
% ============================================================

\subsection{From Extreme Eigenvalues to Intermediate Eigenvalues}
\label{subsec:courant_motivation}

Let $A\in\R^{n\times n}$ be real and symmetric, with
orthonormal eigenvectors $v_1,\ldots,v_n$ and eigenvalues
ordered as
\[
\lambda_1\leq\lambda_2\leq\cdots\leq\lambda_n.
\]

Recall that the Rayleigh quotient is
\[
R_A(x)=\frac{x^TAx}{x^Tx},\qquad x\neq0.
\]

In \cref{subsec:rayleigh_background}, we proved that
\[
\boxed{
\lambda_1=\min_{x\neq0}R_A(x),\qquad
\lambda_n=\max_{x\neq0}R_A(x).
}
\]

These formulas characterize only the two extreme
eigenvalues. To understand an intermediate eigenvalue,
say $\lambda_k$, it is tempting to try to optimize the
Rayleigh quotient while excluding the eigendirections
associated with smaller or larger eigenvalues.

For example, consider
\[
\mathcal S_k=\Span\{v_1,\ldots,v_k\}.
\]
If $x\in\mathcal S_k$ is nonzero, then
\[
x=\sum_{i=1}^k c_i v_i,
\]
and the spectral theorem gives
\[
R_A(x)=
\frac{\sum_{i=1}^k\lambda_i c_i^2}
     {\sum_{i=1}^k c_i^2}.
\]
Since every eigenvalue appearing in this expression
is at most $\lambda_k$,
\[
R_A(x)\leq\lambda_k.
\]
Equality occurs at $x=v_k$. Therefore,
\[
\boxed{
\max_{\substack{x\in\mathcal S_k\\x\neq0}}
R_A(x)=\lambda_k.
}
\]

Now consider instead
\[
\mathcal T_k=\Span\{v_k,\ldots,v_n\},
\]
which has dimension $n-k+1$. A similar argument shows
\[
\boxed{
\min_{\substack{x\in\mathcal T_k\\x\neq0}}
R_A(x)=\lambda_k.
}
\]

These identities suggest that $\lambda_k$ is accessible
by optimizing the Rayleigh quotient over suitable
subspaces. But they refer to spaces constructed from
eigenvectors we may not know in advance.

\textbf{The surprising consequence of the Courant--Fischer
theorem is that we can optimize over \emph{all} subspaces
of the appropriate dimension and still recover
exactly $\lambda_k$.}

% ============================================================
% 4.2 DIMENSION COUNTING
% ============================================================

\subsection{The Geometric Ingredient: Intersecting Subspaces}
\label{subsec:courant_dimension}

Before stating the theorem, let us explain the
geometric observation on which its proof rests.

\begin{lemma}[Dimension-Counting Principle]
\label{lem:courant_dimension}
Let $\mathcal S,\mathcal T$ be subspaces of $\R^n$.
If
\[
\dim\mathcal S+\dim\mathcal T>n,
\]
then
\[
\boxed{\mathcal S\cap\mathcal T\neq\{0\}.}
\]
\end{lemma}

\begin{proof}
The dimension formula states that
\[
\dim(\mathcal S+\mathcal T)
=\dim\mathcal S+\dim\mathcal T
 -\dim(\mathcal S\cap\mathcal T).
\]
Because $\mathcal S+\mathcal T\subseteq\R^n$,
\[
\dim(\mathcal S+\mathcal T)\leq n.
\]
Rearranging the dimension formula therefore gives
\[
\dim(\mathcal S\cap\mathcal T)
\geq\dim\mathcal S+\dim\mathcal T-n.
\]
The right-hand side is strictly positive by
assumption. Hence the intersection has positive
dimension and contains a nonzero vector.
\end{proof}

For our proof, the key dimension count will be
\[
\boxed{k+(n-k+1)=n+1>n.}
\]
Thus a $k$-dimensional subspace and an
$(n-k+1)$-dimensional subspace of $\R^n$ must
share a nonzero vector.

Notice what this does \emph{not} say: the two subspaces
need not be equal, and we do not need to know their
bases. Their dimensions alone guarantee a nonzero
intersection. This is the same idea used in
\cref{rem:different_subspace_dimensions}.

% ============================================================
% 4.3 THEOREM AND PROOF
% ============================================================

\subsection{The Min--Max and Max--Min Principles}
\label{subsec:courant_theorem}

\begin{theorem}[Courant--Fischer]
\label{thm:courant_fischer}
Let $A\in\R^{n\times n}$ be real and symmetric,
with eigenvalues ordered as
$\lambda_1\leq\cdots\leq\lambda_n$.
For every $k\in\{1,\ldots,n\}$,
\[
\boxed{
\lambda_k
=\min_{\substack{\mathcal S\subseteq\R^n\\
                 \dim\mathcal S=k}}
\ \max_{\substack{x\in\mathcal S\\x\neq0}}R_A(x).
}
\]
Equivalently,
\[
\boxed{
\lambda_k
=\max_{\substack{\mathcal T\subseteq\R^n\\
                 \dim\mathcal T=n-k+1}}
\ \min_{\substack{x\in\mathcal T\\x\neq0}}R_A(x).
}
\]
The first formula is the \emph{min--max principle};
the second is the \emph{max--min principle}.
\end{theorem}

\begin{proof}
Choose an orthonormal eigenbasis
$v_1,\ldots,v_n$ such that
\[
Av_i=\lambda_i v_i,
\qquad
\lambda_1\leq\cdots\leq\lambda_n.
\]
We prove the two statements separately.

\medskip
\noindent\textbf{Part I: Proof of the min--max principle.}

We wish to show
\[
\min_{\dim\mathcal S=k}
\ \max_{\substack{x\in\mathcal S\\x\neq0}}R_A(x)
=\lambda_k.
\]

\emph{Step 1: Find a subspace for which the maximum
is $\lambda_k$.}

Take
\[
\mathcal S_0=\Span\{v_1,\ldots,v_k\}.
\]
It has dimension $k$. For any nonzero
$x=\sum_{i=1}^k c_i v_i$,
\[
R_A(x)
=\frac{\sum_{i=1}^k\lambda_i c_i^2}
       {\sum_{i=1}^k c_i^2}
\leq\lambda_k.
\]
Equality is attained at $x=v_k$. Hence
\[
\max_{\substack{x\in\mathcal S_0\\x\neq0}}R_A(x)
=\lambda_k.
\]
Because $\mathcal S_0$ is one of the subspaces
allowed in the outer minimization,
\[
\boxed{
\min_{\dim\mathcal S=k}
\ \max_{\substack{x\in\mathcal S\\x\neq0}}R_A(x)
\leq\lambda_k.
}
\]

\emph{Step 2: Show that no other $k$-dimensional
subspace can have a smaller maximum.}

Let $\mathcal S$ be an arbitrary $k$-dimensional
subspace. Introduce
\[
\mathcal T_0=\Span\{v_k,v_{k+1},\ldots,v_n\}.
\]
This subspace has dimension $n-k+1$. Therefore,
\[
\dim\mathcal S+\dim\mathcal T_0
=k+(n-k+1)=n+1>n.
\]
By \cref{lem:courant_dimension}, there exists a
nonzero vector
\[
\boxed{x\in\mathcal S\cap\mathcal T_0.}
\]
Since $x\in\mathcal T_0$, it has the expansion
\[
x=\sum_{i=k}^n c_i v_i.
\]
All eigenvalues occurring in this expansion are
at least $\lambda_k$. Consequently,
\[
R_A(x)
=\frac{\sum_{i=k}^n\lambda_i c_i^2}
       {\sum_{i=k}^n c_i^2}
\geq\lambda_k.
\]
But $x$ also belongs to $\mathcal S$. Its Rayleigh
quotient cannot exceed the maximum over $\mathcal S$:
\[
\max_{\substack{y\in\mathcal S\\y\neq0}}R_A(y)
\geq R_A(x)\geq\lambda_k.
\]
This holds for \emph{every} $k$-dimensional subspace
$\mathcal S$, so
\[
\boxed{
\min_{\dim\mathcal S=k}
\ \max_{\substack{y\in\mathcal S\\y\neq0}}R_A(y)
\geq\lambda_k.
}
\]
Combining the two bounds proves
\[
\boxed{
\lambda_k
=\min_{\dim\mathcal S=k}
\ \max_{\substack{x\in\mathcal S\\x\neq0}}R_A(x).
}
\]

\medskip
\noindent\textbf{Part II: Proof of the max--min principle.}

We now wish to show
\[
\max_{\dim\mathcal T=n-k+1}
\ \min_{\substack{x\in\mathcal T\\x\neq0}}R_A(x)
=\lambda_k.
\]

\emph{Step 1: Find a subspace for which the minimum
is $\lambda_k$.}

Take
\[
\mathcal T_0=\Span\{v_k,\ldots,v_n\}.
\]
For every nonzero $x\in\mathcal T_0$,
\[
R_A(x)
=\frac{\sum_{i=k}^n\lambda_i c_i^2}
       {\sum_{i=k}^n c_i^2}
\geq\lambda_k.
\]
Equality is attained at $x=v_k$. Hence
\[
\min_{\substack{x\in\mathcal T_0\\x\neq0}}R_A(x)
=\lambda_k.
\]
Since this is one candidate in the outer maximization,
\[
\boxed{
\max_{\dim\mathcal T=n-k+1}
\ \min_{\substack{x\in\mathcal T\\x\neq0}}R_A(x)
\geq\lambda_k.
}
\]

\emph{Step 2: Show that no other subspace of this
dimension can have a larger minimum.}

Let $\mathcal T$ be an arbitrary subspace of dimension
$n-k+1$, and choose
\[
\mathcal S_0=\Span\{v_1,\ldots,v_k\}.
\]
Again,
\[
\dim\mathcal T+\dim\mathcal S_0
=(n-k+1)+k=n+1>n.
\]
There must exist a nonzero vector
\[
\boxed{x\in\mathcal T\cap\mathcal S_0.}
\]
Since $x\in\mathcal S_0$, it is a linear combination
of $v_1,\ldots,v_k$, so
\[
R_A(x)
=\frac{\sum_{i=1}^k\lambda_i c_i^2}
       {\sum_{i=1}^k c_i^2}
\leq\lambda_k.
\]
Because $x\in\mathcal T$ as well,
\[
\min_{\substack{y\in\mathcal T\\y\neq0}}R_A(y)
\leq R_A(x)\leq\lambda_k.
\]
As this is true for every $(n-k+1)$-dimensional
subspace $\mathcal T$, we conclude that
\[
\boxed{
\max_{\dim\mathcal T=n-k+1}
\ \min_{\substack{y\in\mathcal T\\y\neq0}}R_A(y)
\leq\lambda_k.
}
\]
Combining both inequalities establishes
\[
\boxed{
\lambda_k
=\max_{\dim\mathcal T=n-k+1}
\ \min_{\substack{x\in\mathcal T\\x\neq0}}R_A(x).
}
\]
Both formulations are proved.
\end{proof}

\begin{remark}[The Role of Symmetry]
The argument uses the spectral theorem to expand vectors
in an orthonormal eigenbasis and to bound their Rayleigh
quotients by the relevant eigenvalues. Consequently,
these formulas hold as stated for real symmetric matrices;
we must not apply them indiscriminately to nonsymmetric
matrices. Compare \cref{rem:rayleigh_symmetry}.
\end{remark}

% ============================================================
% 4.4 WORKED EXAMPLE
% ============================================================

\subsection{A Concrete Example in Three Dimensions}
\label{subsec:courant_example}

Consider the symmetric matrix
\[
A=\begin{bmatrix}
1&0&0\\
0&3&0\\
0&0&7
\end{bmatrix}.
\]
Its eigenvalues, in ascending order, are
\[
\lambda_1=1,\qquad\lambda_2=3,\qquad\lambda_3=7,
\]
with corresponding orthonormal eigenvectors
$e_1,e_2,e_3$.

We will recover the middle eigenvalue $\lambda_2=3$
without optimizing over individual vectors in all
of $\R^3$.

\paragraph{The min--max viewpoint.}
Here $k=2$, so we consider all two-dimensional
subspaces $\mathcal S$ of $\R^3$ and ask how small
we can make the largest Rayleigh quotient within one.

Take $\mathcal S_0=\Span\{e_1,e_2\}$. Every nonzero
$x\in\mathcal S_0$ has the form $x=ae_1+be_2$, giving
\[
R_A(x)=\frac{a^2+3b^2}{a^2+b^2}\leq3.
\]
Equality holds when $x=e_2$. Thus the maximum on this
plane is exactly $3$.

Why can no other plane have a maximum below $3$?
Any two-dimensional plane $\mathcal S$ intersects the
plane $\Span\{e_2,e_3\}$ in a nonzero vector, because
$2+2>3$. Along that vector the Rayleigh quotient is
at least $3$. Hence
\[
\boxed{
\min_{\dim\mathcal S=2}
\ \max_{\substack{x\in\mathcal S\\x\neq0}}R_A(x)=3.
}
\]

\paragraph{The max--min viewpoint.}
Now take $\mathcal T_0=\Span\{e_2,e_3\}$. For
$x=be_2+ce_3\neq0$,
\[
R_A(x)=\frac{3b^2+7c^2}{b^2+c^2}\geq3,
\]
with equality at $x=e_2$. Hence the minimum on
this plane equals $3$.

Every other two-dimensional plane $\mathcal T$
intersects $\Span\{e_1,e_2\}$ in a nonzero vector.
At that vector the Rayleigh quotient is at most $3$,
so no such plane can have a minimum exceeding $3$.
Therefore,
\[
\boxed{
\max_{\dim\mathcal T=2}
\ \min_{\substack{x\in\mathcal T\\x\neq0}}R_A(x)=3.
}
\]

Both optimization problems select the same eigenvalue,
although one minimizes a maximum and the other
maximizes a minimum.

\begin{remark}[Repeated Eigenvalues]
The Courant--Fischer theorem does not require distinct
eigenvalues. If some eigenvalues are repeated, we may
still choose an orthonormal eigenbasis and order the
eigenvalues with multiplicity. The dimension-counting
proof is unchanged, although the subspaces attaining
the extrema may no longer be unique.
\end{remark}

% ============================================================
% 4.5 COMPUTATIONAL INTERPRETATION
% ============================================================

\subsection{A Connection to Numerical Eigenvalue Computation}
\label{subsec:courant_numerical}

In scientific computing, it is often impractical to
work with every vector in $\R^n$, especially when
$n$ is very large. Instead, numerical methods may
search within a lower-dimensional trial subspace.

Suppose $Z\in\R^{n\times m}$ has orthonormal columns,
so that $Z^TZ=I_m$, and let
\[
\mathcal K=\range(Z).
\]
Every nonzero $x\in\mathcal K$ can be written as
$x=Zy$ with $y\neq0$. Its Rayleigh quotient is
\[
\begin{aligned}
R_A(Zy)
&=\frac{(Zy)^TA(Zy)}{(Zy)^T(Zy)}\\
&=\frac{y^T(Z^TAZ)y}{y^Ty}.
\end{aligned}
\]

Thus optimizing the Rayleigh quotient on the trial
subspace is equivalent to optimizing the Rayleigh
quotient of the smaller symmetric matrix
\[
\boxed{H=Z^TAZ\in\R^{m\times m}.}
\]
The eigenvalues of $H$ are called \emph{Ritz values}
and are fundamental to the Rayleigh--Ritz procedure.
Their usefulness depends on how well the chosen trial
subspace captures the eigendirections of interest.

There is also an immediate connection to perturbation
theory. If $A$ is replaced by $A+E$, then for every
nonzero vector $x$,
\[
R_{A+E}(x)-R_A(x)=\frac{x^TEx}{x^Tx}.
\]
When $E$ is symmetric, its spectral norm gives
\[
\abs{R_{A+E}(x)-R_A(x)}\leq\norm{E}_2.
\]
The Courant--Fischer characterizations allow this
pointwise estimate to be transferred to the ordered
eigenvalues. This observation is the starting point
for Weyl's inequality, which we study next.
\end{comment}

% ============================================================
% SECTION 4: THE COURANT--FISCHER THEOREM (REVISED)
% ============================================================

\section{The Courant--Fischer Theorem}
\label{sec:courant_fischer}

In the preceding sections, we studied the Rayleigh quotient
and the geometry of linear subspaces. We now connect these
ideas through the Courant--Fischer theorem.

Our objective is to answer a natural question: we already
know how to characterize the smallest and largest eigenvalues
by minimizing or maximizing the Rayleigh quotient over all
nonzero vectors. Can we characterize the eigenvalues that
lie between these two extremes?

The answer is yes. The main difficulty is not the algebra;
it is understanding why we optimize twice: first over
vectors within a subspace and then over the subspaces
themselves. We develop this idea gradually.

% ------------------------------------------------------------
% 4.1 FROM EXTREME TO INTERMEDIATE EIGENVALUES
% ------------------------------------------------------------
\subsection{From Extreme to Intermediate Eigenvalues}
\label{subsec:courant_motivation}

Let $A\in\R^{n\times n}$ be \emph{real and symmetric}.
By the spectral theorem, there is an orthonormal eigenbasis
$v_1,\ldots,v_n$ with
\[
Av_i=\lambda_i v_i,\qquad
\lambda_1\leq\lambda_2\leq\cdots\leq\lambda_n.
\]
We keep this ascending eigenvalue convention throughout.

Recall the Rayleigh quotient
\[
R_A(x)=\frac{x^TAx}{x^Tx},\qquad x\neq0.
\]
In \cref{subsec:rayleigh_background}, we proved
\[
\boxed{
\lambda_1=\min_{x\neq0}R_A(x),\qquad
\lambda_n=\max_{x\neq0}R_A(x).
}
\]

\paragraph{How do we find the second-smallest eigenvalue?}
If we minimize $R_A(x)$ over all nonzero vectors, we
obtain $\lambda_1$. To reach $\lambda_2$, a natural
idea is to forbid the direction $v_1$.

Require $x\perp v_1$. Since the eigenvectors form an
orthonormal basis, this forces the coefficient of $v_1$
in the expansion of $x$ to vanish:
\[
x=c_2v_2+\cdots+c_nv_n.
\]
Thus
\[
R_A(x)=
\frac{\lambda_2c_2^2+\cdots+\lambda_nc_n^2}
     {c_2^2+\cdots+c_n^2}\geq\lambda_2.
\]
Equality occurs at $x=v_2$, so
\[
\boxed{
\lambda_2=
\min_{\substack{x\neq0\\x\perp v_1}}R_A(x).
}
\]

The same reasoning works for any $k$. If we forbid all
of $v_1,\ldots,v_{k-1}$, then the remaining nonzero
vectors have the form
\[
x=\sum_{i=k}^n c_i v_i,
\]
and consequently
\[
\boxed{
\lambda_k=
\min_{\substack{x\neq0\\
x\perp v_1,\ldots,v_{k-1}}}R_A(x).
}
\]
This is already a useful variational characterization.
However, it refers to eigenvectors that we might not
know. Courant--Fischer will replace these fixed
constraints by an optimization over \emph{all}
subspaces of a specified dimension.

\paragraph{Two special subspaces.}
For later use, define
\[
\mathcal S_0=\Span\{v_1,\ldots,v_k\},
\qquad \dim\mathcal S_0=k,
\]
and
\[
\mathcal T_0=\Span\{v_k,\ldots,v_n\},
\qquad \dim\mathcal T_0=n-k+1.
\]

Every nonzero $x\in\mathcal S_0$ can be expressed as
$x=\sum_{i=1}^k c_i v_i$. Since $\lambda_i\leq\lambda_k$
for $i\leq k$,
\[
R_A(x)=\frac{\sum_{i=1}^k\lambda_i c_i^2}
             {\sum_{i=1}^k c_i^2}
\leq\lambda_k.
\]
Equality is attained at $v_k$. Therefore,
\[
\boxed{
\max_{\substack{x\in\mathcal S_0\\x\neq0}}R_A(x)
=\lambda_k.
}
\]

Similarly, every nonzero $x\in\mathcal T_0$ is a
linear combination of $v_k,\ldots,v_n$, whose
eigenvalues are all at least $\lambda_k$. Therefore,
\[
\boxed{
\min_{\substack{x\in\mathcal T_0\\x\neq0}}R_A(x)
=\lambda_k.
}
\]

These two observations give us the correct candidates
for the optimizations that follow. The crucial question
is whether \emph{some other} subspace could produce an
even smaller maximum or an even larger minimum.

% ------------------------------------------------------------
% 4.2 WHAT THE TWO OPTIMIZATIONS MEAN
% ------------------------------------------------------------
\subsection{Concept of the Minimum among the Maximums}
\label{subsec:courant_nested}

Fix an integer $k$ between $1$ and $n$.
For each $k$-dimensional subspace $\mathcal S$, define
\[
\boxed{
M(\mathcal S)=
\max_{\substack{x\in\mathcal S\\x\neq0}}R_A(x).
}
\]

There are \emph{two separate choices} in this definition
and in the optimization we will study.

\begin{enumerate}[label=\textbf{Step \arabic*.},leftmargin=*]
\item Choose one $k$-dimensional subspace $\mathcal S$.
      Keep that subspace fixed.
\item Inside that particular subspace, examine all nonzero
      vectors and find the \emph{largest} Rayleigh quotient.
      Call that number $M(\mathcal S)$.
\item Repeat the process for every possible $k$-dimensional
      subspace, and select the \emph{smallest} of the numbers
      $M(\mathcal S)$.
\end{enumerate}

In symbols, we seek
\[
\boxed{
\min_{\dim\mathcal S=k}M(\mathcal S)
=\min_{\dim\mathcal S=k}
\max_{\substack{x\in\mathcal S\\x\neq0}}R_A(x).
}
\]

The inner maximum is computed \emph{first}, with
$\mathcal S$ fixed. The outer minimum then compares
the resulting maximum values across subspaces.
The order matters: this is not the same as simply
minimizing $R_A(x)$ over all nonzero vectors.

\paragraph{An example before the proof.}
Consider
\[
A=\diag(1,3,7),\qquad n=3,\quad k=2.
\]
If $\mathcal S=\Span\{e_1,e_2\}$, then
\[
M(\mathcal S)=3.
\]
If instead $\mathcal S=\Span\{e_1,e_3\}$, then
\[
M(\mathcal S)=7.
\]
These are two different planes with two different
inner maxima. Courant--Fischer asserts that \emph{no}
two-dimensional subspace has an inner maximum below
$3$. Therefore the smallest possible maximum is $3$,
the second eigenvalue.

This assertion is what we must \emph{prove}; checking
a few example subspaces does not establish it.

% ------------------------------------------------------------
% 4.3 DIMENSION COUNTING
% ------------------------------------------------------------
\subsection{Intersection of Subspaces}
\label{subsec:courant_dimension}

The following elementary fact is the geometric key to
the theorem. We encountered the same principle when
comparing subspaces of different dimensions.

\begin{lemma}[Dimension-Counting Principle]
\label{lem:courant_dimension}
Let $\mathcal S,\mathcal T\subseteq\R^n$ be subspaces.
If
\[
\dim\mathcal S+\dim\mathcal T>n,
\]
then they must contain a common nonzero vector:
\[
\boxed{\mathcal S\cap\mathcal T\neq\{0\}.}
\]
\end{lemma}

\begin{proof}
The dimension formula is
\[
\dim(\mathcal S+\mathcal T)
=\dim\mathcal S+\dim\mathcal T
-\dim(\mathcal S\cap\mathcal T).
\]
But $\mathcal S+\mathcal T$ is a subspace of $\R^n$,
so its dimension cannot exceed $n$. Hence
\[
\dim(\mathcal S\cap\mathcal T)
\geq\dim\mathcal S+\dim\mathcal T-n>0.
\]
Therefore the intersection has positive dimension,
and a nonzero vector exists in it.
\end{proof}

The specific calculation that will appear in both
parts of the proof is
\[
\boxed{k+(n-k+1)=n+1>n.}
\]
A $k$-dimensional subspace and an
$(n-k+1)$-dimensional subspace cannot avoid
one another inside $\R^n$.


It's worth noting that the shared vector need not equal $v_k$, and its
Rayleigh quotient need not be exactly $\lambda_k$.
The dimension argument guarantees only that a
\emph{nonzero shared vector exists}. We then use
which special eigenspace it belongs to in order
to establish an \emph{inequality} for its
Rayleigh quotient. That inequality is sufficient
for the proof.


% ------------------------------------------------------------
% 4.4 STATEMENT OF THE THEOREM
% ------------------------------------------------------------
\subsection{Statement of the Courant--Fischer Theorem}
\label{subsec:courant_theorem}

\begin{theorem}[Courant--Fischer]
\label{thm:courant_fischer}
Let $A\in\R^{n\times n}$ be real and symmetric, with
eigenvalues ordered as
$\lambda_1\leq\cdots\leq\lambda_n$.
Then, for each $k=1,\ldots,n$,
\[
\boxed{
\lambda_k=
\min_{\substack{\mathcal S\subseteq\R^n\\
                 \dim\mathcal S=k}}
\ \max_{\substack{x\in\mathcal S\\x\neq0}}R_A(x).
}
\]
This is the \emph{min--max characterization}.

Equivalently,
\[
\boxed{
\lambda_k=
\max_{\substack{\mathcal T\subseteq\R^n\\
                 \dim\mathcal T=n-k+1}}
\ \min_{\substack{x\in\mathcal T\\x\neq0}}R_A(x).
}
\]
This is the \emph{max--min characterization}.
\end{theorem}

We prove the two formulas separately, beginning with
the min--max characterization. The proof repeatedly
uses a familiar strategy: \emph{to prove two numbers
are equal, prove an upper bound and a matching lower
bound}.

% ------------------------------------------------------------
% 4.5 MIN--MAX PROOF: TWO INEQUALITIES
% ------------------------------------------------------------
\subsection{Proof of the Min--Max Characterization}
\label{subsec:courant_minmax_proof}

\begin{proof}[Proof of the min--max formula]
Let
\[
\boxed{
\alpha_k=
\min_{\dim\mathcal S=k}
\max_{\substack{x\in\mathcal S\\x\neq0}}R_A(x).
}
\]
Our goal is to show $\alpha_k=\lambda_k$.
We will establish
\[
\alpha_k\leq\lambda_k
\qquad\text{and}\qquad
\alpha_k\geq\lambda_k.
\]

\medskip
\noindent\textbf{Part A: Proving $\alpha_k\leq\lambda_k$.}\quad
This is the easier direction. Since $\alpha_k$ is
a \emph{minimum over all $k$-dimensional subspaces},
it suffices to exhibit just \emph{one} admissible
subspace whose inner maximum equals $\lambda_k$.

Choose the subspace spanned by the first $k$
eigenvectors:
\[
\boxed{\mathcal S_0=\Span\{v_1,\ldots,v_k\}.}
\]
Its dimension is $k$. Every nonzero $x\in\mathcal S_0$
has the form $x=\sum_{i=1}^k c_i v_i$, so
\[
\begin{aligned}
R_A(x)
&=\frac{\sum_{i=1}^k\lambda_i c_i^2}
         {\sum_{i=1}^k c_i^2}\\[3pt]
&\leq\frac{\sum_{i=1}^k\lambda_k c_i^2}
          {\sum_{i=1}^k c_i^2}\\[3pt]
&=\lambda_k.
\end{aligned}
\]
The inequality follows from $\lambda_i\leq\lambda_k$
and $c_i^2\geq0$. Choosing $x=v_k$ gives equality.
Thus
\[
\boxed{
\max_{\substack{x\in\mathcal S_0\\x\neq0}}
R_A(x)=\lambda_k.
}
\]

Now recall what $\alpha_k$ means: we choose the
\emph{smallest} inner maximum among \emph{all}
$k$-dimensional subspaces. The smallest of a set of
numbers cannot exceed a particular number in that
set. Therefore,
\[
\boxed{\alpha_k\leq\lambda_k.}
\]

\medskip
\noindent\textbf{Part B: Proving $\alpha_k\geq\lambda_k$.}\quad
This direction is more important. We must rule out
the possibility that some clever choice of a
$k$-dimensional subspace has an inner maximum
\emph{smaller} than $\lambda_k$.

Take an \emph{arbitrary} $k$-dimensional subspace
$\mathcal S$. We make no assumption about its basis.
In particular, $\mathcal S$ need not be spanned by
eigenvectors.

Construct the \emph{fixed} subspace
\[
\boxed{\mathcal T_0=\Span\{v_k,v_{k+1},\ldots,v_n\}.}
\]
Why this particular subspace? Every nonzero vector
in $\mathcal T_0$ has a Rayleigh quotient at least
$\lambda_k$, since it uses only eigenvectors
associated with $\lambda_k,\ldots,\lambda_n$.

The dimensions are
\[
\dim\mathcal S=k,
\qquad
\dim\mathcal T_0=n-k+1.
\]
Hence
\[
\dim\mathcal S+\dim\mathcal T_0=n+1>n.
\]
By \cref{lem:courant_dimension}, the two subspaces
must have a nonzero intersection. Choose
\[
\boxed{0\neq w\in\mathcal S\cap\mathcal T_0.}
\]

Because $w\in\mathcal T_0$, it has an expansion
\[
w=\sum_{i=k}^n c_i v_i,
\]
with at least one $c_i\neq0$. Therefore,
\[
\begin{aligned}
R_A(w)
&=\frac{\sum_{i=k}^n\lambda_i c_i^2}
         {\sum_{i=k}^n c_i^2}\\[3pt]
&\geq\frac{\sum_{i=k}^n\lambda_k c_i^2}
          {\sum_{i=k}^n c_i^2}\\[3pt]
&=\lambda_k.
\end{aligned}
\]
This is our first crucial inequality:
\[
\boxed{R_A(w)\geq\lambda_k.}
\]

But remember that $w$ also lies in the \emph{arbitrary}
subspace $\mathcal S$. If we maximize the Rayleigh
quotient over \emph{every} nonzero vector in
$\mathcal S$, the answer must be at least the value
at the particular vector $w$. Thus
\[
\boxed{
\max_{\substack{x\in\mathcal S\\x\neq0}}R_A(x)
\geq R_A(w)\geq\lambda_k.
}
\]

This reasoning works for \emph{every}
$k$-dimensional subspace $\mathcal S$. Therefore,
none of their inner maxima is below $\lambda_k$.
Even after taking the smallest of all those maxima,
we still have
\[
\boxed{\alpha_k\geq\lambda_k.}
\]

\medskip
\noindent\textbf{Conclusion.}\quad
We have proved both
\[
\alpha_k\leq\lambda_k
\qquad\text{and}\qquad
\alpha_k\geq\lambda_k.
\]
Consequently, $\alpha_k=\lambda_k$, or
\[
\boxed{
\lambda_k=
\min_{\dim\mathcal S=k}
\max_{\substack{x\in\mathcal S\\x\neq0}}R_A(x).
}
\]
\end{proof}

\paragraph{The idea in one sentence.}
We found one $k$-dimensional subspace whose maximum
is $\lambda_k$, and we proved that every other
$k$-dimensional subspace has a maximum at least
$\lambda_k$. Therefore the smallest possible
maximum is exactly $\lambda_k$.


We again want to emphasize that the vector $w\in\mathcal S\cap\mathcal T_0$ used in
Part B was \emph{not} assumed to equal $v_k$.
Indeed, $R_A(w)$ may be strictly greater than
$\lambda_k$. We need only the inequality
$R_A(w)\geq\lambda_k$ to establish the lower bound.
Confusing the special subspace $\mathcal S_0$ from
Part A with the arbitrary subspace $\mathcal S$
from Part B obscures the logic of the proof.


% ------------------------------------------------------------
% 4.6 MAX--MIN PROOF
% ------------------------------------------------------------
\subsection{Proof of the Max--Min Characterization}
\label{subsec:courant_maxmin_proof}

Now we reverse the roles of the inner and outer
optimizations. Let
\[
\boxed{
\beta_k=
\max_{\dim\mathcal T=n-k+1}
\min_{\substack{x\in\mathcal T\\x\neq0}}R_A(x).
}
\]

This time we first choose a subspace $\mathcal T$ of
dimension $n-k+1$, find the \emph{smallest} Rayleigh
quotient inside it, and finally choose the
\emph{largest} of those minimum values.
We must prove $\beta_k=\lambda_k$.

\begin{proof}[Proof of the max--min formula]
\medskip
\noindent\textbf{Part A: Proving $\beta_k\geq\lambda_k$.}\quad
Choose
\[
\mathcal T_0=\Span\{v_k,\ldots,v_n\}.
\]
We already established that every nonzero vector in
this subspace has Rayleigh quotient at least
$\lambda_k$, with equality at $v_k$. Hence
\[
\min_{\substack{x\in\mathcal T_0\\x\neq0}}R_A(x)
=\lambda_k.
\]
Since $\beta_k$ is the \emph{largest} inner minimum
among all admissible subspaces, it cannot be smaller
than the minimum for this particular subspace:
\[
\boxed{\beta_k\geq\lambda_k.}
\]

\medskip
\noindent\textbf{Part B: Proving $\beta_k\leq\lambda_k$.}\quad
Let $\mathcal T$ be an \emph{arbitrary} subspace
with $\dim\mathcal T=n-k+1$. Introduce the fixed
$k$-dimensional subspace
\[
\mathcal S_0=\Span\{v_1,\ldots,v_k\}.
\]
Every nonzero vector of $\mathcal S_0$ has Rayleigh
quotient at most $\lambda_k$. Moreover,
\[
\dim\mathcal T+\dim\mathcal S_0
=(n-k+1)+k=n+1>n.
\]
The dimension-counting lemma guarantees a nonzero
vector $w\in\mathcal T\cap\mathcal S_0$.
Because $w\in\mathcal S_0$,
\[
R_A(w)
=\frac{\sum_{i=1}^k\lambda_i c_i^2}
       {\sum_{i=1}^k c_i^2}
\leq\lambda_k.
\]
Because $w\in\mathcal T$ as well, the \emph{minimum}
over $\mathcal T$ cannot exceed the value at $w$:
\[
\boxed{
\min_{\substack{x\in\mathcal T\\x\neq0}}R_A(x)
\leq R_A(w)\leq\lambda_k.
}
\]
This holds for every admissible $\mathcal T$, so
\emph{none} of their inner minima can be greater than
$\lambda_k$. Taking the largest among them still
cannot produce a number above $\lambda_k$:
\[
\boxed{\beta_k\leq\lambda_k.}
\]

\medskip
\noindent\textbf{Conclusion.}\quad
Combining $\beta_k\geq\lambda_k$ and
$\beta_k\leq\lambda_k$, we obtain
\[
\boxed{
\lambda_k=
\max_{\dim\mathcal T=n-k+1}
\min_{\substack{x\in\mathcal T\\x\neq0}}R_A(x).
}
\]
\end{proof}

The two proofs have exactly the same geometric core,
but their inequality directions are reversed. In the
min--max proof, the shared vector produces a quotient
\emph{at least} $\lambda_k$. In the max--min proof,
it produces a quotient \emph{at most} $\lambda_k$.

% ------------------------------------------------------------
% 4.7 WORKED EXAMPLE
% ------------------------------------------------------------
\subsection{A Worked Example in $\R^3$}
\label{subsec:courant_example}

Let
\[
A=\begin{bmatrix}1&0&0\\0&3&0\\0&0&7\end{bmatrix},
\quad
\lambda_1=1,\quad\lambda_2=3,\quad\lambda_3=7.
\]
We recover $\lambda_2=3$ using two-dimensional
subspaces ($k=2$ and $n-k+1=2$).

\paragraph{Min--max.}
On $\mathcal S_0=\Span\{e_1,e_2\}$,
\[
R_A(ae_1+be_2)=\frac{a^2+3b^2}{a^2+b^2}\leq3
\]
for nonzero $(a,b)$, and equality occurs at $e_2$.
Therefore the maximum on $\mathcal S_0$ is $3$.

If $\mathcal S$ is \emph{any} two-dimensional plane,
it intersects $\Span\{e_2,e_3\}$ in a nonzero vector
$w=be_2+ce_3$. At that vector,
\[
R_A(w)=\frac{3b^2+7c^2}{b^2+c^2}\geq3.
\]
Thus every plane has a maximum at least $3$,
and the special plane achieves exactly $3$:
\[
\boxed{
\min_{\dim\mathcal S=2}
\max_{\substack{x\in\mathcal S\\x\neq0}}R_A(x)=3.
}
\]

\paragraph{Max--min.}
On $\mathcal T_0=\Span\{e_2,e_3\}$, the smallest
Rayleigh quotient is $3$. Any other two-dimensional
plane $\mathcal T$ intersects $\Span\{e_1,e_2\}$
in a nonzero vector, whose Rayleigh quotient is at
most $3$. Thus no plane can have a minimum exceeding
$3$, and
\[
\boxed{
\max_{\dim\mathcal T=2}
\min_{\substack{x\in\mathcal T\\x\neq0}}R_A(x)=3.
}
\]

\paragraph{A useful caution.}
For example, take the arbitrary plane
\[
\mathcal S=\Span\{e_1+e_3,\ e_1+e_2\}.
\]
It contains $w=e_3-e_2$, which also belongs to
$\Span\{e_2,e_3\}$. Here
\[
R_A(w)=\frac{7+3}{2}=5>\lambda_2.
\]
The shared vector does \emph{not} have to give
exactly $3$: it only needs to give a value
\emph{at least} $3$ to force the maximum over
$\mathcal S$ to be at least $3$.

\begin{remark}[Repeated Eigenvalues]
Neither formulation requires distinct eigenvalues.
The spectral theorem still supplies an orthonormal
eigenbasis, and the dimension argument remains valid
when eigenvalues repeat. Minimizing or maximizing
subspaces may then be nonunique.
\end{remark}

% ------------------------------------------------------------
% 4.8 COMPUTATIONAL CONNECTION
% ------------------------------------------------------------
\subsection{Interpretation for Scientific Computing}
\label{subsec:courant_numerical}

The Courant--Fischer theorem is a variational
\emph{characterization}, not by itself an efficient
algorithm for enumerating all subspaces. Nevertheless,
it motivates practical subspace-based eigenvalue methods.

Let $Z\in\R^{n\times m}$ have orthonormal columns and
let $\mathcal K=\range(Z)$. For $x=Zy\neq0$,
\[
\begin{aligned}
R_A(Zy)
&=\frac{y^TZ^TAZy}{y^TZ^TZy}\\
&=\frac{y^T(Z^TAZ)y}{y^Ty}.
\end{aligned}
\]
Thus the Rayleigh quotient restricted to
$\mathcal K$ is the Rayleigh quotient of the smaller
symmetric matrix
\[
\boxed{H=Z^TAZ.}
\]
Its eigenvalues are called \emph{Ritz values}.
This observation underlies the Rayleigh--Ritz method.
The quality of the approximation depends on the
choice of the trial subspace $\mathcal K$.



However, if the reader is still having difficulty
understanding the Courant--Fischer theorem and its
proof, there is no reason to be discouraged.
The theorem is conceptually challenging, and its
geometric argument is not always easy to appreciate
on a first reading. Indeed, some textbooks state
the result without providing a detailed proof,
while others pass rather quickly through the
dimension-counting argument. We have tried our
best to explain the underlying mathematical ideas
as clearly and constructively as possible.
Nevertheless, if the argument remains unclear,
we encourage the reader to revisit this section
several times, paying particular attention to
the geometry of intersecting subspaces and the
meaning of the two successive optimization
operations. Sometimes, visualizing what a theorem
is trying to accomplish is more illuminating
than repeatedly manipulating its equations.

It is also worth remembering that the ideas
of minimax and maximin optimization extend far
beyond numerical linear algebra. Related
optimization and selection principles occur
frequently in computer science, algorithm design,
and technical interviews. Let us conclude this
section with a simple problem that may help
make the underlying idea more intuitive.

Suppose an interviewer asks the following question:
\emph{Given a database containing every person
in the United States, along with their estimated
net worth, how would you identify the ten richest
people and report their net worth?}

A natural first response might be to sort the
entire database in descending order of net worth
and return the first ten entries. Although this
approach is mathematically correct, it may be
computationally wasteful, especially when the
database contains millions of records. Sorting
the entire collection requires considerably
more work than is necessary when we only need
the ten largest values. In a real database,
an appropriate index might offer an even
more efficient solution.

Instead, consider a simple one-pass selection
algorithm. We begin by maintaining a collection
of ten candidates, initially populated from
the first ten records. Among these candidates,
we keep track of the person with the
\emph{smallest} net worth.

Now we traverse the remaining records.
Whenever we encounter a person whose net
worth exceeds that of the poorest person
currently in our collection, we replace
that candidate with the new person.
Otherwise, we leave our collection unchanged.

At the end of the traversal, our collection
contains the ten richest people in the
database, and the person with the smallest
net worth among these ten is precisely
the tenth-richest person overall.

This strategy can be implemented efficiently
using a min-heap containing only ten entries.
For a database with $N$ records, maintaining
a heap of size $k=10$ requires
$O(N\log k)$ time and $O(k)$ additional space,
rather than sorting all $N$ records.

Notice the interesting optimization principle
behind this approach. We want to select a
collection of ten people such that the
\emph{minimum} net worth within the selected
collection is as large as possible.

Mathematically, if $\mathcal P$ denotes the
set of people and $w(p)$ denotes the net
worth of person $p$, we are considering
\[
\boxed{
\max_{\substack{\mathcal S\subseteq\mathcal P\\
|\mathcal S|=10}}
\quad
\min_{p\in\mathcal S}w(p).
}
\]

The result is exactly the tenth-largest
net worth in the database.

This is a discrete maximin selection problem,
not an application of the Courant--Fischer
theorem itself. Nevertheless, the resemblance
in the nested optimization structure is
worth appreciating.

In the Courant--Fischer theorem, we optimize
the Rayleigh quotient over subspaces of
prescribed dimensions. In our database
example, we optimize net worth over
collections containing a prescribed
number of people.

The mathematical objects are entirely
different, but the underlying question
is remarkably similar:

\begin{quote}
\emph{
How should we choose a collection of
a fixed size so that its worst-case
value is as large as possible?
}
\end{quote}

Thus, whenever we encounter a maximin
or minimax expression, we should first
identify what is being selected, what
quantity is being optimized within that
selection, and what the outer optimization
is trying to accomplish.

This simple habit can make even a
sophisticated theorem such as
Courant--Fischer considerably easier
to understand.


In the next section, we shall use Courant--Fischer
to answer a different question: if a symmetric
matrix is perturbed by another symmetric matrix,
how much can any individual ordered eigenvalue change?
That leads to Weyl's inequality.

% ============================================================
% SECTION 5: WEYL'S INEQUALITY
% ============================================================

% ============================================================
% SECTION 5: WEYL'S EIGENVALUE PERTURBATION INEQUALITY
% ============================================================

\section{Weyl's Eigenvalue Perturbation Inequality}
\label{sec:weyl}

In the previous section, we learned that the eigenvalues
of a real symmetric matrix can be characterized by
optimizing its Rayleigh quotient over subspaces of fixed
dimensions. We now use this variational point of view to
answer one of the most important questions in scientific
computing: \emph{how much can an eigenvalue change when
the entries of a matrix are perturbed?}

The remarkable conclusion is that, for symmetric matrices,
every ordered eigenvalue is stable under a small symmetric
perturbation. The proof will be considerably shorter than
the Courant--Fischer proof, because almost all the difficult
geometric work has already been done.

% ------------------------------------------------------------
% 5.1 SETTING UP THE PROBLEM
% ------------------------------------------------------------

\subsection{The Perturbation Problem}
\label{subsec:weyl_setup}

Let
\[
A\in\R^{n\times n},\qquad A=A^T,
\]
and suppose that the matrix available to us is
\[
\boxed{\widetilde A=A+E,}
\]
where $E\in\R^{n\times n}$ represents a perturbation.
We assume $E=E^T$, so that $\widetilde A$ is also
real and symmetric.

The entries of $E$ may represent rounding errors,
measurement uncertainty, or errors introduced by an
approximation. We shall measure the size of this
perturbation using its spectral norm $\norm{E}_2$.

Write the ordered eigenvalues of the two matrices as
\[
\begin{aligned}
\lambda_1(A)&\leq\cdots\leq\lambda_n(A),\\
\lambda_1(\widetilde A)&\leq\cdots\leq
\lambda_n(\widetilde A).
\end{aligned}
\]

Our objective is to compare $\lambda_k(A)$ and
$\lambda_k(\widetilde A)$ for the \emph{same index} $k$.
Notice that we are comparing ordered eigenvalues, not
attempting to match individual eigenvectors of $A$ with
those of $\widetilde A$.

% ------------------------------------------------------------
% 5.2 FIRST STEP: ONE VECTOR
% ------------------------------------------------------------

\subsection{How Much Can One Rayleigh Quotient Change?}
\label{subsec:weyl_rayleigh}

Before comparing eigenvalues, let us compare the
Rayleigh quotients of $A$ and $\widetilde A$ at one
fixed nonzero vector $x$.

By definition,
\[
\begin{aligned}
R_{\widetilde A}(x)
&=\frac{x^T(A+E)x}{x^Tx}\\[3pt]
&=\frac{x^TAx}{x^Tx}+\frac{x^TEx}{x^Tx}\\[3pt]
&=R_A(x)+R_E(x).
\end{aligned}
\]
Thus,
\[
\boxed{R_{\widetilde A}(x)-R_A(x)=R_E(x).}
\]

This identity is the starting point of the proof.
It says that the change in the Rayleigh quotient at a
fixed vector is itself the Rayleigh quotient of the
perturbation matrix.

Because $E$ is symmetric, the bounds established in
\cref{subsec:rayleigh_background} apply to $R_E$.
If $\lambda_1(E)\leq\cdots\leq\lambda_n(E)$ are the
eigenvalues of $E$, then
\[
\lambda_1(E)\leq R_E(x)\leq\lambda_n(E).
\]
Combining the two formulas gives
\[
\boxed{
R_A(x)+\lambda_1(E)
\leq R_{\widetilde A}(x)
\leq R_A(x)+\lambda_n(E).
}
\]
This inequality holds for \emph{every} nonzero $x$.

Recall also that, for a real symmetric matrix,
\[
\norm{E}_2=\max_i\abs{\lambda_i(E)}.
\]
Consequently,
\[
-\norm{E}_2\leq\lambda_1(E)
\leq\lambda_n(E)\leq\norm{E}_2.
\]
We obtain the simpler two-sided estimate
\[
\boxed{
R_A(x)-\norm{E}_2
\leq R_{\widetilde A}(x)
\leq R_A(x)+\norm{E}_2.
}
\]

\paragraph{The remaining question.}
We have bounded the change in the Rayleigh quotient
at every vector. But an eigenvalue is not defined by
a single arbitrary Rayleigh quotient: Courant--Fischer
characterizes it through a maximum within a subspace,
followed by a minimum over subspaces.

How do we transfer the bound through these two
optimization operations? Fortunately, the necessary
principle is elementary.

% ------------------------------------------------------------
% 5.3 THE LOGIC OF OPTIMIZATION
% ------------------------------------------------------------

\subsection{Why the Bound Survives Minimization and Maximization}
\label{subsec:weyl_optimization}

Suppose two real-valued functions $f$ and $g$ satisfy
\[
f(x)\leq g(x)+c
\qquad\text{for every }x\text{ in a set }X,
\]
where $c$ is a constant independent of $x$.
Then
\[
\boxed{
\max_{x\in X} f(x)
\leq \max_{x\in X}g(x)+c,
}
\]
provided the maxima exist. Indeed, no value of $f$
can exceed the largest value of $g$ plus $c$.

Similarly,
\[
\boxed{
\min_{x\in X} f(x)
\leq \min_{x\in X}g(x)+c.
}
\]
To see this second statement, choose a point $x_*$
at which $g$ attains its minimum. Then
\[
\min_{x\in X}f(x)
\leq f(x_*)\leq g(x_*)+c
=\min_{x\in X}g(x)+c.
\]

In particular, if the same inequality holds for all
vectors in \emph{every} admissible subspace, we may
first take the maximum over vectors in each subspace
and then take the minimum over the subspaces. Both
operations preserve the inequality.

This observation explains why the Courant--Fischer
theorem is so useful in perturbation theory.

% ------------------------------------------------------------
% 5.4 MAIN THEOREM AND PROOF
% ------------------------------------------------------------

\subsection{Weyl's Inequality}
\label{subsec:weyl_theorem}

\begin{theorem}[Weyl's Eigenvalue Perturbation Bound]
\label{thm:weyl}
Let $A,E\in\R^{n\times n}$ be real symmetric matrices,
and let $\widetilde A=A+E$. Order the eigenvalues of
each symmetric matrix from smallest to largest.
Then, for every $k=1,\ldots,n$,
\[
\boxed{
\lambda_k(A)+\lambda_1(E)
\leq \lambda_k(A+E)
\leq \lambda_k(A)+\lambda_n(E).
}
\]
In particular,
\[
\boxed{
\abs{\lambda_k(A+E)-\lambda_k(A)}
\leq\norm{E}_2.
}
\]
\end{theorem}

\begin{proof}
We use the min--max characterization from
\cref{thm:courant_fischer}:
\[
\lambda_k(A)
=\min_{\dim\mathcal S=k}
  \max_{\substack{x\in\mathcal S\\x\neq0}}R_A(x).
\]
The same formula holds for $\widetilde A=A+E$,
because it is also symmetric.

\medskip
\noindent\textbf{Step 1: Establish the upper bound.}

Fix any $k$-dimensional subspace $\mathcal S$.
For every nonzero vector $x\in\mathcal S$, we already
proved that
\[
R_{\widetilde A}(x)
\leq R_A(x)+\lambda_n(E).
\]
Taking the maximum over the nonzero vectors of
this \emph{same} subspace gives
\[
\begin{aligned}
\max_{\substack{x\in\mathcal S\\x\neq0}}
R_{\widetilde A}(x)
&\leq
\max_{\substack{x\in\mathcal S\\x\neq0}}
\bigl(R_A(x)+\lambda_n(E)\bigr)\\
&=\max_{\substack{x\in\mathcal S\\x\neq0}}
R_A(x)+\lambda_n(E).
\end{aligned}
\]




Now take the minimum over \emph{all} $k$-dimensional
subspaces. Since the inequality holds for each one,
\[
\begin{aligned}
&\min_{\dim\mathcal S=k}
\max_{\substack{x\in\mathcal S\\x\neq0}}
R_{\widetilde A}(x)\\
&\quad\leq
\min_{\dim\mathcal S=k}
\left(
\max_{\substack{x\in\mathcal S\\x\neq0}}
R_A(x)+\lambda_n(E)
\right)\\
&\quad=
\min_{\dim\mathcal S=k}
\max_{\substack{x\in\mathcal S\\x\neq0}}
R_A(x)+\lambda_n(E).
\end{aligned}
\]

By Courant--Fischer, the left side is
$\lambda_k(A+E)$, while the first term on the
right is $\lambda_k(A)$. Thus,
\[
\boxed{
\lambda_k(A+E)
\leq\lambda_k(A)+\lambda_n(E).
}
\]

\medskip
\noindent\textbf{Step 2: Establish the lower bound.}

Again fix an arbitrary $k$-dimensional subspace
$\mathcal S$. The pointwise lower bound is
\[
R_{\widetilde A}(x)
\geq R_A(x)+\lambda_1(E)
\qquad(x\in\mathcal S,\ x\neq0).
\]
Taking the maximum on both sides, we get
\[
\max_{\substack{x\in\mathcal S\\x\neq0}}
R_{\widetilde A}(x)
\geq
\max_{\substack{x\in\mathcal S\\x\neq0}}
R_A(x)+\lambda_1(E).
\]
Since this holds for every $k$-dimensional
subspace, taking the minimum over all such
subspaces preserves the inequality:
\[
\begin{aligned}
&\min_{\dim\mathcal S=k}
\max_{\substack{x\in\mathcal S\\x\neq0}}
R_{\widetilde A}(x)\\
&\quad\geq
\min_{\dim\mathcal S=k}
\max_{\substack{x\in\mathcal S\\x\neq0}}
R_A(x)+\lambda_1(E).
\end{aligned}
\]
Applying Courant--Fischer once more gives
\[
\boxed{
\lambda_k(A+E)
\geq\lambda_k(A)+\lambda_1(E).
}
\]

\medskip
\noindent\textbf{Step 3: Combine the bounds.}

We have proved
\[
\lambda_1(E)
\leq\lambda_k(A+E)-\lambda_k(A)
\leq\lambda_n(E).
\]
Because
\[
-\norm{E}_2\leq\lambda_1(E),
\qquad
\lambda_n(E)\leq\norm{E}_2,
\]
it follows that
\[
-\norm{E}_2
\leq\lambda_k(A+E)-\lambda_k(A)
\leq\norm{E}_2.
\]
Equivalently,
\[
\boxed{
\abs{\lambda_k(A+E)-\lambda_k(A)}
\leq\norm{E}_2.
}
\]
This completes the proof.
\end{proof}

\paragraph{What made the proof work?}
The argument rests on only two ingredients:
for every vector, the Rayleigh quotient changes
by at most the spectral norm of $E$; and the
Courant--Fischer min--max formula preserves
this uniform bound when we optimize over vectors
and subspaces. We did not need to compute any
eigenvector of the perturbed matrix.

% ------------------------------------------------------------
% 5.5 EXAMPLES AND CONSEQUENCES
% ------------------------------------------------------------

\subsection{Examples and Consequences}
\label{subsec:weyl_examples}

\begin{example}[A Perturbation That Mixes Eigenvectors]
\label{ex:weyl_two_by_two}
Consider
\[
A=\begin{bmatrix}1&0\\0&3\end{bmatrix},
\qquad
E=\begin{bmatrix}0&\tfrac12\\
                  \tfrac12&0\end{bmatrix}.
\]
Both matrices are symmetric, and the eigenvalues
of $A$ are $1$ and $3$. The eigenvalues of $E$
are $-\tfrac12$ and $\tfrac12$, so
\[
\boxed{\norm{E}_2=\tfrac12.}
\]
The perturbed matrix is
\[
A+E=\begin{bmatrix}1&\tfrac12\\
                   \tfrac12&3\end{bmatrix}.
\]
Its characteristic polynomial is
\[
\begin{aligned}
\det((A+E)-\lambda I)
&=(1-\lambda)(3-\lambda)-\tfrac14\\
&=\lambda^2-4\lambda+\tfrac{11}{4}.
\end{aligned}
\]
Thus its ordered eigenvalues are
\[
\lambda_1(A+E)=2-\frac{\sqrt5}{2},
\qquad
\lambda_2(A+E)=2+\frac{\sqrt5}{2}.
\]
The changes are equal in magnitude:
\[
\begin{aligned}
\abs{\lambda_1(A+E)-1}
&=\frac{\sqrt5}{2}-1,\\
\abs{\lambda_2(A+E)-3}
&=\frac{\sqrt5}{2}-1.
\end{aligned}
\]
Numerically, this is approximately $0.118$, which
is smaller than $\norm{E}_2=0.5$, as Weyl's
inequality predicts.

Notice that $A+E$ is not diagonal: its eigenvectors
are rotated relative to the coordinate axes.
The eigenvalues remain close even though the
eigendirections have changed.
\end{example}

\begin{remark}[The Bound Is Sharp]
\label{rem:weyl_sharp}
A bound is called \emph{sharp} if equality can occur.
Let $E=tI$ for any real number $t$. Then
\[
(A+tI)v_i=(\lambda_i(A)+t)v_i,
\]
so
\[
\lambda_k(A+tI)-\lambda_k(A)=t.
\]
Since $\norm{tI}_2=\abs{t}$,
\[
\boxed{
\abs{\lambda_k(A+tI)-\lambda_k(A)}
=\norm{tI}_2.
}
\]
Thus the constant $1$ in Weyl's norm bound
cannot be replaced by a smaller universal constant.
\end{remark}

\begin{corollary}[Monotonicity Under Positive Semidefinite Perturbations]
\label{cor:weyl_psd}
If $E$ is positive semidefinite, then for every $k$,
\[
\boxed{\lambda_k(A+E)\geq\lambda_k(A).}
\]
If $E$ is negative semidefinite, the inequality
is reversed.
\end{corollary}
\begin{proof}
When $E$ is positive semidefinite,
$\lambda_1(E)\geq0$. The lower Weyl bound gives
\[
\lambda_k(A+E)\geq\lambda_k(A)+\lambda_1(E)
\geq\lambda_k(A).
\]
For negative semidefinite $E$, use
$\lambda_n(E)\leq0$ in the upper Weyl bound.
\end{proof}

% ------------------------------------------------------------
% 5.6 EIGENVALUES VERSUS EIGENVECTORS
% ------------------------------------------------------------

\subsection{What Weyl's Inequality Does Not Tell Us}
\label{subsec:weyl_limitations}

Weyl's inequality is a statement about the ordered
eigenvalues of symmetric matrices. It does not say
that every individual eigenvector must change by a
similarly small amount. Even for symmetric matrices,
eigenvectors may be sensitive when the corresponding
eigenvalues are close to one another.

It is also important not to apply the theorem in
this form to arbitrary nonsymmetric matrices.
The proof relied on the spectral theorem and
Courant--Fischer, which require symmetry.
For general matrices, eigenvalues can be much more
sensitive to perturbations.

A useful distinction is therefore:
\begin{quote}
\emph{Weyl's inequality controls how far the ordered
eigenvalues move. It does not, by itself, control
how far the associated invariant subspaces rotate.}
\end{quote}

To study the second question, we need additional
information about the separation between eigenvalues.
In the next section, the Davis--Kahan theorem will
connect the size of a perturbation, the spectral gap,
and the subspace distance developed in
\cref{sec:subspace_geometry}.

% ============================================================
% SECTION 6: DAVIS--KAHAN THEOREM
% ============================================================
% ============================================================
% SECTION 6: DAVIS--KAHAN SIN THETA THEOREM
% ============================================================

\section{The Davis--Kahan \texorpdfstring{$\sin\Theta$}{sin Theta} Theorem}
\label{sec:davis_kahan}

Weyl's inequality guarantees that the ordered eigenvalues
of a real symmetric matrix are stable under a small symmetric
perturbation. But what about their eigenvectors? More
precisely, if $\widetilde A=A+E$, must an eigenvector of
$\widetilde A$ point approximately in the same direction
as an eigenvector of $A$?

There is an important obstacle: even normalized eigenvectors
are not uniquely determined, and repeated eigenvalues can
have many different eigenvectors. The appropriate objects
to compare are therefore the \emph{subspaces} generated
by the eigenvectors, using the distance developed in
\cref{sec:subspace_geometry}.

The Davis--Kahan theorem explains when these subspaces
remain close. Its central message is that the stability
of an eigenspace depends on two quantities: the size
of the perturbation and the separation between the
relevant eigenvalues and the rest of the spectrum.

% ------------------------------------------------------------
% 6.1 MOTIVATION
% ------------------------------------------------------------

% ============================================================
% REPLACEMENT FOR SUBSECTION: DK MOTIVATION
% ============================================================
\subsection{Eigenvalues Can Be Stable While Eigenvectors Rotate}
\label{subsec:dk_motivation}

Weyl's inequality tells us how much the eigenvalues of a symmetric
matrix can move under a perturbation. It does not, however, directly
tell us how much the corresponding eigenvectors can rotate. The
following two-dimensional example makes this distinction explicit.


\paragraph{The original matrix and its perturbation.}
Let
\[
A=\begin{bmatrix}0&0\\0&\gamma\end{bmatrix},
\qquad
E=\begin{bmatrix}0&\varepsilon\\\varepsilon&0\end{bmatrix},
\qquad \gamma>0,\quad\varepsilon>0.
\]
Because $A$ is diagonal, its eigenpairs are
\[
\lambda_1(A)=0,\quad v_1=e_1=\begin{bmatrix}1\\0\end{bmatrix},
\qquad
\lambda_2(A)=\gamma,\quad v_2=e_2=\begin{bmatrix}0\\1\end{bmatrix}.
\]
The distance between the two eigenvalues is $\gamma$;
this is the \emph{spectral gap} in our example.

To calculate the size of the perturbation, observe that
\[
\det(E-\mu I)
=\det\begin{bmatrix}-\mu&\varepsilon\\
                   \varepsilon&-\mu\end{bmatrix}
=\mu^2-\varepsilon^2.
\]
Thus the eigenvalues of $E$ are $\pm\varepsilon$.
Since $E$ is symmetric, its spectral norm equals the
largest absolute eigenvalue, so
\[
\boxed{\norm{E}_2=\varepsilon.}
\]

\paragraph{Step 1: Calculate the perturbed eigenvalues.}
The perturbed matrix is
\[
\widetilde A=A+E
=\begin{bmatrix}0&\varepsilon\\\varepsilon&\gamma\end{bmatrix}.
\]
Its characteristic polynomial is
\[
\begin{aligned}
\det(\widetilde A-\lambda I)
&=\det\begin{bmatrix}-\lambda&\varepsilon\\
                    \varepsilon&\gamma-\lambda\end{bmatrix}\\
&=(-\lambda)(\gamma-\lambda)-\varepsilon^2\\
&=\lambda^2-\gamma\lambda-\varepsilon^2.
\end{aligned}
\]
Setting this equal to zero and applying the quadratic formula,
we obtain
\[
\boxed{
\widetilde\lambda_1
=\frac{\gamma-\sqrt{\gamma^2+4\varepsilon^2}}{2},
\qquad
\widetilde\lambda_2
=\frac{\gamma+\sqrt{\gamma^2+4\varepsilon^2}}{2}.
}
\]
Because $\sqrt{\gamma^2+4\varepsilon^2}>\gamma$,
the first eigenvalue is negative and the second exceeds $\gamma$.
In fact, their displacements from the original eigenvalues
have the same magnitude:
\[
\begin{aligned}
|\widetilde\lambda_1-0|
&=\frac{\sqrt{\gamma^2+4\varepsilon^2}-\gamma}{2},\\
|\widetilde\lambda_2-\gamma|
&=\frac{\sqrt{\gamma^2+4\varepsilon^2}-\gamma}{2}.
\end{aligned}
\]
Rationalizing the numerator gives
\[
\boxed{
|\widetilde\lambda_1-\lambda_1(A)|
=|\widetilde\lambda_2-\lambda_2(A)|
=\frac{2\varepsilon^2}
{\sqrt{\gamma^2+4\varepsilon^2}+\gamma}
\leq\varepsilon=\norm{E}_2.
}
\]
The final inequality follows because
$\sqrt{\gamma^2+4\varepsilon^2}+\gamma\geq2\varepsilon$.
Thus we can verify Weyl's bound directly in this example.

\paragraph{Step 2: Calculate an eigenvector}

Let
\[
\widetilde v_1=
\begin{bmatrix}a\\b\end{bmatrix}\neq0
\]
be an eigenvector belonging to the smaller perturbed eigenvalue.

Recall that
\[
\widetilde\lambda_1
=
\frac{\gamma-\sqrt{\gamma^2+4\varepsilon^2}}{2}.
\]

Since $\gamma>0$ and $\varepsilon>0$, we have
\[
\gamma^2+4\varepsilon^2>\gamma^2.
\]

Taking square roots gives
\[
\sqrt{\gamma^2+4\varepsilon^2}>\gamma.
\]

Consequently,
\[
\gamma-\sqrt{\gamma^2+4\varepsilon^2}<0,
\]
and therefore
\[
\boxed{\widetilde\lambda_1<0.}
\]

Now consider the eigenvalue equation
\[
\widetilde A\widetilde v_1
=
\widetilde\lambda_1\widetilde v_1.
\]

Substituting the matrix and eigenvector gives
\[
\begin{bmatrix}
0&\varepsilon\\
\varepsilon&\gamma
\end{bmatrix}
\begin{bmatrix}a\\b\end{bmatrix}
=
\widetilde\lambda_1
\begin{bmatrix}a\\b\end{bmatrix}.
\]

Multiplying out the two rows,
\[
\boxed{
\varepsilon b=\widetilde\lambda_1 a,
\qquad
\varepsilon a+\gamma b=\widetilde\lambda_1 b.
}
\]

The first equation yields
\[
b=\frac{\widetilde\lambda_1}{\varepsilon}a.
\]

Notice that $a$ cannot be zero. Indeed, if $a=0$,
then $\varepsilon b=0$, which implies $b=0$
because $\varepsilon>0$. This would contradict
our assumption that the eigenvector is nonzero.

We are therefore free to scale the eigenvector
so that $a=1$.

An unnormalized eigenvector is
\[
\begin{bmatrix}
1\\
\widetilde\lambda_1/\varepsilon
\end{bmatrix}.
\]

Its Euclidean norm is
\[
\sqrt{
1+\left(\frac{\widetilde\lambda_1}{\varepsilon}\right)^2
}.
\]

Dividing by this norm produces the unit eigenvector
\[
\boxed{
\widetilde v_1
=
\frac{1}{
\sqrt{1+(\widetilde\lambda_1/\varepsilon)^2}
}
\begin{bmatrix}
1\\
\widetilde\lambda_1/\varepsilon
\end{bmatrix}.
}
\]

Since $\widetilde\lambda_1<0$ and $\varepsilon>0$,
the first component of this vector is positive,
while the second component is negative.

Moreover, because the vector is normalized,
the squares of its two components sum to one.

We may therefore introduce an angle
$0<\theta<\pi/2$ by defining
\[
\boxed{
\cos\theta
=
\frac{1}{
\sqrt{1+(\widetilde\lambda_1/\varepsilon)^2}
}
}
\]
and
\[
\boxed{
\sin\theta
=
\frac{-\widetilde\lambda_1/\varepsilon}{
\sqrt{1+(\widetilde\lambda_1/\varepsilon)^2}
}.
}
\]

Both expressions are strictly positive, and
\[
\begin{aligned}
\cos^2\theta+\sin^2\theta
&=
\frac{1}{
1+(\widetilde\lambda_1/\varepsilon)^2
}
+
\frac{(\widetilde\lambda_1/\varepsilon)^2}{
1+(\widetilde\lambda_1/\varepsilon)^2
}\\
&=1.
\end{aligned}
\]



Our normalized eigenvector can now be written as
\[
\boxed{
\widetilde v_1=
\begin{bmatrix}
\cos\theta\\
-\sin\theta
\end{bmatrix}.
}
\]

Notice the negative sign in the second component:
we defined $\sin\theta$ to be positive, but the
second component of $\widetilde v_1$ is negative.

Taking the ratio of our definitions gives
\[
\begin{aligned}
\tan\theta
&=\frac{\sin\theta}{\cos\theta}\\
&=-\frac{\widetilde\lambda_1}{\varepsilon}.
\end{aligned}
\]

Substituting the expression for
$\widetilde\lambda_1$, we obtain
\[
\boxed{
\tan\theta
=
\frac{
\sqrt{\gamma^2+4\varepsilon^2}-\gamma
}{
2\varepsilon
}.
}
\]

Finally, let us determine the range of $\theta$.

We already know that $\tan\theta>0$.
Furthermore, since $\gamma,\varepsilon>0$,
\[
\gamma^2+4\varepsilon^2
<
(\gamma+2\varepsilon)^2.
\]

Taking square roots,
\[
\sqrt{\gamma^2+4\varepsilon^2}
<
\gamma+2\varepsilon.
\]

Therefore,
\[
\sqrt{\gamma^2+4\varepsilon^2}-\gamma
<
2\varepsilon,
\]
which implies
\[
\boxed{0<\tan\theta<1.}
\]

Because tangent is strictly increasing on
$(0,\pi/2)$, we conclude that
\[
\boxed{
0<\theta<\frac{\pi}{4}.
}
\]

Hence, the smaller perturbed eigenvalue
corresponds to a unit eigenvector obtained
by rotating the original direction $e_1$
clockwise through an acute angle
smaller than $45^\circ$.

\paragraph{Step 3: Derive the formula for the rotation angle.}

In Step 2, we established that the normalized eigenvector
corresponding to the smaller perturbed eigenvalue can be
written as
\[
\boxed{
\widetilde v_1=
\begin{bmatrix}
\cos\theta\\
-\sin\theta
\end{bmatrix},
\qquad
0<\theta<\frac{\pi}{4}.
}
\]

We also obtained an expression for $\tan\theta$.
We now derive another useful relationship:
\[
\boxed{
\tan(2\theta)=\frac{2\varepsilon}{\gamma}.
}
\]

This formula will make the dependence of the rotation
angle on the perturbation magnitude $\varepsilon$
and the spectral gap $\gamma$ particularly transparent.

\medskip

\noindent\textbf{Returning to the eigenvalue equation.}

Recall that
\[
\widetilde A\widetilde v_1
=
\widetilde\lambda_1\widetilde v_1,
\]
where
\[
\widetilde A=
\begin{bmatrix}
0&\varepsilon\\
\varepsilon&\gamma
\end{bmatrix}.
\]

Substituting the trigonometric representation of
$\widetilde v_1$, we obtain
\[
\begin{bmatrix}
0&\varepsilon\\
\varepsilon&\gamma
\end{bmatrix}
\begin{bmatrix}
\cos\theta\\
-\sin\theta
\end{bmatrix}
=
\widetilde\lambda_1
\begin{bmatrix}
\cos\theta\\
-\sin\theta
\end{bmatrix}.
\]

Performing the matrix multiplication on the left gives
\[
\begin{bmatrix}
-\varepsilon\sin\theta\\
\varepsilon\cos\theta-\gamma\sin\theta
\end{bmatrix}
=
\begin{bmatrix}
\widetilde\lambda_1\cos\theta\\
-\widetilde\lambda_1\sin\theta
\end{bmatrix}.
\]

Equality of two vectors means that their corresponding
components must be equal.

Therefore, we obtain the following two scalar equations:
\[
\boxed{
-\varepsilon\sin\theta
=
\widetilde\lambda_1\cos\theta
}
\tag{1}
\]
and
\[
\boxed{
\varepsilon\cos\theta-\gamma\sin\theta
=
-\widetilde\lambda_1\sin\theta.
}
\tag{2}
\]

Our objective is to combine these equations in order
to eliminate the unknown eigenvalue
$\widetilde\lambda_1$.

\medskip

\noindent\textbf{First, use Equation (1).}

We have
\[
-\varepsilon\sin\theta
=
\widetilde\lambda_1\cos\theta.
\]

Since $0<\theta<\pi/4$, we know that
$\cos\theta>0$.

Therefore, dividing both sides by $\cos\theta$ gives
\[
-\varepsilon
\frac{\sin\theta}{\cos\theta}
=
\widetilde\lambda_1.
\]

Using
\[
\tan\theta=\frac{\sin\theta}{\cos\theta},
\]
we obtain
\[
\boxed{
\widetilde\lambda_1=-\varepsilon\tan\theta.
}
\tag{3}
\]

Notice that this is consistent with our calculation
in Step 2, where we established that
\[
\tan\theta=-\frac{\widetilde\lambda_1}{\varepsilon}.
\]

We will now substitute Equation (3) into Equation (2)
to eliminate $\widetilde\lambda_1$.

\medskip

\noindent\textbf{Next, use Equation (2).}

Recall that
\[
\varepsilon\cos\theta-\gamma\sin\theta
=
-\widetilde\lambda_1\sin\theta.
\]

Substituting
$\widetilde\lambda_1=-\varepsilon\tan\theta$
into the right-hand side gives
\[
\varepsilon\cos\theta-\gamma\sin\theta
=
-(-\varepsilon\tan\theta)\sin\theta.
\]

The two negative signs cancel:
\[
\varepsilon\cos\theta-\gamma\sin\theta
=
\varepsilon\tan\theta\sin\theta.
\]

Now divide the entire equation by
$\cos\theta>0$:
\[
\frac{\varepsilon\cos\theta}{\cos\theta}
-
\frac{\gamma\sin\theta}{\cos\theta}
=
\frac{\varepsilon\tan\theta\sin\theta}
{\cos\theta}.
\]

Simplifying each term,
\[
\varepsilon-\gamma\tan\theta
=
\varepsilon\tan\theta
\frac{\sin\theta}{\cos\theta}.
\]

Since $\sin\theta/\cos\theta=\tan\theta$,
the right-hand side becomes
$\varepsilon\tan^2\theta$.

Therefore,
\[
\boxed{
\varepsilon-\gamma\tan\theta
=
\varepsilon\tan^2\theta.
}
\tag{4}
\]

We have now eliminated the eigenvalue
$\widetilde\lambda_1$ completely.

Equation (4) contains only the rotation angle
$\theta$, the perturbation magnitude $\varepsilon$,
and the spectral gap $\gamma$.

\medskip

\noindent\textbf{Rearranging the equation.}

Starting from
\[
\varepsilon-\gamma\tan\theta
=
\varepsilon\tan^2\theta,
\]
move the term involving $\gamma$ to the right
and the squared tangent term to the left:
\[
\varepsilon-\varepsilon\tan^2\theta
=
\gamma\tan\theta.
\]

Factor $\varepsilon$ from the left-hand side:
\[
\boxed{
\varepsilon(1-\tan^2\theta)
=
\gamma\tan\theta.
}
\tag{5}
\]

Because Step 2 established that
$0<\theta<\pi/4$, we have
\[
0<\tan\theta<1.
\]

Consequently,
\[
1-\tan^2\theta>0.
\]

We may therefore divide Equation (5) by
$\gamma(1-\tan^2\theta)$, obtaining
\[
\boxed{
\frac{\varepsilon}{\gamma}
=
\frac{\tan\theta}{1-\tan^2\theta}.
}
\tag{6}
\]

\medskip

\noindent\textbf{Recognizing the double-angle identity.}

Recall the trigonometric identity
\[
\boxed{
\tan(2\theta)
=
\frac{2\tan\theta}{1-\tan^2\theta}.
}
\]

The right-hand side of Equation (6) is
almost identical to this expression.
The only missing factor is $2$.

Multiplying both sides of Equation (6) by $2$
gives
\[
\frac{2\varepsilon}{\gamma}
=
\frac{2\tan\theta}{1-\tan^2\theta}.
\]

But the right-hand side is precisely
$\tan(2\theta)$.

Hence,
\[
\boxed{
\tan(2\theta)=\frac{2\varepsilon}{\gamma}.
}
\]

Since
\[
0<2\theta<\frac{\pi}{2},
\]
we may apply the inverse tangent function
without ambiguity:
\[
2\theta
=
\tan^{-1}\left(\frac{2\varepsilon}{\gamma}\right).
\]

Dividing by two, we obtain the explicit formula
\[
\boxed{
\theta
=
\frac12
\tan^{-1}\left(\frac{2\varepsilon}{\gamma}\right).
}
\]

\paragraph{Geometric interpretation.}

The formula shows that the rotation angle depends
on the ratio
\[
\boxed{
\frac{\varepsilon}{\gamma}
=
\frac{\text{perturbation magnitude}}
{\text{spectral gap}}.
}
\]

For a fixed spectral gap $\gamma$, increasing
$\varepsilon$ increases the rotation angle.

Conversely, for a fixed perturbation magnitude
$\varepsilon$, decreasing $\gamma$ increases
the rotation angle.

In particular, when $\varepsilon\ll\gamma$,
the ratio $2\varepsilon/\gamma$ is small.

Using the small-angle approximation
$\tan^{-1} t\approx t$ for $t$ near zero,
we obtain
\[
\begin{aligned}
\theta
&=
\frac12
\tan^{-1}\left(\frac{2\varepsilon}{\gamma}\right)\\
&\approx
\frac12\left(\frac{2\varepsilon}{\gamma}\right)\\
&=
\frac{\varepsilon}{\gamma}.
\end{aligned}
\]

Thus,
\[
\boxed{
\theta\approx\frac{\varepsilon}{\gamma}
\qquad\text{when }\varepsilon\ll\gamma.
}
\]

\textbf{This calculation provides our first concrete
indication that eigenvector stability depends
not only on the size of a perturbation but also
on the separation between eigenvalues.}

The Davis--Kahan theorem will formalize this
observation for general symmetric matrices
and invariant subspaces.
\paragraph{Step 4: Compute the subspace distance explicitly.}
The original unit eigenvector is $e_1=(1,0)^T$, whereas
the perturbed unit eigenvector is
$\widetilde v_1=(\cos\theta,-\sin\theta)^T$.
Their inner product is $e_1^T\widetilde v_1=\cos\theta$.
By the one-dimensional projection-distance formula from
\cref{thm:one_dimensional_distance},
\[
\boxed{
\dist\bigl(\Span\{e_1\},\Span\{\widetilde v_1\}\bigr)
=\sqrt{1-(e_1^T\widetilde v_1)^2}
=\sin\theta.
}
\]
We can also express this distance entirely in terms of
$\gamma$ and $\varepsilon$. Since
$0<2\theta<\pi/2$ and
$\tan(2\theta)=2\varepsilon/\gamma$,
\[
\cos(2\theta)
=\frac{\gamma}{\sqrt{\gamma^2+4\varepsilon^2}}.
\]
The half-angle identity gives
\[
\begin{aligned}
\sin^2\theta
&=\frac{1-\cos(2\theta)}{2}\\
&=\frac12\left(1-
\frac{\gamma}{\sqrt{\gamma^2+4\varepsilon^2}}\right).
\end{aligned}
\]
Therefore,
\[
\boxed{
\dist\bigl(\Span\{e_1\},\Span\{\widetilde v_1\}\bigr)
=\sin\theta
=\sqrt{\frac12\left(1-
\frac{\gamma}{\sqrt{\gamma^2+4\varepsilon^2}}\right)}.
}
\]

\paragraph{Step 5: What happens as the spectral gap changes?}
Keep the perturbation size fixed at $\varepsilon=1$.
When $\gamma=10$, the angle equation gives
\[
\theta=\frac12\tan^{-1}\left(\frac{2}{10}\right)
\approx5.65^\circ.
\]
When $\gamma=1$, the \emph{same} perturbation gives
\[
\theta=\frac12\tan^{-1}(2)
\approx31.72^\circ.
\]
Thus a fixed perturbation can produce a much larger
eigenvector rotation when the unperturbed eigenvalues
are closer together. In this particular example,
$\theta\to45^\circ$ as $\gamma\to0^+$.
At $\gamma=0$, the original matrix has a repeated eigenvalue,
so its individual eigenvector directions are no longer unique.

When $\varepsilon\ll\gamma$, small-angle and Taylor
approximations further give
\[
\boxed{
\theta\approx\frac{\varepsilon}{\gamma},
\qquad
|\widetilde\lambda_1-\lambda_1(A)|
\approx\frac{\varepsilon^2}{\gamma}.
}
\]
The rotation is approximately first order in the
ratio $\varepsilon/\gamma$, whereas the eigenvalue
shift in this example is second order in $\varepsilon$
for a fixed positive gap.

This example explains why Weyl's inequality, despite
its strong eigenvalue bound, cannot by itself control
the rotation of eigenvectors. A meaningful eigenvector
bound must also account for \emph{spectral separation}.
That is precisely the role of the gap in the
Davis--Kahan theorem.

% ------------------------------------------------------------
% 6.2 SPECTRAL GAPS
% ------------------------------------------------------------

\subsection{The Spectral Gap and the Meaning of Stability}
\label{subsec:dk_gap}

Let $A=A^T\in\R^{n\times n}$ with $n\geq2$ have eigenvalues, counted with multiplicity,
\[
\lambda_1\leq\lambda_2\leq\cdots\leq\lambda_n,
\]
and corresponding orthonormal eigenvectors
$v_1,\ldots,v_n$.

\begin{definition}[Spectral Gap of a Simple Eigenvalue]
\label{def:dk_gap}
If $\lambda_i$ is a \emph{simple} eigenvalue of $A$,
its spectral gap is
\[
\boxed{
\gamma_i=\min_{j\neq i}|\lambda_j-\lambda_i|>0.
}
\]
Equivalently, for an interior index,
\[
\gamma_i=\min\{\lambda_i-\lambda_{i-1},
                    \lambda_{i+1}-\lambda_i\}.
\]
At either endpoint only the existing neighbor is used.
\end{definition}

The gap measures the distance from the chosen eigenvalue
to the closest \emph{other} eigenvalue. If two eigenvalues
are nearly equal, the gap is small and their eigenvectors
can be sensitive. If the gap is large, a perturbation of
the same size produces less rotation.

When an eigenvalue is repeated, its individual eigenvectors
are not uniquely distinguished. In that situation it is
more meaningful to study the invariant subspace belonging
to the entire group of repeated or clustered eigenvalues.

% ------------------------------------------------------------
% 6.3 AN ELEMENTARY EIGENVECTOR BOUND
% ------------------------------------------------------------

\subsection{A Step-by-Step Bound for a Simple Eigenvector}
\label{subsec:dk_simple}

We begin with a one-dimensional version, because it
exposes the main idea without introducing additional
matrix notation.

\begin{theorem}[Perturbation of a Simple Eigenvector]
\label{thm:dk_simple}
Let $A=A^T\in\R^{n\times n}$ and
$\widetilde A=A+E$, where $E=E^T$.
Suppose $\lambda_i(A)$ is simple, with normalized
eigenvector $v_i$ and spectral gap $\gamma_i>0$.
Set $\varepsilon=\norm{E}_2$ and assume
\[
\varepsilon<\frac{\gamma_i}{2}.
\]
Let $\widetilde v_i$ be a unit eigenvector associated
with the $i$th ordered eigenvalue of $\widetilde A$.
Then
\[
\boxed{
\dist\bigl(\Span\{v_i\},\Span\{\widetilde v_i\}\bigr)
=\sqrt{1-(v_i^T\widetilde v_i)^2}
\leq\frac{\varepsilon}{\gamma_i-\varepsilon}
\leq\frac{2\varepsilon}{\gamma_i}.
}
\]
\end{theorem}

\begin{proof}
Write $\widetilde\lambda_i=\lambda_i(\widetilde A)$.
Our goal is to estimate how much of $\widetilde v_i$
points \emph{outside} the original eigendirection.

\medskip
\noindent\textbf{Step 1: Use Weyl's inequality to locate the perturbed eigenvalue.}

Since $\widetilde A=A+E$, Weyl's inequality gives
\[
\boxed{
|\widetilde\lambda_i-\lambda_i|\leq\varepsilon.
}
\]
For every $j\neq i$, the reverse triangle inequality yields
\[
\begin{aligned}
|\lambda_j-\widetilde\lambda_i|
&\geq|\lambda_j-\lambda_i|
   -|\lambda_i-\widetilde\lambda_i|\\
&\geq\gamma_i-\varepsilon>0.
\end{aligned}
\]
Therefore the perturbed eigenvalue is separated from
all the \emph{other original eigenvalues} by at least
$\gamma_i-\varepsilon$.
The stronger assumption $2\varepsilon<\gamma_i$
also ensures that the $i$th perturbed eigenvalue remains
simple and does not exchange its order with a neighbor.

\medskip
\noindent\textbf{Step 2: Expand the perturbed eigenvector in the original eigenbasis.}

Since $v_1,\ldots,v_n$ form an orthonormal basis,
\[
\boxed{
\widetilde v_i=\sum_{j=1}^n c_jv_j,
\qquad \sum_{j=1}^n c_j^2=1.
}
\]
The component along $v_i$ is $c_i v_i$.
All remaining components are perpendicular to $v_i$:
\[
z=\sum_{j\neq i}c_jv_j,
\qquad z\perp v_i.
\]
By our one-dimensional subspace-distance formula,
\[
\boxed{
\dist\bigl(\Span\{v_i\},\Span\{\widetilde v_i\}\bigr)
=\norm{z}_2
=\left(\sum_{j\neq i}c_j^2\right)^{1/2}.
}
\]
Thus, we only need to bound the coefficients with
indices $j\neq i$.

\medskip
\noindent\textbf{Step 3: Use the two eigenvalue equations.}

The perturbed eigenvector satisfies
\[
(A+E)\widetilde v_i
=\widetilde\lambda_i\widetilde v_i.
\]
Rearranging,
\[
\boxed{
(A-\widetilde\lambda_i I)\widetilde v_i
=-E\widetilde v_i.
}
\]
Take the inner product with $v_j$ for $j\neq i$.
Because $A$ is symmetric and $Av_j=\lambda_jv_j$,
\[
\begin{aligned}
v_j^T(A-\widetilde\lambda_i I)\widetilde v_i
&=(\lambda_j-\widetilde\lambda_i)
  v_j^T\widetilde v_i\\
&=(\lambda_j-\widetilde\lambda_i)c_j.
\end{aligned}
\]
Consequently,
\[
\boxed{
(\lambda_j-\widetilde\lambda_i)c_j
=-v_j^TE\widetilde v_i,
\qquad j\neq i.
}
\]

\medskip
\noindent\textbf{Step 4: Combine the gap and perturbation bounds.}

Squaring the preceding equation and summing over
$j\neq i$ gives
\[
\sum_{j\neq i}
(\lambda_j-\widetilde\lambda_i)^2c_j^2
=
\sum_{j\neq i}(v_j^TE\widetilde v_i)^2.
\]
The lower bound from Step 1 implies
\[
(\gamma_i-\varepsilon)^2
\sum_{j\neq i}c_j^2
\leq
\sum_{j\neq i}(v_j^TE\widetilde v_i)^2.
\]
By orthonormality and the Pythagorean theorem,
\[
\sum_{j\neq i}(v_j^TE\widetilde v_i)^2
\leq\norm{E\widetilde v_i}_2^2
\leq\norm{E}_2^2=\varepsilon^2.
\]
Therefore,
\[
\left(\sum_{j\neq i}c_j^2\right)^{1/2}
\leq\frac{\varepsilon}{\gamma_i-\varepsilon}.
\]
Using Step 2, we conclude that
\[
\boxed{
\dist\bigl(\Span\{v_i\},\Span\{\widetilde v_i\}\bigr)
\leq\frac{\varepsilon}{\gamma_i-\varepsilon}.
}
\]
Finally, $\varepsilon<\gamma_i/2$ implies
$\gamma_i-\varepsilon>\gamma_i/2$, so
\[
\frac{\varepsilon}{\gamma_i-\varepsilon}
\leq\frac{2\varepsilon}{\gamma_i}.
\]
This completes the proof.
\end{proof}

\paragraph{What should we remember from this proof?}

The argument combined two ideas. Weyl's inequality
controlled the movement of the \emph{eigenvalue},
ensuring that it remained separated from the other
eigenvalues. The perturbed eigenvalue equation then
controlled how much of the new eigenvector could lie
outside the original eigendirection.

The relevant quantity is therefore the ratio
\[
\boxed{
\frac{\text{perturbation size}}
     {\text{spectral separation}}.
}
\]
The denominator matters: eigenvectors are not equally
stable for all symmetric matrices.

% ------------------------------------------------------------
% 6.4 SUBSPACE SIN THETA THEOREM
% ------------------------------------------------------------

\subsection{From One Eigenvector to an Invariant Subspace}
\label{subsec:dk_subspaces}

We now extend the same reasoning to an entire invariant
subspace. As in Section~\ref{sec:subspace_geometry},
let $W\in\R^{n\times r}$ have orthonormal columns.
The subspace $\mathcal W=\range(W)$ is invariant under
$A$ when
\[
\boxed{AW=W\Lambda,}
\]
where $\Lambda\in\R^{r\times r}$ is symmetric. We
may take its columns to be eigenvectors, so that
$\Lambda$ is diagonal.

Let $\widetilde W\in\R^{n\times r}$ contain the
orthonormal eigenvectors of $\widetilde A$ corresponding
to the eigenvalues we wish to compare. Complete its
columns to an orthonormal basis
$[\widetilde W,\widetilde W_\perp]$ of $\R^n$.
Then
\[
\widetilde A\widetilde W_\perp
=\widetilde W_\perp\widetilde\Lambda_\perp,
\]
where $\widetilde\Lambda_\perp$ is diagonal and contains
the remaining perturbed eigenvalues.

\begin{theorem}[Davis--Kahan $\sin\Theta$ Theorem]
\label{thm:davis_kahan_interval}
Let $A=A^T$ and $\widetilde A=A+E$ with $E=E^T$.
Suppose $\mathcal W=\range(W)$ is an $r$-dimensional
invariant subspace of $A$, with $1\leq r<n$, and
its associated eigenvalues lie in an interval $[a,b]$.
Suppose the eigenvalues of $\widetilde A$ associated
with the orthogonal complement of
$\widetilde{\mathcal W}=\range(\widetilde W)$ lie
outside the enlarged interval
\[
(a-\delta,b+\delta),\qquad\delta>0.
\]
Assume $\dim\widetilde{\mathcal W}=r$. Then
\[
\boxed{
\begin{aligned}
\dist(\mathcal W,\widetilde{\mathcal W})
&=\norm{WW^T-\widetilde W\widetilde W^T}_2\\
&=\sin\theta_{\max}\\
&\leq\frac{\norm{E}_2}{\delta}.
\end{aligned}
}
\]
\end{theorem}

\begin{proof}
We explain the key calculation rather than assuming
any prior knowledge of subspace perturbation theory.

\medskip
\noindent\textbf{Step 1: Measure the component outside the perturbed subspace.}

The columns of $\widetilde W_\perp$ span
$\widetilde{\mathcal W}^{\perp}$. Define
\[
\boxed{X=\widetilde W_\perp^T W.}
\]
The entries of $X$ measure how much the original
basis vectors overlap with directions perpendicular
to the perturbed subspace. Since both subspaces have
dimension $r$, the projection-distance identity
from \cref{subsec:distance_principal_angles} gives
\[
\boxed{
\dist(\mathcal W,\widetilde{\mathcal W})
=\norm{(I-\widetilde W\widetilde W^T)W}_2
=\norm{X}_2.
}
\]
The problem is therefore reduced to bounding $\norm{X}_2$.

\medskip
\noindent\textbf{Step 2: Derive an equation satisfied by $X$.}

From $\widetilde A=A+E$ and $AW=W\Lambda$,
\[
EW=\widetilde A W-W\Lambda.
\]
Multiply on the left by $\widetilde W_\perp^T$:
\[
\widetilde W_\perp^T EW
=\widetilde W_\perp^T\widetilde A W
  -\widetilde W_\perp^T W\Lambda.
\]
Because $\widetilde A$ is symmetric and
$\widetilde A\widetilde W_\perp
=\widetilde W_\perp\widetilde\Lambda_\perp$,
we have
$\widetilde W_\perp^T\widetilde A
=\widetilde\Lambda_\perp\widetilde W_\perp^T$.
Thus
\[
\boxed{
\widetilde W_\perp^T EW
=\widetilde\Lambda_\perp X-X\Lambda.
}
\]
The left side has norm at most $\norm{E}_2$, since
$W$ and $\widetilde W_\perp$ have orthonormal columns.

\medskip
\noindent\textbf{Step 3: Explain how spectral separation controls $X$.}

Let $c=(a+b)/2$ and $h=(b-a)/2$. The spectrum of
$\Lambda$ lies in $[c-h,c+h]$, so
\[
\norm{\Lambda-cI}_2\leq h.
\]
Every eigenvalue of $\widetilde\Lambda_\perp$ is
at least $\delta$ outside $[a,b]$. Consequently,
\[
\abs{\widetilde\lambda-c}\geq h+\delta
\quad\text{for all }
\widetilde\lambda\in\spec(\widetilde\Lambda_\perp).
\]
It follows that
\[
\norm{(\widetilde\Lambda_\perp-cI)X}_2
\geq(h+\delta)\norm{X}_2.
\]
Indeed, $\widetilde\Lambda_\perp-cI$ is invertible,
and its inverse has norm at most $(h+\delta)^{-1}$.

Now rewrite the equation from Step 2 as
\[
\widetilde W_\perp^T EW
=(\widetilde\Lambda_\perp-cI)X
-X(\Lambda-cI).
\]
The reverse triangle inequality gives
\[
\begin{aligned}
\norm{\widetilde W_\perp^T EW}_2
&\geq\norm{(\widetilde\Lambda_\perp-cI)X}_2
      -\norm{X(\Lambda-cI)}_2\\
&\geq(h+\delta)\norm{X}_2-h\norm{X}_2\\
&=\delta\norm{X}_2.
\end{aligned}
\]
Since $\norm{\widetilde W_\perp^T EW}_2\leq\norm{E}_2$,
\[
\boxed{\norm{X}_2\leq\frac{\norm{E}_2}{\delta}.}
\]
Step 1 identifies $\norm{X}_2$ with the projection
distance and with $\sin\theta_{\max}$, completing
the proof.
\end{proof}

\paragraph{Why does the interval condition matter?}

The theorem compares an original group of eigenvalues
with the \emph{complement} of a chosen perturbed group.
There must be a positive separation $\delta$ between
the original selected interval and the perturbed
complement. Without this hypothesis, the denominator
might vanish and the estimate would provide no control.

% ------------------------------------------------------------
% 6.5 A GAP CONDITION INVOLVING ONLY A
% ------------------------------------------------------------

\subsection{A Practical Bound Using the Original Spectral Gap}
\label{subsec:dk_population_gap}

The preceding theorem involves eigenvalues of both
$A$ and $\widetilde A$. In applications, we may prefer
a condition that refers only to the original matrix.
Weyl's inequality allows us to obtain such a bound.

\begin{corollary}[Stability of an Isolated Eigenvalue Cluster]
\label{cor:dk_cluster}
Let $A=A^T$, $E=E^T$ and
$\widetilde A=A+E$. Select a consecutive group
\[
\lambda_p(A),\ldots,\lambda_q(A),
\qquad 1\leq p\leq q\leq n,
\]
with $r=q-p+1<n$. Assume that this group is
separated from the remaining eigenvalues by a
positive gap
\[
\boxed{
\gamma=
\min\!\left\{
\begin{array}{ll}
\lambda_p-\lambda_{p-1},&\text{if }p>1,\\
\lambda_{q+1}-\lambda_q,&\text{if }q<n
\end{array}\right\}>0,
}
\]
where a missing endpoint term is simply omitted.
Let $\mathcal W$ be the subspace spanned by the
corresponding original eigenvectors and let
$\widetilde{\mathcal W}$ be spanned by the perturbed
eigenvectors with the same ordered indices $p,\ldots,q$.

If $\varepsilon=\norm{E}_2<\gamma/2$, then
\[
\boxed{
\begin{aligned}
\dist(\mathcal W,\widetilde{\mathcal W})
&\leq\frac{\varepsilon}{\gamma-\varepsilon}\\
&\leq\frac{2\varepsilon}{\gamma}.
\end{aligned}
}
\]
\end{corollary}

\begin{proof}
Set $a=\lambda_p(A)$ and $b=\lambda_q(A)$.
Weyl's inequality gives
\[
\abs{\lambda_j(\widetilde A)-\lambda_j(A)}
\leq\varepsilon
\qquad (j=1,\ldots,n).
\]
If $p>1$, every perturbed eigenvalue with $j<p$
satisfies
\[
\lambda_j(\widetilde A)
\leq\lambda_{p-1}(A)+\varepsilon
\leq a-(\gamma-\varepsilon).
\]
If $q<n$, every perturbed eigenvalue with $j>q$
satisfies
\[
\lambda_j(\widetilde A)
\geq\lambda_{q+1}(A)-\varepsilon
\geq b+(\gamma-\varepsilon).
\]
Thus the perturbed eigenvalues outside the chosen
group lie outside
$(a-\delta,b+\delta)$ with
$\delta=\gamma-\varepsilon>0$.
The smallness condition also keeps the selected group
separated from its perturbed neighbors.
Applying \cref{thm:davis_kahan_interval} yields
\[
\dist(\mathcal W,\widetilde{\mathcal W})
\leq\frac{\varepsilon}{\gamma-\varepsilon}.
\]
Since $\varepsilon<\gamma/2$, we also have
$1/(\gamma-\varepsilon)\leq2/\gamma$.
\end{proof}

\begin{remark}[Why the Gap Must Be Defined for a Group]
If the selected eigenvalues contain repetitions,
there may be no unique way to choose their individual
eigenvectors. The corollary therefore measures the
invariant subspace associated with the \emph{whole}
cluster. Eigenvalues can be very close to one another
\emph{inside} the group; what matters is that the
group is separated from eigenvalues \emph{outside}
it. This is particularly useful in PCA, where one
usually cares about a leading principal-component
subspace rather than any particular orthonormal
basis for that subspace.
\end{remark}

% ------------------------------------------------------------
% 6.6 A WORKED EXAMPLE AND INTERPRETATION
% ------------------------------------------------------------

\subsection{A Numerical Illustration and the Main Lesson}
\label{subsec:dk_example}

Return to the $2\times2$ example with
$\gamma=1$ and $\varepsilon=0.1$:
\[
A=\begin{bmatrix}0&0\\0&1\end{bmatrix},
\qquad
E=\begin{bmatrix}0&0.1\\0.1&0\end{bmatrix}.
\]
The original eigenvalue gap is $\gamma_1=1$, and
$\norm{E}_2=0.1<\gamma_1/2$.
The smaller perturbed eigenvalue is
\[
\widetilde\lambda_1
=\frac{1-\sqrt{1.04}}{2}\approx-0.009902.
\]
The corresponding eigendirection is rotated by
\[
\theta=\frac12\arctan(0.2)\approx5.655^\circ.
\]
Therefore,
\[
\boxed{
\dist\bigl(\Span\{e_1\},\Span\{\widetilde v_1\}\bigr)
=\sin\theta\approx0.09854.
}
\]
Our perturbation bound predicts
\[
\dist\bigl(\Span\{e_1\},\Span\{\widetilde v_1\}\bigr)
\leq\frac{0.1}{1-0.1}=\frac19\approx0.11111,
\]
which is consistent with the direct calculation.

We can summarize the distinction between Weyl's
inequality and Davis--Kahan as follows:
\[
\boxed{
\begin{array}{c}
\text{Weyl: eigenvalue movement is controlled by }\norm{E}_2,\\[4pt]
\text{Davis--Kahan: subspace rotation is controlled by}
\ \norm{E}_2/\text{gap}.
\end{array}
}
\]
The two theorems address different aspects of
spectral stability. Weyl tells us that a small
perturbation cannot move the ordered eigenvalues
very far. Davis--Kahan explains why small
perturbations may nevertheless rotate eigenvectors
substantially when the relevant spectral gap is small.

In the next section we will examine these conclusions
numerically in MATLAB, varying both the perturbation
magnitude and the spectral gap and comparing the
observed subspace distance with the theoretical bounds.

% ============================================================
% SECTION 7: NUMERICAL EXPERIMENTS IN MATLAB
% ============================================================
\section{Numerical Experiments in MATLAB}
\label{sec:numerical_experiments}

The preceding sections developed three complementary ideas:
Courant--Fischer identifies ordered eigenvalues through
variational optimization, Weyl controls the movement of
ordered eigenvalues under symmetric perturbations, and
Davis--Kahan controls invariant-subspace rotations when
there is sufficient spectral separation. We now check
these statements numerically and observe how their
hypotheses influence the results.


All experiments use real symmetric matrices, and all
eigenvalues are assumed to be ordered increasingly:
\[
\lambda_1(A)\leq\lambda_2(A)\leq\cdots\leq\lambda_n(A).
\]

\subsection{How to Reproduce the Experiments}
\label{subsec:matlab_setup}

Download the MATLAB scripts accompanying this note
and place all seven files in one working directory:
\begin{lstlisting}[language=Matlab]
courant_fisher_minmax.m
courant_fisher_maxmin.m
experiment_courant_fischer.m
experiment_principal_angles.m
experiment_weyl.m
experiment_davis_kahan.m
run_all_experiments.m
\end{lstlisting}

Open MATLAB, change its Current Folder to the directory
containing these files, and type, for example,
\begin{lstlisting}[language=Matlab]
experiment_principal_angles
experiment_courant_fischer
experiment_weyl
experiment_davis_kahan
\end{lstlisting}
Each file is a script that can run on its own. To execute
all four, enter \texttt{run\_all\_experiments}. The plots
are displayed in MATLAB and saved as PNG images in the
current directory. Numerical values also appear in the
Command Window. The scripts use MATLAB's built-in
\texttt{eig}, \texttt{svd}, and \texttt{norm} as independent
checks of the mathematical claims. The Courant--Fischer
experiment additionally calls the two educational
implementations developed locally.

For reproducibility, the scripts use \texttt{rng} when
random matrices or trial subspaces are needed. However,
the experimental eigenpair solvers should not be treated
as guaranteed convergent algorithms. Inspect the printed
errors and residuals rather than assuming that an iterate
has converged because its objective value changes slowly.

\begin{remark}[How to Read the Numerical Results]
For an approximate eigenpair $(\widehat\lambda,\widehat v)$,
where $\norm{\widehat v}_2=1$, the residual
\[
\boxed{\norm{A\widehat v-\widehat\lambda\widehat v}_2}
\]
measures how closely it satisfies the original eigenvalue
equation. A small eigenvalue error alone does not imply
a small eigenvector error. We therefore report both when
we test a proposed eigensolver. For the perturbation
experiments we use the spectral norm and the distance
between \emph{projectors}, not the Euclidean distance
between arbitrarily signed eigenvectors.
\end{remark}

\subsection{The Original Courant--Fischer Implementations}
\label{subsec:matlab_cf_original}

The following two MATLAB programs are the educational
implementations developed during our study of the
Courant--Fischer theorem. We reproduce them here without
changing their numerical methods, so that readers can
study and experiment with their actual steps.

Recall the minimax characterization:
\[
\lambda_k(A)=
\min_{\dim\mathcal S=k}\;\max_{x\in\mathcal S\setminus\{0\}}R_A(x).
\]
If $Q\in\R^{n\times k}$ has orthonormal columns and
$\mathcal S=\range(Q)$, then writing $x=Qy$ gives
\[
R_A(Qy)=\frac{y^T(Q^TAQ)y}{y^Ty},
\qquad
\max_{x\in\mathcal S\setminus\{0\}}R_A(x)
=\lambda_{\max}(Q^TAQ).
\]
The minimax code searches over trial subspaces by
updating an orthonormal basis $Q$ in a descent direction.
Its inner eigenvalue estimate comes from shifted power
iteration on $Q^TAQ$; modified Gram--Schmidt is used
after each step to restore orthonormality.

\paragraph{Original minimax MATLAB function.}
The source below is saved as
\texttt{courant\_fisher\_minmax.m}.
\begin{lstlisting}[language=Matlab,caption={Original educational Courant--Fischer minimax implementation.},label={lst:cf_minmax}]
function [curr_lambda, curr_v] = courant_fisher_minmax(A, k, step_size, max_outer_iter)

	% This function gets the k-th eigenvalue and eigenvector pair of a symmetrix Matrix A.
	% usage : [lambda_k, v_k] = courant_fisher_minmax(A, k, step_size, max_outer_iter)
	% step_size and max_outer_iter are optional arguments.
	% lambda1 <= lambda2 <= ... <= lambda_k <= ...<= lambda_n
	
	if nargin < 3
		step_size = 0.05;
	end

	if nargin < 4 
		max_outer_iter = 2000;
	end

	A = (A + A') / 2;
	[n, ~] = size(A);

	V = randn(n, k); % We limit our search on k dimensional subspace. we get the maximum eigenvalue on this subspace, and then minimize it over all such k dimensional subspace--> the core of Courant_Fisher Algorithm.
	% code to be written here after completion of the local functions.

	Q = local_modified_GS(V);
	t = 0;

	for iter = 1 : max_outer_iter 
		M = Q' * A * Q;
		if iter > 1
			old_iter = curr_lambda;
		end
		[curr_lambda, curr_v] = local_power_iter(M);

		if iter > 1 
			if abs(old_iter - curr_lambda) < 1e-6
				t = t + 1;
			else 
				t = 0;
			end
		end
		 
		if t == 2 
			break;
		end

		grad_q = A * Q * curr_v * curr_v';
		Q = Q - step_size * grad_q;
		Q = local_modified_GS(Q);

		% just tracking the progress 
		if mod(iter, 100) == 0
			fprintf("Iteration %d : max R(x) = %f\n", iter, curr_lambda);
		end
	end

	M = Q' * A * Q;
	[curr_lambda, curr_v] = local_power_iter(M, 200);

	% mapping back to R^{n};
	curr_v = Q * curr_v;
	curr_v = curr_v / norm(curr_v);

end

% Local functions 


function [lambda, y] = local_power_iter(A, n)

	if nargin < 2 
		n = 100;
	end

	k = size(A, 1);
	y = randn(k, 1);
	y = y / norm(y);

	s = max(sum(abs(A), 2));
	Ashifted = A + s * eye(k);
	
	lambda_old = y' * Ashifted * y;


	for index = 1 : n
		y = Ashifted * y;
		y = y / norm(y);
		lambda = y' * Ashifted * y;
		if abs(lambda - lambda_old) < 1e-6 
			break;
		end

		lambda_old = lambda;
	end

	lambda = lambda - s;
	
end


function Q = local_modified_GS(A)
	[n, k] = size(A);
	Q = zeros(n, k);

	for index = 1 : k 
		current_vector = A(:, index);

		for jj = 1 : index - 1 
			r = Q(:, jj)' * current_vector;
			current_vector = current_vector - r * Q(:, jj);
		end
		Q(:, index) = current_vector / norm(current_vector);
	end
end
\end{lstlisting}

For the maximin principle we instead use
\[
\lambda_k(A)=
\max_{\dim\mathcal T=n-k+1}\;\min_{x\in\mathcal T\setminus\{0\}}R_A(x).
\]
Thus the trial basis now has $n-k+1$ columns, the
inner calculation searches for the \emph{smallest}
eigenvalue of $Q^TAQ$, and the outer update is an
\emph{ascent} step.

\paragraph{Original maximin MATLAB function.}
The source below is saved as
\texttt{courant\_fisher\_maxmin.m}.
\begin{lstlisting}[language=Matlab,caption={Original educational Courant--Fischer maximin implementation.},label={lst:cf_maxmin}]
function [curr_lambda, curr_v] = courant_fisher_maxmin(A, k, step_size, max_outer_iter)
	% This function gets the kth eigenvalue-eigenvector
	% pair of a symmetric matrix using Courant_Fisher Maxmin procedure.
	% usage : [lambda_k, v_k] = courant_fisher_maxmin(A, k, step_size, max_outer_iter)
	% step_size and max_outer_iter are optional arguments.
	% lambda1 <= lambda2 <= ... <= lambda_k <= ...<= lambda_n

	if nargin < 3
		step_size = 0.05;
	end

	if nargin < 4 
		max_outer_iter = 2000;
	end

	A = (A + A') / 2;
	[n, ~] = size(A);

	% subspace dimension for maxmin is (n - k + 1)
	m = n - k + 1;
	V = randn(n, m);
	Q = local_modified_GS(V);
	t = 0;

	for iter = 1 : max_outer_iter
		M = Q' * A * Q; 
		if iter > 1
			old_iter = curr_lambda;
		end

		[curr_lambda, curr_v] = local_power_iter_min(M);

		if iter > 1
			if abs(old_iter - curr_lambda) < 1e-6
				t = t + 1;
			else 
				t = 0;
			end

		end

		if t == 2
			break;
		end

		grad_q = A * (Q * curr_v) * curr_v';

		% Now we do the gradient ascent
		Q = Q + step_size * grad_q;

		Q = local_modified_GS(Q);

		if mod(iter, 100) == 0
			fprintf("Iteration %d : min R(x) = %f\n", iter, curr_lambda);
		end
	end

	M = Q' * A * Q;
	[curr_lambda, curr_v] = local_power_iter_min(M, 200);
	curr_v = Q * curr_v; % x = Py;
	curr_v = curr_v / norm(curr_v);
end





% Local function for getting the minimum eigenvalue of 
% a symmetric matrix
% we will first shift all eigenvalue to the negative side
% then we get the maximum eigenvalue, and then shift to the positive side

function [lambda, y] = local_power_iter_min(M, max_iter)

	if nargin < 2
		max_iter = 100;
	end

	m = size(M, 1);
	y = randn(m, 1);
	y = y / norm(y);
	% shift
	s = max(sum(abs(M), 2));
	M = M - s * eye(m);
	lambda_old = y' * M * y;

	for index = 1 : max_iter
		y = M * y;
		y = y / norm(y);
		lambda = y' * M * y;

		if abs(lambda - lambda_old) < 1e-6
			break;
		end

		lambda_old = lambda;
	end

	lambda = lambda + s;
end


% Local function for doing Gram-schmidt procedure
function Q = local_modified_GS(A)
	[n, k] = size(A);
	Q = zeros(n, k);

	for index = 1 : k 
		current_vector = A(:, index);

		for jj = 1 : index - 1 
			r = Q(:, jj)' * current_vector;
			current_vector = current_vector - r * Q(:, jj);
		end
		Q(:, index) = current_vector / norm(current_vector);
	end
end
\end{lstlisting}

\begin{remark}[Why These Are Educational Rather Than Production Solvers]
The two programs illustrate how to turn a mathematical
variational principle into an iterative search. But a
variational characterization is \emph{not} itself a
convergence theorem for a particular algorithm.
The objective $Q\mapsto\lambda_{\max}(Q^TAQ)$ can be
nondifferentiable at multiple leading Ritz values;
repeated orthogonalization may also lose accuracy if
its input columns become nearly dependent. Shifted
power iteration can converge slowly, and monitoring
changes in a Rayleigh quotient alone can trigger
premature stopping. Moreover, reaching the correct
Ritz value does not by itself guarantee that the
reported vector has a small residual for the full
matrix. These implementations are therefore best
used for small illustrative matrices and accompanied
by independent residual and eigenvalue checks.
\end{remark}

\subsection{Experiment 1: Checking Both Variational Inequalities}
\label{subsec:matlab_cf_experiment}

Let $A=\diag(1,3,7)$ and $k=2$. The exact answer is
$\lambda_2(A)=3$. For every two-dimensional subspace
with orthonormal basis $Q$, the Courant--Fischer theorem
implies
\[
\boxed{\lambda_{\max}(Q^TAQ)\geq3.}
\]
For every two-dimensional subspace with orthonormal
basis $Z$, its maximin counterpart implies
\[
\boxed{\lambda_{\min}(Z^TAZ)\leq3.}
\]
The extreme cases are attained by the \emph{special}
subspaces $\Span\{e_1,e_2\}$ and
$\Span\{e_2,e_3\}$, respectively. The first part
of the script samples 250 random planes and checks
both inequalities. Random planes typically give
strict inequalities; randomly sampling the \emph{exact}
optimal plane has probability zero in exact arithmetic.
The second part tests the original educational solvers,
prints their estimates, and calculates the eigenpair
residuals $\norm{A\widehat v-\widehat\lambda\widehat v}_2$.

\begin{lstlisting}[language=Matlab,caption={Experiment 1: random trial subspaces and the two educational solvers.},label={lst:experiment_cf}]
% experiment_courant_fischer.m
% Courant--Fischer: verify min-max/max-min and test educational solvers.
% Put this file beside courant_fisher_minmax.m and courant_fisher_maxmin.m.
clear; clc; close all;
rng(7); % Reproducible random trial subspaces and initial guesses
A = diag([1, 3, 7]);
n = size(A,1);
k = 2;
[V,D] = eig(A);
[lambda,order] = sort(diag(D),'ascend');
V = V(:,order);
truth = lambda(k);

% For any k-dimensional Q: lambda_max(Q'*A*Q) >= lambda_k(A).
% For any (n-k+1)-dimensional Z: lambda_min(Z'*A*Z) <= lambda_k(A).
N = 250;
trial_minmax = zeros(N,1);
trial_maxmin = zeros(N,1);
for j = 1:N
    [Q,~] = qr(randn(n,k),0);
    [Z,~] = qr(randn(n,n-k+1),0);
    trial_minmax(j) = max(eig(Q'*A*Q));
    trial_maxmin(j) = min(eig(Z'*A*Z));
end
Qopt = V(:,1:k);
Zopt = V(:,k:n);
opt_minmax = max(eig(Qopt'*A*Qopt));
opt_maxmin = min(eig(Zopt'*A*Zopt));

fprintf('True lambda_%d = %.8f\n', k, truth);
fprintf('Optimal min-max: %.8f; optimal max-min: %.8f\n', opt_minmax,opt_maxmin);
fprintf('Smallest sampled max: %.8f (should be >= %.8f)\n',min(trial_minmax),truth);
fprintf('Largest sampled min: %.8f (should be <= %.8f)\n',max(trial_maxmin),truth);
assert(all(trial_minmax >= truth-1e-10));
assert(all(trial_maxmin <= truth+1e-10));

figure('Name','Courant-Fischer trial subspaces');
plot(1:N,trial_minmax,'.','DisplayName','Max on random 2D subspace'); hold on;
plot(1:N,trial_maxmin,'.','DisplayName','Min on random 2D subspace');
yline(truth,'k--','DisplayName','True \lambda_2');
xlabel('Trial subspace'); ylabel('Restricted Rayleigh extremum');
legend('Location','best'); grid on;
saveas(gcf,'fig_courant_fischer.png');

% Test the original educational gradient implementations.
% These routines are nonconvex, initialization-sensitive numerical
% experiments and do not guarantee accurate eigenpairs.
step_size = 0.02;
max_outer_iter = 400;
[lam_mm, v_mm] = courant_fisher_minmax(A,k,step_size,max_outer_iter);
[lam_mn, v_mn] = courant_fisher_maxmin(A,k,step_size,max_outer_iter);
res_mm = norm(A*v_mm-lam_mm*v_mm,2);
res_mn = norm(A*v_mn-lam_mn*v_mn,2);
fprintf('\nEducational solvers (results may vary):\n');
fprintf('min-max estimate %.8f, abs error %.2e, eigenpair residual %.2e\n', ...
    lam_mm,abs(lam_mm-truth),res_mm);
fprintf('max-min estimate %.8f, abs error %.2e, eigenpair residual %.2e\n', ...
    lam_mn,abs(lam_mn-truth),res_mn);
if (res_mm > 1e-3) || (res_mn > 1e-3)
    fprintf('A small objective change does not necessarily imply an accurate eigenpair.\n');
end
\end{lstlisting}

\paragraph{What should the reader observe?}
All randomly sampled \emph{minimax} values should lie
at or above the horizontal line $\lambda_2=3$, while
all sampled \emph{maximin} values should lie at or
below it, up to rounding errors. Inserting the
special eigenvector subspaces gives equality exactly,
within floating-point accuracy. The two iterative
routines may approach $3$, but their printed residuals
are the appropriate check that their output actually
approximates an eigenpair. Change $k$, the initial
random seed, the step size, or the iteration limit
and examine how the results change.

\subsection{Experiment 2: Principal Angles and Projector Distance}
\label{subsec:matlab_angles}

Before studying perturbations, verify the subspace
distance formula established in
\cref{subsec:distance_principal_angles}. Consider
$\mathcal W=\Span\{e_1,e_2\}\subset\R^4$ and
\[
\mathcal U_\theta
=\Span\{e_1,\cos\theta\,e_2+\sin\theta\,e_3\}.
\]
These are both two-dimensional subspaces. Their
principal angles are $0$ and $\theta$, so
\[
\boxed{
\norm{WW^T-U_\theta U_\theta^T}_2
=\sin\theta
=\sqrt{1-\sigma_{\min}(W^TU_\theta)^2}.
}
\]
The following script evaluates both sides for angles
between $0^\circ$ and $85^\circ$.

\begin{lstlisting}[language=Matlab,caption={Experiment 2: projector distance and principal angles.},label={lst:experiment_angles}]
% experiment_principal_angles.m
% Compare projector distance with sin(largest principal angle).
clear; clc; close all;
angles_deg = 0:5:85;
proj_dist = zeros(size(angles_deg));
svd_dist = zeros(size(angles_deg));
for j = 1:numel(angles_deg)
    theta = angles_deg(j)*pi/180;
    W = [1 0;0 1;0 0;0 0];
    U = [1 0;0 cos(theta);0 sin(theta);0 0];
    % Orthonormal columns, equal dimension 2.
    proj_dist(j) = norm(W*W' - U*U',2);
    sigmas = svd(W'*U);
    svd_dist(j) = sqrt(max(0,1-min(sigmas)^2));
end
fprintf('Max |projector distance - SVD formula| = %.3e\n', ...
    max(abs(proj_dist-svd_dist)));
assert(max(abs(proj_dist-svd_dist)) < 1e-10);
figure('Name','Subspace distance and principal angles');
plot(angles_deg,proj_dist,'o-'); hold on;
plot(angles_deg,sin(angles_deg*pi/180),'k--');
xlabel('Largest principal angle (degrees)');
ylabel('Subspace distance');
legend('||WW'' - UU''||_2','sin(\theta_{max})','Location','best');
grid on;
saveas(gcf,'fig_principal_angles.png');
\end{lstlisting}

\paragraph{What should the reader observe?}
The computed projector distance, the singular-value
formula, and $\sin\theta$ should agree to floating-point
precision. The distance increases from $0$ toward $1$
as the largest principal angle increases toward $90^\circ$.
This agreement also illustrates that the distance does
not depend on the signs of basis vectors.

\subsection{Experiment 3: Numerical Verification of Weyl's Inequality}
\label{subsec:matlab_weyl}

Choose a symmetric $6\times6$ matrix $A$ and a
symmetric matrix $H$ normalized so that
$\norm{H}_2=1$. Set $E=tH$ for several values
of $t\geq0$, so $\norm{E}_2=t$. The theorem gives
\[
\boxed{
\abs{\lambda_i(A+tH)-\lambda_i(A)}\leq t,
\qquad i=1,\ldots,6.
}
\]
The stronger version gives
\[
\lambda_i(A)+\lambda_1(E)
\leq\lambda_i(A+E)
\leq\lambda_i(A)+\lambda_n(E).
\]
The following program checks both claims and plots
all six eigenvalue errors against the same bound.

\begin{lstlisting}[language=Matlab,caption={Experiment 3: Weyl's inequality for several perturbation sizes.},label={lst:experiment_weyl}]
% experiment_weyl.m
% Weyl's eigenvalue perturbation bound for real symmetric matrices.
clear; clc; close all;
rng(12);
n = 6;
X = randn(n);
A = (X+X')/2;
H = randn(n);
H = (H+H')/2;
H = H/norm(H,2); % ||H||_2 = 1
strengths = [0 0.01 0.03 0.06 0.10 0.20 0.40];
lambdaA = sort(eig(A),'ascend');
errors = zeros(numel(strengths),n);

fprintf('strength    ||E||_2     max eigenvalue error    within bound?\n');
for j = 1:numel(strengths)
    E = strengths(j)*H;
    lambdaPerturbed = sort(eig(A+E),'ascend');
    errors(j,:) = abs(lambdaPerturbed-lambdaA).';
    upper = norm(E,2);
    % Strong Weyl: lambda_k(A)+lambda_min(E) <= lambda_k(A+E)
    %                         <= lambda_k(A)+lambda_max(E).
    lambdaE = sort(eig(E),'ascend');
    strong_ok = all(lambdaPerturbed >= lambdaA+lambdaE(1)-1e-10) && ...
                all(lambdaPerturbed <= lambdaA+lambdaE(end)+1e-10);
    simple_ok = max(errors(j,:)) <= upper + 1e-10;
    fprintf('%8.3f  %11.5f  %20.7f   %d\n', ...
        strengths(j),upper,max(errors(j,:)),simple_ok && strong_ok);
    assert(simple_ok && strong_ok,'A Weyl bound failed.');
end

figure('Name','Weyl bound');
plot(strengths,errors,'o-'); hold on;
plot(strengths,strengths,'k--','LineWidth',1.8);
xlabel('Perturbation spectral norm ||E||_2');
ylabel('Absolute change in an ordered eigenvalue');
title('Weyl: every eigenvalue error lies below the dashed bound');
grid on;
saveas(gcf,'fig_weyl.png');
\end{lstlisting}

\paragraph{What should the reader observe?}
All six eigenvalue-error curves should lie on or below
the dashed line $y=t$, apart from negligible rounding
error. Increasing $t$ permits larger changes, but
Weyl's inequality is a \emph{bound}, not a prediction
that every eigenvalue must move by exactly $t$.
Try replacing $H$ with $I$: then every eigenvalue
moves by exactly $t$, demonstrating that the norm
bound can be sharp. Try also changing the random seed
or matrix dimension; the inequality must continue
to hold as long as $A$ and $E$ are symmetric.

\subsection{Experiment 4: Davis--Kahan and the Spectral Gap}
\label{subsec:matlab_dk}

We return to the explicit $2\times2$ example from
\cref{subsec:dk_motivation}. With $\varepsilon=0.1$,
let
\[
A=\begin{bmatrix}0&0\\0&\gamma\end{bmatrix},
\qquad
E=\begin{bmatrix}0&\varepsilon\\\varepsilon&0\end{bmatrix}.
\]
We hold the perturbation norm fixed at $0.1$ and
vary the original spectral gap $\gamma>0$.
The exact rotation satisfies
\[
\theta=\frac12\arctan\left(\frac{2\varepsilon}{\gamma}\right),
\qquad
\dist(\Span\{e_1\},\Span\{\widetilde v_1\})=\sin\theta.
\]
Under the sufficient condition
$\varepsilon<\gamma/2$, the bound from
\cref{cor:dk_cluster} becomes
\[
\boxed{
\dist(\Span\{e_1\},\Span\{\widetilde v_1\})
\leq\frac{\varepsilon}{\gamma-\varepsilon}.
}
\]
Notice that this estimate must only be \emph{asserted}
when its hypotheses are satisfied. A numerical
comparison outside that range is not a test of the
stated theorem. The program marks such entries as
\texttt{NaN} (not applicable).

\begin{lstlisting}[language=Matlab,caption={Experiment 4: eigenvector rotation, Davis--Kahan bounds, and a two-dimensional cluster.},label={lst:experiment_dk}]
% experiment_davis_kahan.m
% Rotate eigenvectors while changing the spectral gap, then test bounds.
clear; clc; close all;
epsilon = 0.1;
gaps = [0.1 0.15 0.2 0.3 0.5 1 2 5 10];
observed = zeros(size(gaps));
exact = zeros(size(gaps));
bound = nan(size(gaps)); % NaN when hypotheses do not apply
fprintf(' gap       distance      exact sin(theta)    DK bound\n');
for j = 1:numel(gaps)
    gamma = gaps(j);
    A = diag([0,gamma]);
    E = [0,epsilon;epsilon,0];
    [U,D] = eig(A+E);
    [~,idx] = sort(diag(D),'ascend');
    u = U(:,idx(1));
    v = [1;0];
    % Projector distance is independent of eigenvector sign.
    observed(j) = norm(v*v' - u*u',2);
    theta = 0.5*atan(2*epsilon/gamma);
    exact(j) = sin(theta);
    if epsilon < gamma/2 % condition in the original-gap corollary
        bound(j) = epsilon/(gamma-epsilon);
        assert(observed(j) <= bound(j)+1e-10);
    end
    assert(abs(observed(j)-exact(j)) < 1e-10);
    assert(max(abs(sort(eig(A+E))-sort(eig(A)))) <= epsilon+1e-10);
    fprintf('%6.2f     %10.6f      %10.6f       %10.6f\n', ...
        gamma,observed(j),exact(j),bound(j));
end

figure('Name','Davis-Kahan: gap sensitivity');
semilogx(gaps,observed,'o-','LineWidth',1.5); hold on;
semilogx(gaps,exact,'k--','LineWidth',1.3);
semilogx(gaps,bound,'s:','LineWidth',1.3);
xlabel('Original eigenvalue gap \gamma');
ylabel('Distance between one-dimensional eigenspaces');
legend('Observed projector distance','Exact sin(\theta)', ...
    'Davis-Kahan bound (where applicable)','Location','best');
grid on;
saveas(gcf,'fig_davis_kahan.png');

% Optional invariant-subspace experiment: a tightly clustered pair.
rng(21);
A = diag([0,0.05,3,4,5]);
Z = randn(5); Z = (Z+Z')/2;
E = 0.05*Z/norm(Z,2);
[V0,D0] = eig(A);
[~,i0] = sort(diag(D0),'ascend'); V0 = V0(:,i0);
[V1,D1] = eig(A+E);
[~,i1] = sort(diag(D1),'ascend'); V1 = V1(:,i1);
W = V0(:,1:2); Wtilde = V1(:,1:2);
cluster_dist = norm(W*W' - Wtilde*Wtilde',2);
gamma_cluster = 3-0.05;
bound_cluster = norm(E,2)/(gamma_cluster-norm(E,2));
fprintf('\nTwo-dimensional cluster: distance %.7f; bound %.7f\n', ...
    cluster_dist,bound_cluster);
assert(cluster_dist <= bound_cluster+1e-10);
\end{lstlisting}

\paragraph{What should the reader observe?}
The observed projector distance agrees with the exact
$\sin\theta$ expression and becomes larger as the
eigenvalue gap decreases. The following approximate
values are predicted by the analytic formula:
\[
\begin{array}{c|c|c}
\gamma&\sin\theta&\varepsilon/(\gamma-\varepsilon)\\\hline
0.3&0.289784&0.500000\\
1&0.098538&0.111111\\
2&0.049814&0.052632\\
5&0.019988&0.020408
\end{array}
\]
For $\gamma\leq2\varepsilon$, the stated smallness
condition fails, so the script does not display this
particular Davis--Kahan bound. Notice that Weyl's
eigenvalue bound continues to hold regardless of the
gap: it is the \emph{eigenvector rotation} that is
especially sensitive to near-colliding eigenvalues.

The final part of the script selects a two-dimensional
\emph{cluster} $\{0,0.05\}$ in the matrix
$A=\diag(0,0.05,3,4,5)$ and adds a symmetric
perturbation of norm $0.05$. The two selected
eigenvalues are close to one another, but the cluster
is separated from the other eigenvalues by
$\gamma=3-0.05=2.95$. The relevant distance is
\[
\norm{WW^T-\widetilde W\widetilde W^T}_2,
\]
which compares \emph{the two-dimensional spaces}, not
individual eigenvectors chosen within them. It should
satisfy the bound $0.05/(2.95-0.05)$, approximately
$0.01724$. This illustrates why Davis--Kahan is often
most naturally stated for invariant subspaces.

\subsection{Optional: Run the Complete Suite}
\label{subsec:matlab_all}

A short driver executes the four experiments in a
fixed order. Because each standalone script begins
with \texttt{clear; clc; close all}, the final figure
may be the only one left on screen after this driver
finishes, but all four PNG files will have been saved.
Readers may instead run the scripts individually
to inspect the plots one at a time.

\begin{lstlisting}[language=Matlab,caption={Driver for running all four MATLAB experiments.},label={lst:experiment_all}]
% run_all_experiments.m
% Run experiments one at a time to read the Command Window outputs.
% If you only want guaranteed theorem-verification experiments,
% run the last three scripts without the educational solver demo.
experiment_principal_angles;
experiment_weyl;
experiment_davis_kahan;
experiment_courant_fischer;
\end{lstlisting}

\paragraph{Concluding observations.}
These experiments reveal two different notions of
numerical stability. Weyl's inequality bounds
\emph{eigenvalue} displacement by the norm of a
symmetric perturbation, with no spectral-gap
requirement. Davis--Kahan controls the distance
between \emph{invariant subspaces}, provided the
selected part of the spectrum is sufficiently
separated from its complement. Courant--Fischer,
in turn, explains why the ordered eigenvalues can
be studied through extremal Rayleigh quotients.
The numerical tests demonstrate these principles,
but do not replace their proofs. In particular,
a successful run of an experimental iteration
is not by itself evidence of convergence for
all matrices or all initial guesses.



% ============================================================
% SECTION 8: CONCLUSION
% ============================================================

\section{Conclusion}
\label{sec:conclusion}

\subsection{What Have We Learned?}

The central question motivating this note was simple to
state but surprisingly subtle to answer: if a symmetric
matrix changes slightly, what can we say about its
eigenvalues and the subspaces spanned by its eigenvectors?
We have approached this question progressively, beginning
with the geometry of vectors and subspaces and ending with
rigorous perturbation bounds and their numerical
illustrations.

Our mathematical background section established the tools
needed for this journey. The spectral theorem guarantees
an orthonormal eigenbasis for every real symmetric matrix,
while the Rayleigh quotient expresses spectral information
through the scalar quantity
\[
R_A(x)=\frac{x^TAx}{x^Tx},\qquad x\neq0.
\]
The power method introduced the computational viewpoint:
eigenvalue problems can be approached by iteration rather
than by explicitly solving a characteristic polynomial.
Singular value decomposition, orthogonal projections, and
matrix norms supplied a natural language for comparing
vectors, matrices, and subspaces.

In \cref{sec:subspace_geometry}, we developed a notion of
distance between two subspaces that does not depend on the
particular bases used to represent them. If $P_{\mathcal U}$
and $P_{\mathcal W}$ denote their orthogonal projectors,
we defined
\[
\boxed{
\dist(\mathcal U,\mathcal W)
=\norm{P_{\mathcal U}-P_{\mathcal W}}_2.
}
\]
We verified that this is a metric and connected it to
principal angles. When the subspaces have equal dimension,
\[
\boxed{
\dist(\mathcal U,\mathcal W)=\sin\theta_{\max}.
}
\]
This was more than an exercise in geometry: it gave us a
basis-independent way to describe the rotation of an
invariant subspace. In particular, it avoids the misleading
comparison of two eigenvectors that differ only by a sign.

\begin{samepage}
We next studied the Courant--Fischer theorem, which extends
Rayleigh's extremal characterization from the smallest and
largest eigenvalues to every intermediate eigenvalue. With
our convention $\lambda_1\leq\cdots\leq\lambda_n$, we
proved the two equivalent characterizations
\begin{gather*}
\boxed{
\lambda_k(A)=
\min_{\dim\mathcal S=k}
\max_{\substack{x\in\mathcal S\\x\neq0}}R_A(x)
}\\[1.2ex]
\boxed{
\lambda_k(A)=
\max_{\dim\mathcal T=n-k+1}
\min_{\substack{x\in\mathcal T\\x\neq0}}R_A(x)
}.
\end{gather*}
\end{samepage}
The essential insight was geometric: subspaces whose
dimensions add up to more than the ambient dimension must
intersect nontrivially. That elementary observation made
it possible to turn an apparently difficult optimization
over all subspaces into an exact characterization of
ordered eigenvalues. The theorem also provided the bridge
from spectral geometry to perturbation theory.

In \cref{sec:weyl}, we used this bridge to derive Weyl's
inequality. For symmetric matrices $A$ and $E$, it states
\[
\boxed{
\abs{\lambda_k(A+E)-\lambda_k(A)}
\leq\norm{E}_2,
\qquad k=1,\ldots,n.
}
\]
Thus, the displacement of every \emph{ordered eigenvalue}
is controlled by the size of the perturbation, without any
assumption about gaps between eigenvalues. This is a
powerful guarantee of stability and a fundamental reason
that symmetric eigenvalue problems occupy such an
important place in numerical linear algebra.

But Weyl's inequality does not settle the question of
eigenvectors. In \cref{sec:davis_kahan}, we saw explicitly
that even when eigenvalues change very little, their
associated eigendirections can rotate substantially if
neighboring eigenvalues are close. The Davis--Kahan
$\sin\Theta$ theorem identifies the missing ingredient:
\emph{spectral separation}. In its appropriate
spectral-gap setting, it bounds the distance between
invariant subspaces in terms of the size of the perturbation
relative to the gap separating the selected part of the
spectrum from its complement. This is especially important
for clusters of nearby or repeated eigenvalues, where
individual eigenvectors may be ambiguous but the invariant
subspace formed by the cluster can remain well determined.

Finally, in \cref{sec:numerical_experiments}, we translated
theory into MATLAB experiments. We revisited our educational
Courant--Fischer minimax and maximin implementations,
compared numerical eigenpair estimates with reference
results and residuals, verified the relation between
principal angles and projector distance, checked Weyl's
perturbation bound, and examined how eigenspace rotation
responds to a changing spectral gap. These experiments
help us see the theorems at work, although numerical
agreement in a finite collection of tests cannot replace
a proof or guarantee convergence of every iterative method.

Taken together, the sections reveal a coherent chain of
ideas:
\[
\boxed{
\begin{gathered}
\text{spectral theorem and Rayleigh quotient}\\
\Downarrow\\
\text{subspace geometry and variational principles}\\
\Downarrow\\
\text{eigenvalue stability (Weyl)}\\
\Downarrow\\
\text{invariant-subspace stability (Davis--Kahan).}
\end{gathered}
}
\]
In scientific computing, we should not ask only whether an
algorithm produces an answer. We should also ask which
aspects of that answer are stable, how the data affect its
accuracy, and what mathematical assumptions make a
reliability guarantee possible.

\subsection{A Historical Epilogue}

As is our custom, we close with a few historical remarks
about the people whose names have accompanied us through
this note. Their stories remind us that mathematics is not
only a collection of theorems; it is also a human enterprise
built through curiosity, generosity, teaching, and the
creation of institutions in which others can work.

\paragraph{Lord Rayleigh: a physicist and a patron of science.}
John William Strutt, the third Baron Rayleigh
(1842--1919), made major contributions to the theory of
waves, acoustics, optics, and experimental physics. The
Rayleigh quotient is only one part of his scientific
legacy. His investigations of the densities of gases led
to the discovery of argon, for which he received the 1904
Nobel Prize in Physics. Yet there is another side to his
story that deserves attention. As head of the Cavendish
Laboratory at Cambridge, Rayleigh contributed his own money
toward laboratory equipment and the support of experimental
instruction. He later donated his Nobel Prize money to
the University of Cambridge for the expansion of the
Cavendish Laboratory. He was therefore not merely a
scientist who benefited from the institutions around him;
he also heavily invested in the resources through which future
scientists could conduct research.\footnote{See the
\href{https://mathshistory.st-andrews.ac.uk/Biographies/Rayleigh/}{MacTutor biography of Lord Rayleigh},
the \href{https://www.nobelprize.org/prizes/physics/1904/summary/}{Nobel Prize account of the 1904 Physics Prize},
and the\href{https://mathshistory.st-andrews.ac.uk/RS/rayleigh_rs.pdf}{Royal Society's biographical memoir of Rayleigh}.}

\paragraph{Ernst Fischer and Richard Courant: from a theorem to a school of mathematics.}
Ernst Fischer (1875--1954) belongs in this history because
the finite-dimensional minimax principle is associated
with his early twentieth-century work. Richard Courant
(1888--1972) helped establish the broader variational
point of view as a central tool of mathematical physics,
including the study of differential equations and their
eigenvalue problems.\footnote{For historical attribution,
see discussions of the minimax principle in the literature
on variational eigenvalue methods, including the historical
survey presented in
\href{https://www.sciencedirect.com/topics/mathematics/dirichlet-eigenvalue}{ScienceDirect Topics: Dirichlet Eigenvalue}.}

Courant's influence, however, cannot be measured by a
single theorem or even by the number of papers he wrote.
Forced to leave his position in G\"ottingen during the Nazi
period, he arrived at New York University in 1934 and
helped build a research community devoted to the unity of
pure mathematics, applied mathematics, and the physical
sciences. He recruited colleagues, encouraged younger
mathematicians, and helped other scholars displaced from
Europe find a new intellectual home. The research institute
that eventually bore his name became an enduring center
for applied mathematics and computation. It is in this
sense that Courant's contribution was larger than any one
institution: he helped create a community that could
continue producing mathematics long after its founder.
\footnote{See
\href{https://cims.nyu.edu/dynamic/about/history/}{NYU Courant's institutional history}
and the
\href{https://cims.nyu.edu/about/courant_bio.html}{NYU biography of Richard Courant}.}

\paragraph{Hermann Weyl: unity across mathematics and physics.}
The inequality we studied in Section~5 bears the name of
Hermann Weyl (1885--1955), one of the remarkably versatile
mathematicians of the twentieth century. His work ranged
across spectral theory, analysis, differential geometry,
group representation theory, the foundations of mathematics,
relativity, and quantum mechanics. In particular, he helped
make symmetry and group theory part of the language of
modern theoretical physics. After leaving G\"ottingen in
1933, Weyl joined the Institute for Advanced Study in
Princeton. The breadth of his work is a useful reminder
that a theorem which looks like a statement about matrix
eigenvalues may be part of a much larger intellectual
story.\footnote{See the
\href{https://plato.stanford.edu/entries/weyl/}{Stanford Encyclopedia of Philosophy biography of Hermann Weyl}
and the
\href{https://www.ias.edu/hermann-weyl-work}{Institute for Advanced Study account of his work}.}

\paragraph{Chandler Davis and William Kahan: making stability precise.}
The Davis--Kahan theorem bears the names of Chandler Davis
and William M. Kahan. Their 1970 paper, \emph{The Rotation
of Eigenvectors by a Perturbation. III}, studied precisely
the issue that motivated the final part of this note: how
much an invariant subspace may rotate when a symmetric or
Hermitian operator is perturbed. Their work expressed this
change through angles between subspaces and related it to
spectral separation. Kahan later became an influential
figure in numerical computation more broadly, helping
shape the IEEE standard for floating-point arithmetic.
There is a particularly fitting connection here: both
perturbation theory and floating-point arithmetic concern
what we can \emph{trust} about the results of a computation.
\footnote{See
\href{https://epubs.siam.org/doi/abs/10.1137/0707001}{Davis and Kahan, \emph{SIAM Journal on Numerical Analysis} 7 (1970), 1--46},
and the
\href{https://www2.eecs.berkeley.edu/Faculty/Homepages/kahan.html}{UC Berkeley biography of William M. Kahan}.}

There is a common thread running through these lives.
Rayleigh combined discovery with material support for
scientific research; Fischer and Courant helped turn
spectral optimization into a lasting mathematical method;
Weyl moved freely among geometry, analysis, and physics;
and Davis and Kahan helped explain why numerical spectral
information can, under suitable assumptions, be trusted.
Their achievements were different, but each contributed to
a mathematical culture in which ideas are not merely
proved and stored away: they are taught, refined, tested,
and used to understand new problems.

We hope that this note has followed that spirit. The
formulas themselves are important, but their greatest
value lies in the habits of thought they encourage:
visualize the geometry, question the assumptions, follow
the proof carefully, and verify the consequences through
computation. In the end, understanding \emph{why} a result
is true makes it much more useful than remembering only
\emph{what} it says.


% ============================================================
% REFERENCES
% ============================================================

% ============================================================
% REFERENCES
% ============================================================

\clearpage
\phantomsection
\addcontentsline{toc}{section}{References}

\begin{thebibliography}{9}

\bibitem{murthy_matrix}
C.~R. Murthy,
\emph{E2 212: Matrix Theory---Course Lectures and Notes},
Indian Institute of Science, Bengaluru.
\href{https://ece.iisc.ac.in/~cmurthy/e2-212-matrix-theory-video-lectures-and-notes-2/}
{Course materials (online)}.

\bibitem{spielman_sgt}
D.~A. Spielman,
\emph{CPSC 462/562: Spectral Graph Theory},
course lectures and notes, Yale University, Spring 2025.
\href{https://www.cs.yale.edu/homes/spielman/462/462schedule.html}
{Course materials (online)}.

\bibitem{horn_johnson}
R.~A. Horn and C.~R. Johnson,
\emph{Matrix Analysis}, 2nd ed.,
Cambridge University Press, 2013.

\bibitem{golub_vanloan}
G.~H. Golub and C.~F. Van Loan,
\emph{Matrix Computations}, 4th ed.,
Johns Hopkins University Press, 2013.

\bibitem{stewart_sun}
G.~W. Stewart and J.-G. Sun,
\emph{Matrix Perturbation Theory},
Academic Press, 1990.

\bibitem{davis_kahan}
C. Davis and W.~M. Kahan,
``The Rotation of Eigenvectors by a Perturbation. III,''
\emph{SIAM Journal on Numerical Analysis},
vol.~7, no.~1, pp.~1--46, 1970.
\href{https://doi.org/10.1137/0707001}{doi:10.1137/0707001}.

\end{thebibliography}

\end{document}